% Leçon 39 — Vocabulaire et compréhension de texte : rédaction d’une lettre formelle — Chinois 1ère A — Cours signé OnBuch+
% Programme officiel MINESEC. Compiler avec : tectonic ZHA39-vocabulaire-redaction-dune-lettre-formelle.tex
\newcommand{\DOCMATIERE}{Chinois}
\newcommand{\DOCNIVEAU}{1ère A}
\newcommand{\DOCMODULE}{Module 5 : Mass médias et télécommunication}
\newcommand{\DOCLECON}{ZHA39}
\newcommand{\DOCTITRE}{Leçon 39 — Vocabulaire et compréhension de texte : rédaction d’une lettre formelle}
% OnBuch+ — préambule « cours signés » ENRICHI (compilé avec tectonic / XeLaTeX)
% Système visuel « Pop » d'OnBuch+ : fond crème (jamais blanc), bordures encre
% épaisses, ombres « dures » décalées, titres Archivo Black en capitales.
% Version MATHÉMATIQUES : graphes (pgfplots/tikz), boîtes pédagogiques
% (définition, propriété, méthode, attention, le savais-tu, exercice résolu,
% exercice, TP, bilan), page de garde et signature OnBuch+ ancrée en bas.
%
% Macros attendues dans main.tex AVANT \input{preamble} :
%   \DOCMATIERE  \DOCNIVEAU  \DOCMODULE  \DOCLECON (numéro)  \DOCTITRE
\documentclass[11pt]{article}

\usepackage{fontspec}
\usepackage[french]{babel} % français = langue principale ; le chinois via \zh{..} / textebox (xeCJK)
\usepackage{xeCJK}
\usepackage{pifont} % \ding{51} \ding{55} utilisés par les modèles
\usepackage[a4paper,margin=2cm,top=3.4cm,bottom=2.8cm,headheight=1.4cm,footskip=1.3cm]{geometry}
\usepackage{amsmath,amssymb,amsfonts}
\usepackage{mathtools} % \xmapsto, \coloneqq... souvent utilisés par les modèles
\usepackage{cancel} % simplifications barrées (bilans de forces)
% \abs{z} (module, valeur absolue) et \norm{u} (norme), utilisés par les modèles
\providecommand{\abs}{}\let\abs\relax\DeclarePairedDelimiter{\abs}{\lvert}{\rvert}
\providecommand{\norm}{}\let\norm\relax\DeclarePairedDelimiter{\norm}{\lVert}{\rVert}
\usepackage[table]{xcolor} % [table] : fournit aussi \rowcolors (alternance de lignes)
\usepackage{pagecolor}
\usepackage{tikz}
\usetikzlibrary{arrows.meta,positioning,calc,shapes.geometric,shapes.misc,shapes.arrows,decorations.pathmorphing,decorations.pathreplacing,decorations.markings,patterns,shadows,mindmap,trees,fit,backgrounds,matrix,intersections,angles,quotes}
\usepackage{pgfplots}
\usepackage[version=4]{mhchem} % équations nucléaires \ce{^{235}_{92}U}
\usepackage[european,siunitx]{circuitikz} % schémas électriques
\pgfplotsset{compat=1.18}
\usepgfplotslibrary{fillbetween,groupplots} % aires entre courbes (calcul intégral)
\usepackage{enumitem}
\usepackage{fancyhdr}
% [most] charge toutes les bibliothèques tcolorbox (listings, minted, raster,
% documentation...) et gonfle inutilement la mémoire TeX sur un document long
% (~300 boîtes) : on ne charge que ce qui est réellement utilisé.
\usepackage[skins,breakable]{tcolorbox}
\usepackage{graphicx}
\usepackage{colortbl}
\usepackage{booktabs}
\usepackage{tabularx}
\usepackage{diagbox} % case d'en-tête diagonale des tableaux à double entrée
% Colonnes extensibles centrée / gauche / droite (C, L, R), souvent utilisées par les modèles
\newcolumntype{C}{>{\centering\arraybackslash}X}
\newcolumntype{L}{>{\raggedright\arraybackslash}X}
\newcolumntype{R}{>{\raggedleft\arraybackslash}X}
\usepackage{array}
\usepackage{multicol}
\usepackage{multirow}
\usepackage{siunitx}
% parse-numbers=false : accepte aussi \SI{2,5 \times 10^{-3}}{...}
\sisetup{locale=FR,output-decimal-marker={,},group-minimum-digits=4,parse-numbers=false}
\DeclareSIUnit\ohm{\ensuremath{\symup{\Omega}}} % U+2126 absent de la police
\DeclareSIUnit\division{div} % réglages d'oscilloscope (ms/div, V/div)
\DeclareSIUnit\tr{tr} % tours (vitesse de rotation, ex. \SI{1500}{\tr\per\minute})
\DeclareSIUnit\molar{M} % concentration molaire (ex. \SI{3}{\molar}, \SI{1}{\micro\molar})
\DeclareSIUnit\an{an} % année (le modèle confond parfois avec \year, un compteur TeX) — ex. \SI{5}{\kilo\meter\per\an}
\DeclareSIUnit\FCFA{FCFA} % franc CFA (ex. \SI{500}{\FCFA\per\kilo\gram})
\DeclareSIUnit\degreeBrix{\textdegree Brix} % teneur en sucre d'un sirop/jus (réfractométrie)
\DeclareSIUnit\month{mois} % le modèle utilise parfois \month, un compteur TeX, comme unité de durée
\usepackage{qrcode}
% \degree : commande fréquemment utilisée par les modèles pour un angle ou une
% température en dehors de \SI{}{} (non fournie par les paquets chargés).
\providecommand{\degree}{^{\circ}}
% \qed (fin de démonstration), utilisé par les modèles sans amsthm
\providecommand{\qed}{\hfill$\square$}
% \barre : double barre des valeurs interdites dans un tableau de variations (tkz-tab)
\providecommand{\barre}{\ensuremath{\|}}
% environnement proof (amsthm non chargé) utilisé par les modèles
\ifdefined\proof\else\newenvironment{proof}[1][Démonstration]{\par\noindent\textit{#1.}\ }{\qed\par}\fi
\usepackage{lastpage}
\usepackage{hyperref}
\hypersetup{hidelinks}

% ============================================================
% Palette « Pop » OnBuch+ — mêmes hex que la classe P (plus_ui.dart)
% ============================================================
\definecolor{poporange}{HTML}{F59321}
\definecolor{popdark}{HTML}{E07A0C}
\definecolor{popcream}{HTML}{FFF6E8}
\definecolor{popink}{HTML}{1C1714}
\definecolor{poppurple}{HTML}{7A5AE0}
\definecolor{popgreen}{HTML}{2E9E6A}
\definecolor{poppink}{HTML}{E0457B}
\definecolor{popblue}{HTML}{2F7DE1}
\definecolor{popgold}{HTML}{C98A12}
\definecolor{popmuted}{HTML}{8A7F76}
\definecolor{popline}{HTML}{EADFD2}
\definecolor{popgray}{HTML}{8A7F76}
\colorlet{popgrey}{popgray}
% Alias tolérants (le modèle nomme parfois les couleurs différemment)
\colorlet{popwhite}{white}
\colorlet{popblack}{popink}
\colorlet{popred}{poppink}
\colorlet{popyellow}{popgold}
% Teintes claires (fonds de tableaux/encadrés) — le modèle nomme parfois
% les couleurs avec un suffixe L (ex. popblueL) sans les avoir définies.
\colorlet{poporangeL}{poporange!20}
\colorlet{popdarkL}{popdark!20}
\colorlet{poppurpleL}{poppurple!20}
\colorlet{popgreenL}{popgreen!20}
\colorlet{poppinkL}{poppink!20}
\colorlet{popblueL}{popblue!20}
\colorlet{popgoldL}{popgold!20}
\colorlet{popredL}{popred!20}
\colorlet{popgrayL}{popgray!20}
% Teintes claires pour fonds de tableaux / zones
\colorlet{popblueL}{popblue!10!white}
\colorlet{poporangeL}{poporange!14!white}
\colorlet{popgreenL}{popgreen!12!white}
\colorlet{poppurpleL}{poppurple!12!white}
\colorlet{poppinkL}{poppink!12!white}
\colorlet{popgoldL}{popgold!14!white}

\pagecolor{popcream}

% ============================================================
% Typographie
% ============================================================
\defaultfontfeatures{Ligatures=TeX}
\setmainfont{PlusJakartaSans-Medium.ttf}[Path=../../../fonts/,BoldFont=PlusJakartaSans-Bold.ttf]
% Symboles, API et emojis absents de la police principale : repli sur DejaVu Sans (caractères actifs)
\newfontface\symfont{DejaVuSans.ttf}[Path=../../../fonts/]
\makeatletter
\newcommand\symactive[1]{\catcode#1=13 \begingroup\lccode`\~=#1 \lowercase{\endgroup\def~}{{\symfont\char#1}}}
\newcommand\symblank[1]{\catcode#1=13 \begingroup\lccode`\~=#1 \lowercase{\endgroup\def~}{}}
\newcommand\symhyph[1]{\catcode#1=13 \begingroup\lccode`\~=#1 \lowercase{\endgroup\def~}{\mbox{-}}}
\newcount\symc
\newcommand\symrange[2]{\symc=#1 \@whilenum\symc<#2 \do{\expandafter\symactive\expandafter{\the\symc}\advance\symc by 1 }}
\newcommand\symblankrange[2]{\symc=#1 \@whilenum\symc<#2 \do{\expandafter\symblank\expandafter{\the\symc}\advance\symc by 1 }}
\makeatother
\symrange{"0250}{"02FF}\symactive{"2713}\symactive{"2714}\symactive{"2717}\symactive{"2718}\symactive{"2610}\symactive{"2611}\symactive{"2612}
\symrange{"2600}{"27BF}
\catcode"2705=13 \begingroup\lccode`\~="2705 \lowercase{\endgroup\def~}{{\symfont\char"2713}}\symblank{"26BD}\symblank{"26A1}
\symrange{"2460}{"257F}\symactive{"223C}\symactive{"2248}\symactive{"2260}
\symactive{"01F8}\symactive{"01F9}\symactive{"1E3E}\symactive{"1E3F}
\symhyph{"2011}\symhyph{"2E1A}\symblankrange{"1F300}{"1FAFF}\symblank{"FE0F}
% Caractères chinois : WenQuanYi Zen Hei (sous-ensemble GB2312 + pinyin), gras/italique simulés
\setCJKmainfont{WenQuanYiZenHei-subset.ttf}[Path=../../../fonts/,AutoFakeBold=3,AutoFakeSlant=0.2]
\xeCJKsetup{CJKspace=false}
% Pinyin : voyelles à caron (ǎ ǐ ǒ ǔ ǖ ǘ ǚ ǜ) absentes de la police principale -> caractères actifs
\newfontface\pyfont{WenQuanYiZenHei-subset.ttf}[Path=../../../fonts/]
\newcommand\pyactive[1]{\catcode#1=13 \begingroup\lccode`\~=#1 \lowercase{\endgroup\def~}{{\pyfont\char#1}}}
\pyactive{"01CD}\pyactive{"01CE}\pyactive{"01CF}\pyactive{"01D0}\pyactive{"01D1}\pyactive{"01D2}\pyactive{"01D3}\pyactive{"01D4}\pyactive{"01D5}\pyactive{"01D6}\pyactive{"01D7}\pyactive{"01D8}\pyactive{"01D9}\pyactive{"01DA}\pyactive{"01DB}\pyactive{"01DC}
% Maths en sans-serif (Fira Math) avec chiffres et lettres droites de Plus
% Jakarta, pour que les formules se fondent dans le texte.
\usepackage[math-style=ISO,bold-style=ISO]{unicode-math}
\setmathfont{FiraMath-Regular.otf}[Path=../../../fonts/]
\setmathfont{PlusJakartaSans-Medium.ttf}[Path=../../../fonts/,range={up/num,up/{Latin,latin},\mathup}]
\usepackage{icomma} % virgule décimale sans espace en mode math
% Caractères mathématiques Unicode absents des polices de texte (Plus Jakarta,
% Archivo Black) : rendus par leur équivalent mathématique, en texte comme en
% formule — sinon ils s'affichent en carré vide (ex. « Arithmétique dans ℤ »).
\usepackage{newunicodechar}
\newunicodechar{ℝ}{\ensuremath{\mathbb{R}}}
\newunicodechar{ℤ}{\ensuremath{\mathbb{Z}}}
\newunicodechar{ℂ}{\ensuremath{\mathbb{C}}}
\newunicodechar{ℕ}{\ensuremath{\mathbb{N}}}
\newunicodechar{ℚ}{\ensuremath{\mathbb{Q}}}
\newunicodechar{𝔻}{\ensuremath{\mathbb{D}}}
\newunicodechar{α}{\ensuremath{\alpha}}
\newunicodechar{β}{\ensuremath{\beta}}
\newunicodechar{γ}{\ensuremath{\gamma}}
\newunicodechar{δ}{\ensuremath{\delta}}
\newunicodechar{ε}{\ensuremath{\varepsilon}}
\newunicodechar{θ}{\ensuremath{\theta}}
\newunicodechar{λ}{\ensuremath{\lambda}}
\newunicodechar{μ}{\ensuremath{\mu}}
\newunicodechar{σ}{\ensuremath{\sigma}}
\newunicodechar{τ}{\ensuremath{\tau}}
\newunicodechar{φ}{\ensuremath{\varphi}}
\newunicodechar{ψ}{\ensuremath{\psi}}
\newunicodechar{ω}{\ensuremath{\omega}}
\newunicodechar{Σ}{\ensuremath{\Sigma}}
% majuscules produites par \MakeUppercase dans les titres (α-aminés -> Α)
\newunicodechar{Α}{\ensuremath{\alpha}}
\newunicodechar{Β}{\ensuremath{\beta}}
\newunicodechar{Γ}{\ensuremath{\Gamma}}
\newunicodechar{Δ}{\ensuremath{\Delta}}
\newunicodechar{Ε}{\ensuremath{\varepsilon}}
\newunicodechar{Θ}{\ensuremath{\Theta}}
\newunicodechar{Λ}{\ensuremath{\Lambda}}
\newunicodechar{Μ}{\ensuremath{\mu}}
\newunicodechar{Τ}{\ensuremath{\tau}}
\newunicodechar{Φ}{\ensuremath{\Phi}}
\newunicodechar{Ψ}{\ensuremath{\Psi}}
\newunicodechar{Ω}{\ensuremath{\Omega}}
\newunicodechar{↦}{\ensuremath{\mapsto}}
\newunicodechar{∈}{\ensuremath{\in}}
\newunicodechar{∉}{\ensuremath{\notin}}
\newunicodechar{∧}{\ensuremath{\wedge}}
\newunicodechar{⇔}{\ensuremath{\Leftrightarrow}}
\newunicodechar{⟺}{\ensuremath{\Longleftrightarrow}}
\newunicodechar{⇒}{\ensuremath{\Rightarrow}}
\newunicodechar{⟹}{\ensuremath{\Longrightarrow}}
\newunicodechar{∘}{\ensuremath{\circ}}
\newunicodechar{∪}{\ensuremath{\cup}}
\newunicodechar{∩}{\ensuremath{\cap}}
\newunicodechar{≡}{\ensuremath{\equiv}}
\newunicodechar{⊂}{\ensuremath{\subset}}
\newunicodechar{⊥}{\ensuremath{\perp}}
\newunicodechar{∃}{\ensuremath{\exists}}
\newunicodechar{∀}{\ensuremath{\forall}}
\newunicodechar{∛}{\ensuremath{\sqrt[3]{\,}}}
\newunicodechar{∜}{\ensuremath{\sqrt[4]{\,}}}
\newunicodechar{⃗}{\ensuremath{{}^{\to}}}
\newunicodechar{ⁿ}{\ifmmode^{n}\else\textsuperscript{\ensuremath{n}}\fi}
\newunicodechar{ˣ}{\ifmmode^{x}\else\textsuperscript{\ensuremath{x}}\fi}
\newunicodechar{ᵇ}{\ifmmode^{b}\else\textsuperscript{\ensuremath{b}}\fi}
\newunicodechar{ʳ}{\ifmmode^{r}\else\textsuperscript{\ensuremath{r}}\fi}
\newunicodechar{ᵘ}{\ifmmode^{u}\else\textsuperscript{\ensuremath{u}}\fi}
\newunicodechar{ʸ}{\ifmmode^{y}\else\textsuperscript{\ensuremath{y}}\fi}
\newunicodechar{ᵃ}{\ifmmode^{a}\else\textsuperscript{\ensuremath{a}}\fi}
\newunicodechar{ᵉ}{\ifmmode^{e}\else\textsuperscript{\ensuremath{e}}\fi}
\newunicodechar{ᵏ}{\ifmmode^{k}\else\textsuperscript{\ensuremath{k}}\fi}
\newunicodechar{ᵖ}{\ifmmode^{p}\else\textsuperscript{\ensuremath{p}}\fi}
\newunicodechar{ᴺ}{\ifmmode^{N}\else\textsuperscript{\ensuremath{N}}\fi}
\newunicodechar{ᶜ}{\ifmmode^{c}\else\textsuperscript{\ensuremath{c}}\fi}
\newunicodechar{ʰ}{\ifmmode^{h}\else\textsuperscript{\ensuremath{h}}\fi}
\newunicodechar{ᵠ}{\ifmmode^{q}\else\textsuperscript{\ensuremath{q}}\fi}
\newunicodechar{ₙ}{\ifmmode_{n}\else\textsubscript{\ensuremath{n}}\fi}
\newunicodechar{ₐ}{\ifmmode_{a}\else\textsubscript{\ensuremath{a}}\fi}
\newunicodechar{ᵢ}{\ifmmode_{i}\else\textsubscript{\ensuremath{i}}\fi}
\newunicodechar{ᵣ}{\ifmmode_{r}\else\textsubscript{\ensuremath{r}}\fi}
\newunicodechar{ₖ}{\ifmmode_{k}\else\textsubscript{\ensuremath{k}}\fi}
\newunicodechar{ᵦ}{\ifmmode_{\beta}\else\textsubscript{\ensuremath{\beta}}\fi}
\newunicodechar{ₓ}{\ifmmode_{x}\else\textsubscript{\ensuremath{x}}\fi}
\newfontfamily\popdisplay{ArchivoBlack-Regular.ttf}[Path=../../../fonts/]
\newfontfamily\poplogo{SpaceGrotesk-Bold.ttf}[Path=../../../fonts/]
\newfontfamily\popsemi{PlusJakartaSans-SemiBold.ttf}[Path=../../../fonts/]
\newfontfamily\popxbold{PlusJakartaSans-ExtraBold.ttf}[Path=../../../fonts/]
\newfontfamily\popxboldls{PlusJakartaSans-ExtraBold.ttf}[Path=../../../fonts/,LetterSpace=3.0]

\setlist[enumerate]{leftmargin=1.7em,itemsep=5pt,topsep=4pt}
\setlist[itemize]{leftmargin=1.7em,itemsep=4pt,topsep=4pt,label={\color{poporange}\textbullet}}
\setlength{\parindent}{0pt}
\setlength{\parskip}{5pt}
\renewcommand{\arraystretch}{1.25}

% pgfplots : style Pop par défaut
\pgfplotsset{
  every axis/.append style={axis line style={popink,line width=1pt},tick style={popink},
    grid style={popline,line width=0.5pt},label style={font=\small},tick label style={font=\footnotesize}},
  popcurve/.style={poporange,line width=1.8pt,no markers,smooth},
}
% Même style disponible dans un \draw TikZ simple (hors pgfplots)
\tikzset{popcurve/.style={poporange,line width=1.8pt}}
\tikzset{popgrid/.style={popline,very thin}}
\colorlet{popcurve}{poporange} % le nom du style est parfois utilisé comme couleur

% ============================================================
% Pill / squiggle
% ============================================================
\newcommand{\poppill}[3]{%
  \tikz[baseline=-0.55ex]\node[draw=popink,line width=1.2pt,rounded corners=2.6pt,%
    fill=#2,inner xsep=6.5pt,inner ysep=3pt]{%
    \popxboldls\fontsize{7}{7}\selectfont\color{#3}\MakeUppercase{#1}};%
}
\newcommand{\popsquiggle}[1]{%
  \tikz[baseline=-0.4ex]{
    \draw[#1,line width=2.6pt,line cap=round]
      plot [smooth] coordinates {(0,0.09) (0.34,0.01) (0.68,0.21) (1.02,0.01) (1.36,0.10)};
  }%
}
% Mot-clé surligné (marqueur orange sous le texte)
\newcommand{\cle}[1]{{\popxbold\color{popdark}#1}}

% ============================================================
% En-tête / pied
% ============================================================
\pagestyle{fancy}
\fancyhf{}
\renewcommand{\headrulewidth}{0pt}
\renewcommand{\footrulewidth}{0pt}
\lhead{\raisebox{-0.15ex}{\poplogo\fontsize{13}{13}\selectfont\color{popink}OnBuch\color{poporange}+}}
\rhead{\poppill{\DOCMATIERE\ \textbullet\ \DOCNIVEAU\ \textbullet\ Leçon \DOCLECON}{white}{popink}}
\lfoot{\popxbold\fontsize{7.6}{7.6}\selectfont\color{popmuted}\MakeUppercase{Cours signé OnBuch+}\ \textbullet\ {\popsemi\DOCTITRE}}
\rfoot{\tikz[baseline=-0.6ex]\node[rounded corners=8.5pt,fill=popink,minimum height=17pt,minimum width=17pt,inner xsep=5pt,inner sep=0]{\popxbold\fontsize{8}{8}\selectfont\color{popcream}\thepage\,/\,\pageref*{LastPage}};}

% ============================================================
% Titres
% ============================================================
\newcommand{\chaptitre}[2]{%
  \begin{tcolorbox}[enhanced,colback=poporange,colframe=popink,boxrule=2.6pt,arc=11pt,
      drop shadow=popink,left=15pt,right=15pt,top=12pt,bottom=11pt,before skip=3pt,after skip=18pt]
    \poppill{#2}{popink}{popcream}\par\vspace{9pt}
    {\raggedright\hyphenpenalty=10000\exhyphenpenalty=10000\popdisplay\fontsize{22}{24}\selectfont\color{popcream}\MakeUppercase{#1}\par}
  \end{tcolorbox}
}
% Section numérotée : pastille orange avec le numéro + titre Archivo
\newcounter{popsec}
\newcommand{\coursec}[1]{%
  \stepcounter{popsec}\setcounter{popsub}{0}%
  \par\vspace{16pt}\noindent
  \begin{minipage}{\linewidth}
  \tikz[baseline=-0.7ex]\node[draw=popink,line width=1.6pt,fill=poporange,rounded corners=4pt,minimum size=22pt,inner sep=0]{\popdisplay\fontsize{12}{12}\selectfont\color{popcream}\thepopsec};%
  \hspace{8pt}{\popdisplay\fontsize{15}{17}\selectfont\color{popink}\MakeUppercase{#1}}\par
  \vspace{4pt}\noindent{\color{poporange}\rule{36pt}{3.6pt}}
  \end{minipage}\par\nopagebreak\vspace{8pt}
}
\newcounter{popsub}
\newcommand{\courssub}[1]{%
  \stepcounter{popsub}%
  \par\vspace{9pt}\noindent{\popxbold\fontsize{11.5}{13}\selectfont\color{popink}{\color{poporange}\thepopsec.\thepopsub}\hspace{6pt}#1}\par\nopagebreak\vspace{4pt}
}

% ============================================================
% Boîtes pédagogiques — même recette : fond blanc, bordure épaisse colorée,
% ombre dure encre, badge plein. Titre optionnel [#1].
% ============================================================
\tcbset{popbase/.style={breakable,enhanced,colback=white,boxrule=2.3pt,arc=9pt,
  shadow={3pt}{-3pt}{0pt}{popink},left=14pt,right=14pt,top=11pt,bottom=11pt,before skip=11pt,after skip=15pt,
  fontupper=\popsemi\fontsize{10.6}{14.5}\selectfont\color{popink},
  fontlower=\popsemi\fontsize{10.6}{14.5}\selectfont\color{popink}}}
% Coche dessinée (le glyphe ✓ n'existe pas dans les polices du document)
\newcommand{\popcheck}[1]{\tikz[baseline=-0.5ex]\draw[#1,line width=1.6pt,line cap=round,line join=round](-0.13,0.0)--(-0.04,-0.09)--(0.13,0.1);}
\providecommand{\coche}{\popcheck{popgreen}}
\providecommand{\resultat}[1]{\par\smallskip\noindent\fcolorbox{popgreen}{popgreenL}{\parbox{\dimexpr\linewidth-2\fboxsep-2\fboxrule\relax}{\textbf{Résultat.} #1}}\par\smallskip}
\newcommand{\popboxhead}[3]{\poppill{#1}{#2}{white}\ifx\relax#3\relax\else\hspace{6pt}{\popxbold\fontsize{10.6}{12}\selectfont\color{popink}#3}\fi\par\vspace{7pt}}

\newtcolorbox{aretenir}[1][]{popbase,colframe=popgreen,before upper={\popboxhead{À retenir}{popgreen}{#1}}}
% Texte en chinois (document de lecture, dialogue, extrait) : cadre
% d'étude. Un titre optionnel (source, auteur) s'affiche dans l'en-tête.
\newtcolorbox{textebox}[1][]{popbase,colframe=popblue,before upper={\popboxhead{文本 · Texte}{popblue}{#1}},after upper={}}
% Marques ✓ / ✗ (forme correcte / fautive) souvent utilisées par les modèles
\providecommand{\cross}{\textcolor{poppink}{\ensuremath{\times}}}
\providecommand{\xmark}{\cross}\providecommand{\crossmark}{\cross}\providecommand{\cmark}{\ensuremath{\checkmark}}
% Ligne de réponse / justification à compléter (inventée par les modèles)
\providecommand{\justification}[1]{\par\noindent\textit{Justification~:} \dotfill\par}
% \textipa (package tipa non chargé) : la police ne couvre pas l'API -> simple passage
\providecommand{\textipa}[1]{#1}
% Soulignement (ulem non chargé) : \uline, \ul
\providecommand{\uline}[1]{\underline{#1}}\providecommand{\ul}[1]{\underline{#1}}
% \glossaire{\item ...} inventée par les modèles : liste de glossaire
\providecommand{\glossaire}[1]{\begin{itemize}#1\end{itemize}}
% \alert (beamer) inventée par les modèles : texte mis en évidence
\providecommand{\alert}[1]{\textbf{\textcolor{poppink}{#1}}}
% variantes ASCII d'ordinaux écrites en macro
\providecommand{\iere}{\textsuperscript{re}}\providecommand{\ieme}{\textsuperscript{e}}\providecommand{\ere}{\textsuperscript{re}}\providecommand{\eme}{\textsuperscript{e}}\providecommand{\er}{\textsuperscript{er}}\providecommand{\ie}{\textsuperscript{e}}
% Ordinaux français écrits en macro par les modèles : 1\ière, 3\ième
\newcommand{\ière}{\textsuperscript{re}}\newcommand{\ième}{\textsuperscript{e}}
% \exemple (inventée par les modèles) : étiquette « Exemple : »
\providecommand{\exemple}{\textbf{Exemple~:}\ }
% \square utilisable aussi hors mode math (cases à cocher)
\AtBeginDocument{\let\squareMath\square\renewcommand{\square}{\ensuremath{\squareMath}}}
% Mot ou phrase en chinois à l'intérieur d'un paragraphe français
\newcommand{\zh}[1]{\ifmmode\text{#1}\else{#1}\fi}
\let\eng\zh \let\esp\zh \let\all\zh % alias tolérés
% flèches d'intonation inventées par les modèles
\providecommand{\textsearrow}{}\renewcommand{\textsearrow}{\ensuremath{\searrow}}\providecommand{\textnearrow}{}\renewcommand{\textnearrow}{\ensuremath{\nearrow}}
% \fr{..} : mot français (inventé par les modèles)
\providecommand{\fr}[1]{\textit{#1}}
% zhbox : encadré de texte chinois inventé par les modèles
\newenvironment{zhbox}{\begin{quote}}{\end{quote}}
% \courssubsub : titre de niveau 3 inventé par les modèles
\providecommand{\courssubsub}[1]{\par\medskip\noindent\textbf{#1}\par\smallskip}
% \zhbf : texte chinois en gras inventé par les modèles
\providecommand{\zhbf}[1]{\textbf{#1}}
% \vs : « versus » inventé par les modèles
\providecommand{\vs}{\textit{vs}\ }
% \blank : case à compléter inventée par les modèles
\providecommand{\blank}[1][]{\underline{\hspace{1.6em}}}
\newtcolorbox{exemplebox}[1][]{popbase,colframe=poporange,before upper={\popboxhead{Exemple}{poporange}{#1}}}
\newtcolorbox{definition}[1][]{popbase,colframe=popblue,before upper={\popboxhead{Définition}{popblue}{#1}}}
\newtcolorbox{propriete}[1][]{popbase,colframe=poppurple,before upper={\popboxhead{Propriété}{poppurple}{#1}}}
\newtcolorbox{methode}[1][]{popbase,colframe=popgold,before upper={\popboxhead{Méthode}{popgold}{#1}}}
\newtcolorbox{attention}[1][]{popbase,colframe=poppink,before upper={\popboxhead{Attention}{poppink}{#1}}}
\newtcolorbox{experience}[1][]{popbase,colframe=popblue,colback=popblueL,before upper={\popboxhead{Expérience}{popblue}{#1}}}
% « Le savais-tu ? » : carte violette pleine (contexte, culture, Cameroun)
\newtcolorbox{savaistu}[1][]{popbase,colback=poppurple,colframe=popink,
  fontupper=\popsemi\fontsize{10.4}{14.2}\selectfont\color{white},
  before upper={\poppill{Le savais-tu\,?}{white}{poppurple}\ifx\relax#1\relax\else\hspace{6pt}{\popxbold\color{white}#1}\fi\par\vspace{7pt}}}
% Exercice résolu : énoncé en haut, \tcblower puis la solution sur fond crème
\newcounter{popexo}\newcounter{popexr}
\newtcolorbox{exoresolu}[1][]{popbase,colframe=poporange,colbacklower=popcream,
  segmentation style={popink,line width=1.2pt,dashed},
  before upper={\stepcounter{popexr}\popboxhead{Exercice résolu \thepopexr}{poporange}{#1}},
  before lower={\poppill{Solution}{popink}{popcream}\par\vspace{6pt}}}
% Exercice (non résolu en place) — la correction est dans la partie Corrigés
\newtcolorbox{exercice}[1][]{popbase,colframe=popink,
  before upper={\stepcounter{popexo}\popboxhead{Exercice \thepopexo}{popink}{#1}}}
% Corrigé
\newtcolorbox{corrige}[1][]{popbase,colframe=popgreen,colback=popgreenL,
  before upper={\popboxhead{Corrigé}{popgreen}{#1}}}
% Objectifs / prérequis / bilan
\newtcolorbox{objectifs}{popbase,colframe=popink,colback=poporangeL,
  before upper={\popboxhead{Objectifs de la leçon}{popink}{}}}
\newtcolorbox{prerequis}{popbase,colframe=popink,colback=popblueL,
  before upper={\popboxhead{Prérequis}{popblue}{}}}
\newtcolorbox{bilan}{popbase,colframe=popink,colback=white,boxrule=2.8pt,
  before upper={\popboxhead{Fiche bilan}{poporange}{L'essentiel à connaître pour l'examen}}}
% Figure encadrée avec légende
\newtcolorbox{popfigure}[1][]{enhanced,breakable=false,colback=white,colframe=popline,boxrule=1.4pt,arc=8pt,
  left=10pt,right=10pt,top=10pt,bottom=8pt,before skip=10pt,after skip=12pt,halign=center,
  lower separated=false,halign lower=center,
  fontlower=\popsemi\fontsize{9}{11.5}\selectfont\color{popmuted},
  #1}
\newcommand{\legende}[1]{\par\vspace{6pt}{\popsemi\fontsize{9}{11.5}\selectfont\color{popmuted}#1}}

% ============================================================
% Signature OnBuch+
% ============================================================
\newcommand{\poplogoinline}[1]{{\poplogo\fontsize{#1}{#1}\selectfont\color{popink}OnBuch\color{poporange}+}}
\newcommand{\signatureonbuch}{%
  \begin{tcolorbox}[enhanced,colback=popink,colframe=popink,boxrule=2.2pt,arc=10pt,
      drop shadow=poporange,left=14pt,right=14pt,top=10pt,bottom=10pt,before skip=0pt,after skip=0pt]
    \tikz[baseline=-0.7ex]\node[circle,fill=popgreen,draw=popcream,line width=1.3pt,minimum size=22pt,inner sep=0]{\popcheck{white}};%
    \hspace{9pt}\begin{minipage}[c]{0.62\linewidth}
      {\popxbold\fontsize{12}{13}\selectfont\color{popcream}Signé\ \textbullet\ {\poplogo OnBuch\color{poporange}+}}\par\vspace{2pt}
      {\raggedright\popsemi\fontsize{8}{10.5}\selectfont\color{popline}\MakeUppercase{Équipe pédagogique OnBuch+}\\\MakeUppercase{Conforme au programme officiel MINESEC}\par}
    \end{minipage}\hfill
    {\poplogo\fontsize{22}{22}\selectfont\color{popcream}OnBuch\color{poporange}+}
  \end{tcolorbox}%
}

% ============================================================
% Page de garde
% #1 titre · #2 matière · #3 niveau · #4 module · #5 n° leçon · #6 durée ·
% #7 description / objectifs (texte ou itemize)
% ============================================================
\newcommand{\pagegarde}[7]{%
  \thispagestyle{empty}
  \newgeometry{a4paper,margin=1.7cm,top=1.6cm,bottom=1.4cm}%
  \begin{tikzpicture}[remember picture,overlay]
    \fill[poporange!18] ([xshift=-3.2cm,yshift=-3.6cm]current page.north east) circle (4.6cm);
    \fill[poppurple!14] ([xshift=2.4cm,yshift=7.6cm]current page.south west) circle (3.8cm);
  \end{tikzpicture}%
  {\poplogo\fontsize{30}{30}\selectfont\color{popink}OnBuch\color{poporange}+}\hfill
  \raisebox{0.6ex}{\poppill{Cours signé \textbullet\ Premium}{popink}{popcream}}\par
  \vspace{1pt}\popsquiggle{poppurple}\par
  \vspace{8pt}
  \begin{tcolorbox}[enhanced,colback=poporange,colframe=popink,boxrule=2.8pt,arc=13pt,
      drop shadow=popink,left=18pt,right=18pt,top=12pt,bottom=13pt,after skip=10pt]
    \poppill{#2 \textbullet\ #3}{popink}{popcream}\hspace{5pt}\poppill{#4}{white}{popink}\par\vspace{8pt}
    {\popdisplay\fontsize{14}{14}\selectfont\color{popink}LEÇON #5}\par\vspace{4pt}
    {\raggedright\hyphenpenalty=10000\exhyphenpenalty=10000\popdisplay\fontsize{25}{27}\selectfont\color{popcream}\MakeUppercase{#1}\par}
    \vspace{8pt}
    {\popxbold\fontsize{9.5}{9.5}\selectfont\color{popink}\MakeUppercase{Durée conseillée : #6}}
  \end{tcolorbox}
  \begin{tcolorbox}[enhanced,colback=white,colframe=popink,boxrule=2.2pt,arc=10pt,
      drop shadow=popink,left=16pt,right=16pt,top=10pt,bottom=10pt,after skip=8pt]
    \poppill{Dans cette leçon}{poppurple}{white}\par\vspace{5pt}
    \popsemi\fontsize{9.3}{12.6}\selectfont\color{popink}#7
  \end{tcolorbox}
  \noindent\begin{minipage}[t]{0.487\linewidth}
    \begin{tcolorbox}[enhanced,colback=popgreen,colframe=popink,boxrule=2.2pt,arc=10pt,
        drop shadow=popink,left=11pt,right=11pt,top=10pt,bottom=10pt,height=3.1cm,valign=center]
      \tcbox[colback=white,colframe=popink,boxrule=1.3pt,arc=5pt,boxsep=0pt,nobeforeafter,
        left=3pt,right=3pt,top=3pt,bottom=3pt,valign=center]{\qrcode[height=1.9cm]{https://play.google.com/store/apps/details?id=com.luvvix.onbuch&hl=fr}}%
      \hspace{8pt}\begin{minipage}[c]{0.5\linewidth}
        {\popdisplay\fontsize{11}{12.5}\selectfont\color{white}\MakeUppercase{Télécharge OnBuch}}\par\vspace{4pt}
        {\raggedright\popsemi\fontsize{8.4}{11}\selectfont\color{white}Gratuit \textbullet\ Google Play\\Tuteur IA, annales, concours, résultats\par}
      \end{minipage}
    \end{tcolorbox}
  \end{minipage}\hfill
  \begin{minipage}[t]{0.487\linewidth}
    \begin{tcolorbox}[enhanced,colback=poppurple,colframe=popink,boxrule=2.2pt,arc=10pt,
        drop shadow=popink,left=11pt,right=11pt,top=10pt,bottom=10pt,height=3.1cm,valign=center]
      \tikz[baseline=-0.7ex]\node[circle,fill=white,draw=popink,line width=1.3pt,minimum size=1.6cm,inner sep=0]{\popdisplay\fontsize{24}{24}\selectfont\color{poppurple}+};%
      \hspace{8pt}\begin{minipage}[c]{0.6\linewidth}
        {\popdisplay\fontsize{11}{12.5}\selectfont\color{white}\MakeUppercase{Passe à OnBuch+}}\par\vspace{4pt}
        {\raggedright\popsemi\fontsize{8.4}{11}\selectfont\color{white}Tous les cours signés, Tuteur IA illimité, fiches PDF — onglet OnBuch+ de l'app\par}
      \end{minipage}
    \end{tcolorbox}
  \end{minipage}
  \vspace{8pt}
  \popcontenu
  \vfill
  \signatureonbuch
  \restoregeometry
  \newpage
}

% Rangée « Ce cours contient » (page de garde)
\newcommand{\poptile}[3]{% #1 couleur #2 gros texte #3 légende
  \begin{tcolorbox}[enhanced,colback=#1,colframe=popink,boxrule=2pt,arc=9pt,shadow={2.5pt}{-2.5pt}{0pt}{popink},
      left=6pt,right=6pt,top=9pt,bottom=9pt,halign=center,valign=center,height=1.75cm,width=\linewidth,nobeforeafter]
    {\popdisplay\fontsize{12.5}{13}\selectfont\color{white}\MakeUppercase{#2}}\par\vspace{3pt}
    {\centering\popsemi\fontsize{7.8}{9.5}\selectfont\color{white}#3\par}
  \end{tcolorbox}}
\newcommand{\popcontenu}{%
  \par\noindent{\popxboldls\fontsize{7.5}{7.5}\selectfont\color{popmuted}CE COURS SIGNÉ CONTIENT}\par\vspace{7pt}
  \noindent\begin{minipage}[t]{0.235\linewidth}\poptile{popblue}{Cours}{complet, illustré, pas à pas}\end{minipage}\hfill
  \begin{minipage}[t]{0.235\linewidth}\poptile{poporange}{Méthodes}{et exercices résolus}\end{minipage}\hfill
  \begin{minipage}[t]{0.235\linewidth}\poptile{poppink}{Exercices}{type examen + corrigés}\end{minipage}\hfill
  \begin{minipage}[t]{0.235\linewidth}\poptile{popgreen}{Fiche bilan}{l'essentiel à retenir}\end{minipage}\par}

% Sommaire
\newtcolorbox{sommaire}{enhanced,colback=white,colframe=popink,boxrule=2.2pt,arc=10pt,
  drop shadow=popink,left=18pt,right=18pt,top=13pt,bottom=13pt,before skip=6pt,after skip=20pt,
  before upper={\poppill{Sommaire}{poppurple}{white}\par\vspace{9pt}\popsemi\fontsize{10.6}{14.5}\selectfont\color{popink}}}

% Signature de fin de cours — poussée tout en bas de la dernière page
\newcommand{\findecours}{%
  \par\vfill\nopagebreak
  \begin{center}\popsquiggle{poporange}\end{center}\vspace{4pt}
  \signatureonbuch
}
\AtBeginDocument{\def\llbracket{\mathopen{[\mkern-3.5mu[}}\def\rrbracket{\mathclose{]\mkern-3.5mu]}}} % intervalles d'entiers (glyphe ⟦ absent de la police)
% Glyphes absents de Fira Math : redéfinis à partir de glyphes disponibles
\AtBeginDocument{%
  \def\setminus{\mathbin{\backslash}}\let\smallsetminus\setminus
  \def\vdots{\mathord{\vbox{\baselineskip3.2pt\lineskiplimit0pt\kern4pt\hbox{.}\hbox{.}\hbox{.}}}}%
  \def\ddots{\mathinner{\mkern1mu\raise6.5pt\hbox{.}\mkern2mu\raise3.6pt\hbox{.}\mkern2mu\raise0.7pt\hbox{.}\mkern1mu}}%
  \def\triangle{\mathord{\scalebox{0.9}{$\symup{\Delta}$}}}%
  \def\nabla{\mathord{\raisebox{\depth}{\scalebox{1}[-1]{$\symup{\Delta}$}}}}%
  \def\checkmark{\popcheck{popdark}}%
  \def\star{\mathbin{\ast}}\def\bigstar{\mathord{\ast}}%
  \def\diamond{\mathbin{\scalebox{0.6}{$\Diamond$}}}%
  \def\bigcup{\mathop{\mathchoice{\raisebox{-0.3em}{\scalebox{1.7}{$\cup$}}}{\scalebox{1.25}{$\cup$}}{\cup}{\cup}}\displaylimits}%
  \def\bigcap{\mathop{\mathchoice{\raisebox{-0.3em}{\scalebox{1.7}{$\cap$}}}{\scalebox{1.25}{$\cap$}}{\cap}{\cap}}\displaylimits}%
  \def\lesseqqgtr{\mathrel{\lessgtr}}\def\bigoplus{\mathop{\oplus}\displaylimits}%
}
\providecommand{\ssi}{si et seulement si} % abréviation fréquente
\AtBeginDocument{\providecommand{\wideparen}{\overparen}} % arc de cercle

\begin{document}

\pagegarde{Leçon 39 — Vocabulaire et compréhension de texte : rédaction d’une lettre formelle}{Chinois}{1ère A}{Module 5 : Mass médias et télécommunication}{ZHA39}{2 h}{Cette leçon propose l'étude complète d'un modèle de lettre formelle chinoise : texte support bilingue, analyse lexicale détaillée, décomposition graphique des caractères clés, structure rhétorique type et mise en pratique guidée. Elle prépare l'élève à l'examen du Baccalauréat en développant les compétences de lecture, d'analyse linguistique et de production écrite administrative.
\vspace{4pt}
\begin{multicols}{2}\small
\begin{itemize}
\item À la fin de cette leçon, je sais identifier et reproduire la structure standard en 6 parties d'une lettre formelle chinoise (en-tête, appel, corps, formule de politesse, signature, date).
\item Je sais reconnaître, prononcer correctement (pinyin et tons) et écrire les 20 mots de vocabulaire administratif et épistolaire du thème.
\item Je sais analyser la composition graphique (radicaux, ordre des traits) de 8 caractères clés pour en faciliter la mémorisation et la reconnaissance.
\item Je sais expliquer la fonction des formules de politesse spécifiques (敬语) comme 贵公司, 敬启, 顺颂商祺 et les distinguer du langage courant.
\item Je sais répondre par écrit en chinois à des questions de compréhension globale et détaillée sur le texte support.
\item Je sais employer les connecteurs logiques formels (如 贵方, 关于, 烦请, 专此) dans des phrases types de correspondance.
\end{itemize}\end{multicols}}

\begin{sommaire}
\begin{multicols}{2}
\begin{enumerate}
\item 1. Mise en situation \& Texte support : Lettre de demande de stage
\item 2. Glossaire thématique : Vocabulaire de la correspondance administrative
\item 3. Atelier Caractères : Radicaux, Ordre des traits \& Mémorisation
\item 4. Architecture rhétorique : Connecteurs \& Structures types du corps de lettre
\item 5. Compréhension de texte \& Collocations : Du passif à l'actif
\item 6. Production guidée : Atelier d'écriture 'Ma première lettre officielle'
\item 7. Extension Bac \& Évaluation formative (Complément examen)
\item Activité d'intégration
\item Méthodes et exercices résolus
\item Exercices
\item Corrigés des exercices
\item Fiche bilan
\end{enumerate}
\end{multicols}
\end{sommaire}

\begin{objectifs}
\begin{itemize}
\item À la fin de cette leçon, je sais identifier et reproduire la structure standard en 6 parties d'une lettre formelle chinoise (en-tête, appel, corps, formule de politesse, signature, date).
\item Je sais reconnaître, prononcer correctement (pinyin et tons) et écrire les 20 mots de vocabulaire administratif et épistolaire du thème.
\item Je sais analyser la composition graphique (radicaux, ordre des traits) de 8 caractères clés pour en faciliter la mémorisation et la reconnaissance.
\item Je sais expliquer la fonction des formules de politesse spécifiques (敬语) comme 贵公司, 敬启, 顺颂商祺 et les distinguer du langage courant.
\item Je sais répondre par écrit en chinois à des questions de compréhension globale et détaillée sur le texte support.
\item Je sais employer les connecteurs logiques formels (如 贵方, 关于, 烦请, 专此) dans des phrases types de correspondance.
\item Je sais rédiger une lettre formelle complète et cohérente (demande de stage, réclamation, demande d'information) en respectant les conventions de mise en page chinoises.
\end{itemize}
\end{objectifs}

\begin{prerequis}
\begin{itemize}
\item Maîtrise des 300 caractères HSK 1-2 (niveau Seconde) et notions de base de l'ordre des traits.
\item Connaissance des structures de base : 是, 有, 在, 想/要/会, 了/过/着.
\item Notions de politesse de base : 您, 请, 谢谢, 对不起, et distinction 你/您.
\item Capacité à rédiger un paragraphe simple en chinois (80-100 caractères) avec connecteurs simples (因为...所以, 虽然...但是).
\item Connaissance culturelle : notion de 'Face' (面子) et importance de la hiérarchie en contexte sinophone.
\end{itemize}
\end{prerequis}

\newpage

\coursec{Mise en situation \& Texte support : Lettre de demande de stage}

\courssub{Situation d'accroche et problématique}

Imagine que tu souhaites postuler à une bourse d'études à l'Université de Pékin (\zh{北京大学}) ou écrire à l'Ambassade de Chine à Yaoundé (\zh{中国驻喀麦隆大使馆}) pour obtenir un visa d'affaires (\zh{商务签证}). Au Cameroun, comme dans tout contexte administratif international, la maîtrise des codes de la correspondance formelle en chinois (\zh{正式信函写作}) est une clé qui ouvre les portes. Contrairement à un message WeChat (\zh{微信}) entre amis, une lettre formelle exige une structure rigide (\zh{格式固定}), un vocabulaire spécifique (\zh{敬语} \emph{jìngyǔ}, « termes de respect ») et un respect absolu de la hiérarchie (\zh{尊卑有序}). Une maladresse dans l'appel (\zh{称呼}) ou la formule de politesse finale (\zh{结语}) peut faire échouer ta démarche avant même qu'elle ne soit lue.

Dans cette leçon, nous suivons \textbf{Wang Wei} (\zh{王伟}), un élève camerounais de la section bilingue du Lycée de Yaoundé (\zh{雅温得第一中学}), qui écrit à une entreprise technologique de Shenzhen (\zh{深圳科技有限公司}) pour solliciter un stage d'été (\zh{暑期实习}). À travers sa lettre, tu vas apprendre à repérer les six blocs structurels d'une lettre formelle chinoise, à manier les marques de déférence et à éviter les pièges typiques des francophones.

\textbf{Problématique :} comment structurer et formuler une lettre administrative en chinois pour qu'elle soit lue, comprise et prise au sérieux ?

\courssub{Texte support : la lettre de Wang Wei}

\begin{textebox}[文本内容 : Lettre de demande de stage à Shenzhen]
\textbf{Texte chinois :}\\
招聘部：\\
\u3000\u3000贵公司：\\
\u3000\u3000敬启！\\
\u3000\u3000我是喀麦隆雅温得第一中学高二(1)班的学生王伟。关于贵公司招聘暑期实习生一事，我非常感兴趣。我学习刻苦，成绩优异，且通过了HSK四级考试，能用中文进行日常交流。烦请贵方予以考虑。\\
\u3000\u3000顺颂\\
\u3000\u3000商祺！\\
\u3000\u3000\u3000\u3000\u3000\u3000王伟 启\\
\u3000\u3000\u3000\u3000\u3000\u30002025年10月15日\\[4pt]
\textbf{Pinyin :}\\
Zhāopìn bù : Guì gōngsī : Jìng qǐ ! Wǒ shì Kāmàilóng Yǎwēndé Dì-yī Zhōngxué gāo èr (yī) bān de xuésheng Wáng Wěi. Guānyú guì gōngsī zhāopìn shǔqī shíxíshēng yī shì, wǒ fēicháng gǎn xìngqù. Wǒ xuéxí kèkǔ, chéngjì yōuyì, qiě tōngguò le HSK sì jí kǎoshì, néng yòng Zhōngwén jìnxíng rìcháng jiāoliú. Fán qǐng guìfāng yǔyǐ kǎolǜ. Shùn sòng shāng qí ! Wáng Wěi qǐ. Èr líng èr wǔ nián shí yuè shíwǔ rì.\\[4pt]
\textbf{Traduction française :}\\
Service du recrutement : À votre honorable société : Salutations respectueuses ! Je suis Wang Wei, élève de la classe de Première (1) du Premier Lycée de Yaoundé, au Cameroun. Concernant le recrutement de stagiaires d'été par votre honorable société, je suis très intéressé. J'étudie avec assiduité, j'obtiens d'excellents résultats, et j'ai réussi l'examen HSK niveau 4 ; je peux communiquer au quotidien en chinois. Je vous prie respectueusement de bien vouloir prendre ma candidature en considération. Avec mes salutations distinguées (formule commerciale) ! Wang Wei, respectueusement. Le 15 octobre 2025.
\end{textebox}

\textbf{Lecture guidée.} L'enseignant lit d'abord la lettre à voix haute, en marquant les tons et les pauses (groupes respiratoires). Les élèves relisent ensuite en chorale, phrase par phrase, puis individuellement. Pendant la lecture, observe : où commence chaque partie ? Quels mots reviennent pour exprimer le respect ? Pourquoi \zh{贵} apparaît-il deux fois (\zh{贵公司}, \zh{贵方}) ?

\courssub{Observation : les six blocs structurels}

Avant toute règle, observons la lettre. Elle n'est pas un bloc unique : elle se découpe en \textbf{six blocs} toujours présents, dans un ordre immuable (de haut en bas).

\begin{definition}[格式 : Le format de la lettre formelle]
Une lettre formelle chinoise (\zh{正式信函} \emph{zhèngshì xìnhán}) suit une mise en page codifiée : pas de ligne blanche entre l'appel et le corps, indentation de deux caractères chinois (\u3000\u3000, équivalent à deux espaces pleins) pour chaque paragraphe du corps, formules de politesse finales alignées à gauche sur deux lignes distinctes (ex. \zh{顺颂} / \zh{商祺}), signature et date rejetées à droite.
\end{definition}

\begin{center}
\begin{tabularx}{\linewidth}{|X|X|X|X|}
\hline
\rowcolor{popblueL} Bloc & Chinois (pinyin) & Contenu dans la lettre & Équivalent français \\
\hline
1. En-tête & 抬头 \emph{táitóu} & 招聘部：\newline 贵公司： & « À l'attention de... » \\
\hline
2. Appel & 称呼 \emph{chēnghu} & 敬启！ & « Monsieur, Madame, » (solennel) \\
\hline
3. Corps & 正文 \emph{zhèngwén} & 我是...予以考虑。 & Corps de la lettre (indenté) \\
\hline
4. Formule finale & 敬语/结语 \emph{jìngyǔ/jiéyǔ} & 顺颂\newline 商祺！ & « Veuillez agréer... » (commercial) \\
\hline
5. Signature & 落款 \emph{luòkuǎn} & 王伟 启 & Nom + formule de clôture \\
\hline
6. Date & 日期 \emph{rìqī} & 2025年10月15日 & Date (année-mois-jour) \\
\hline
\end{tabularx}
\end{center}

\begin{popfigure}
\begin{tikzpicture}[
    node distance=1.2cm,
    box/.style={fill=popblueL, draw=popblue, thick, rounded corners, inner sep=6pt, align=center, text width=5cm, font=\small},
    arrow/.style={->, thick, draw=popdark}
]
\node[box, align=center] (h){1. En-tête (\emph{táitóu})\\ \zh{招聘部：贵公司：}};
\node[box, align=center] (c1) [below of=h]{2. Appel (\emph{chēnghu})\\ \zh{敬启！}};
\node[box, align=center] (c2) [below of=c1]{3. Corps (\emph{zhèngwén})\\ \zh{我是...予以考虑。} \emph{(indenté)}};
\node[box, align=center] (c3) [below of=c2]{4. Formule finale (\emph{jiéyǔ})\\ \zh{顺颂 商祺！}};
\node[box, align=center] (c4) [below of=c3]{5. Signature (\emph{luòkuǎn})\\ \zh{王伟 启} \emph{(à droite)}};
\node[box, align=center] (c5) [below of=c4]{6. Date (\emph{rìqī})\\ \zh{2025年10月15日} \emph{(à droite)}};
\draw[arrow] (h) -- (c1);
\draw[arrow] (c1) -- (c2);
\draw[arrow] (c2) -- (c3);
\draw[arrow] (c3) -- (c4);
\draw[arrow] (c4) -- (c5);
\end{tikzpicture}
\legende{Les six blocs structurels d'une lettre formelle chinoise, dans l'ordre obligatoire de haut en bas.}
\end{popfigure}

\courssub{La date chinoise : 年 / 月 / 日}

Observe la dernière ligne : \zh{2025年10月15日} (\emph{èr líng èr wǔ nián shí yuè shíwǔ rì}). Le chinois procède \textbf{du plus grand au plus petit} : année (\zh{年}) → mois (\zh{月}) → jour (\zh{日}). C'est l'inverse de l'habitude française « 15 octobre 2025 ».

\begin{propriete}[Structure de la date]
Schéma : \cle{chiffre + 年 + chiffre + 月 + chiffre + 日}.\\
Exemple : \zh{2025年10月15日} \emph{èr líng èr wǔ nián shí yuè shíwǔ rì} « le 15 octobre 2025 ».\\
En chinois parlé, on lit les années chiffre par chiffre (\zh{二零二五年}), jamais « deux mille vingt-cinq ». Le caractère \zh{日} (jour) est obligatoire dans l'écrit formel ; \zh{号} (\emph{hào}) s'utilise à l'oral et dans les écrits informels (SMS, agenda).
\end{propriete}

\textbf{Alignement.} Deux formats coexistent : le \textbf{format bloc moderne} (signature et date alignées à droite, comme dans notre lettre) et le \textbf{format traditionnel} où l'appel peut être surélevé (\zh{抬头}, littéralement « tête levée » : on laissait un blanc ou on montait d'une ligne devant le nom du destinataire pour l'honorer). En Première, retiens le format bloc : appel à gauche au bord, corps indenté de deux caractères, signature et date à droite.

\coursec{Glossaire thématique : Vocabulaire de la correspondance administrative}

\begin{center}
\begin{tabularx}{\linewidth}{|X|X|X|X|}
\hline
\rowcolor{popblueL} Caractères & Pinyin & Nature & Traduction \\
\hline
招聘 & zhāopìn & verbe/nom & recruter / recrutement \\
\hline
部 & bù & nom (suffixe) & service, département \\
\hline
贵 & guì & adj. (honorifique) & honorable (préfixe pour l'interlocuteur) \\
\hline
公司 & gōngsī & nom & société, entreprise \\
\hline
敬启 & jìng qǐ & locution & salutations respectueuses (appel) \\
\hline
喀麦隆 & Kāmàilóng & nom propre & Cameroun \\
\hline
雅温得 & Yǎwēndé & nom propre & Yaoundé \\
\hline
第一中学 & Dì-yī Zhōngxué & nom propre & Premier Lycée (Lycée n°1) \\
\hline
高二 & gāo èr & nom & classe de Première (2e année lycée) \\
\hline
关于 & guānyú & préposition & à propos de, concernant \\
\hline
暑期 & shǔqī & nom & période estivale, vacances d'été \\
\hline
实习生 & shíxíshēng & nom & stagiaire \\
\hline
一事 & yī shì & nom & cette affaire, ce dossier \\
\hline
感兴趣 & gǎn xìngqù & verbe (séparable) & s'intéresser à (\zh{对...感兴趣}) \\
\hline
刻苦 & kèkǔ & adj. & assidu, acharné (au travail/études) \\
\hline
优异 & yōuyì & adj. & excellent, brillant (résultats) \\
\hline
且 & qiě & conjonction & de plus, et (formel, écrit) \\
\hline
通过 & tōngguò & verbe & réussir (examen), passer par \\
\hline
日常交流 & rìcháng jiāoliú & nom & communication quotidienne \\
\hline
烦请 & fán qǐng & verbe (formel) & je vous prie (litt. « déranger pour demander ») \\
\hline
贵方 & guìfāng & nom (honorifique) & votre partie / votre côté (votre entreprise) \\
\hline
予以 & yǔyǐ & verbe (formel) & accorder, donner (structure \zh{予以+V}) \\
\hline
考虑 & kǎolǜ & verbe & considérer, examiner \\
\hline
顺颂 & shùn sòng & locution & profiter de l'occasion pour souhaiter \\
\hline
商祺 & shāng qí & nom & prospérité des affaires \\
\hline
启 & qǐ & verbe (formel) & respectueusement (après signature, inf.→sup.) \\
\hline
此致 & cǐ zhì & locution & c'est tout / par la présente (formule finale alt.) \\
\hline
敬礼 & jìnglǐ & verbe/nom & salutations respectueuses (formule finale alt.) \\
\hline
尊敬的 & zūnjìng de & adj. & honoré, respecté (devant titre/nom) \\
\hline
\end{tabularx}
\end{center}

\coursec{Atelier Caractères : Radicaux, Ordre des traits \& Mémorisation}

\begin{definition}[Clés (radicaux) et structure des caractères clés]
Maîtriser l'écriture des caractères formels permet de gagner en lisibilité et en crédibilité. Voici l'analyse de 5 caractères à haute fréquence dans ce thème.
\end{definition}

\begin{center}
\begin{tabularx}{\linewidth}{|X|X|X|X|X|}
\hline
\rowcolor{popblueL} Caractère & Clé (Radical) & Nb traits & Ordre des traits (règles principales) & Astuce mnémotechnique \\
\hline
\zh{敬} \emph{jìng} & \zh{攵} (pū, « frapper/action ») & 13 & 1. Haut \zh{苟} (herbe + deux mains) 2. Bas \zh{攵} (action respectueuse) & « Agir (\zh{攵}) avec la plante du respect sur la tête » \\
\hline
\zh{启} \emph{qǐ} & \zh{户} (hù, « porte ») & 7 & 1. \zh{户} (cadre porte) 2. Traits intérieurs (deux points, horizontale, verticale) & « Ouvrir la porte (\zh{户}) pour commencer » \\
\hline
\zh{贵} \emph{guì} & \zh{贝} (bèi, « coquillage/monnaie ») & 12 & 1. Haut \zh{⺁} + \zh{口} + \zh{田} 2. Bas \zh{贝} & « Deux mains (\zh{⺁}) tiennent un trésor (\zh{贝}) = précieux/honorable » \\
\hline
\zh{颂} \emph{sòng} & \zh{页} (yè, « tête/page ») & 13 & 1. Gauche \zh{公} (public) 2. Droite \zh{页} (tête) & « Louer publiquement (\zh{公}) de tête (\zh{页}) » \\
\hline
\zh{烦} \emph{fán} & \zh{火} (huǒ, « feu ») & 13 & 1. Bas \zh{火} (4 traits) 2. Haut \zh{页} (9 traits) & « Le feu (\zh{火}) brûle la tête (\zh{页}) = ennui, déranger » \\
\hline
\end{tabularx}
\end{center}

\begin{methode}[Écrire \zh{敬} et \zh{启} correctement]
\begin{enumerate}
\item \textbf{\zh{敬} (jìng)} : Respecter l'ordre « haut → bas ». Écrire le composant supérieur \zh{苟} (herbe \zh{艹} + \zh{句}) complètement avant le radical \zh{攵} (point, horizontale, point, trait oblique). Ne pas confondre \zh{攵} (4 traits) et \zh{文} (4 traits mais structure différente).
\item \textbf{\zh{启} (qǐ)} : Règle « extérieur avant intérieur » pour le radical \zh{户}. Écrire le cadre \zh{户} (verticale, horizontale courbe, verticale) PUIS les traits intérieurs (deux petits points, une horizontale, une verticale finale qui ne perce pas le bas).
\end{enumerate}
\end{methode}

\begin{experience}[Dictée progressive \& Auto-correction]
L'enseignant dicte : 1. \zh{贵公司} 2. \zh{敬启} 3. \zh{烦请} 4. \zh{顺颂商祺} 5. \zh{予以考虑}. Les élèves écrivent sur ardoise. Correction collective : focus sur l'ordre des traits de \zh{敬} et \zh{启}, et la distinction \zh{贵} vs \zh{遗} (yí, perdre).
\end{experience}

\coursec{Architecture rhétorique : Connecteurs \& Structures types du corps de lettre}

Le corps de la lettre (\zh{正文}) suit une logique rhétorique standardisée : \textbf{Identité} → \textbf{Objet} → \textbf{Arguments} → \textbf{Demande}. Les connecteurs (mots-outils) assurent la cohésion.

\begin{propriete}[Connecteurs logiques du style administratif]
\begin{itemize}
\item \zh{关于} (\emph{guānyú}) + [Sujet] + \zh{一事/一节} : « Concernant [Sujet]... » (Introduit l'objet). Ex: \zh{关于贵公司招聘暑期实习生一事}。
\item \zh{且} (\emph{qiě}) : « De plus, par ailleurs » (écrit, formel). Relie deux qualités/arguments. Ex: \zh{学习刻苦，成绩优异，且通过了HSK四级}。
\item \zh{烦请} (\emph{fán qǐng}) + [Destinataire honorifique \zh{贵方/您}] + \zh{予以} (\emph{yǔyǐ}) + [Verbe d'action] : Structure figée de demande polie. « Je vous prie de bien vouloir [Verbe] ». Ex: \zh{烦请贵方予以考虑}。
\end{itemize}
\end{propriete}

\begin{center}
\begin{tabularx}{\linewidth}{|X|X|X|}
\hline
\rowcolor{popblueL} Structure type & Exemple (Caractères + Pinyin) & Traduction \\
\hline
\zh{我是...} (Identité) & \zh{我是喀麦隆雅温得第一中学学生王伟。} \emph{Wǒ shì Kāmàilóng Yǎwēndé Dì-yī Zhōngxué xuésheng Wáng Wěi.} & Je suis Wang Wei, élève du Lycée n°1 de Yaoundé, Cameroun. \\
\hline
\zh{关于...一事} (Objet) & \zh{关于贵公司招聘暑期实习生一事，...} \emph{Guānyú guì gōngsī zhāopìn shǔqī shíxíshēng yī shì,...} & Concernant le recrutement de stagiaires d'été par votre société,... \\
\hline
\zh{我...，且...} (Arguments) & \zh{我学习刻苦，成绩优异，且通过了HSK四级考试。} \emph{Wǒ xuéxí kèkǔ, chéngjì yōuyì, qiě tōngguò le HSK sì jí kǎoshì.} & J'étudie assidûment, mes résultats sont excellents, et j'ai réussi le HSK 4. \\
\hline
\zh{烦请...予以...} (Demande) & \zh{烦请贵方予以考虑。} \emph{Fán qǐng guìfāng yǔyǐ kǎolǜ.} & Je vous prie de bien vouloir examiner ma candidature. \\
\hline
\end{tabularx}
\end{center}

\begin{attention}[Pièges syntaxiques pour francophones]
\begin{enumerate}
\item \textbf{Place de \zh{关于}} : En français « Concernant X, je... ». En chinois \zh{关于X，我...} (Sujet \zh{我} après la virgule). Ne pas dire \zh{关于贵公司，我感兴趣} (trop vague), préférer \zh{关于...一事}。
\item \textbf{\zh{且} vs \zh{和/跟}} : \zh{且} lie des propositions (sujet + verbe), style écrit. \zh{和/跟} lient des noms. \textbf{Erreur :} \zh{我刻苦和成绩好} → \textbf{Correct :} \zh{我学习刻苦，成绩优异} ou \zh{我既刻苦，成绩又优异}。
\item \textbf{\zh{予以} + Verbe} : \zh{予以} est un mot-outil formel qui « transpose » le verbe suivant. \zh{予以考虑} = « donner considération » = « considérer ». Ne pas ajouter \zh{给} après : \zh{*予以给考虑} (faux).
\item \textbf{Objet de \zh{感兴趣}} : Structure \zh{对...感兴趣}. Dans la lettre, \zh{对} est omis car l'objet est le sujet du thème (\zh{关于...一事}), mais à l'oral : \zh{我对贵公司的业务很感兴趣}。
\end{enumerate}
\end{attention}

\coursec{Compréhension de texte \& Collocations : Du passif à l'actif}

\courssub{Exercices de compréhension (QCM et Questions ouvertes)}

\begin{exercice}[Compréhension globale et détaillée]
\begin{enumerate}
\item \textbf{QCM :} Quel est le niveau d'études actuel de Wang Wei au moment de l'écriture ?
      A. \zh{高一} (Seconde) \quad B. \zh{高二} (Première) \quad C. \zh{高三} (Terminale) \quad D. \zh{大学生} (Étudiant)
\item \textbf{QCM :} Quelle formule indique que Wang Wei a un niveau de chinois certifié ?
      A. \zh{我学习刻苦} \quad B. \zh{成绩优异} \quad C. \zh{通过了HSK四级考试} \quad D. \zh{烦请贵方予以考虑}
\item \textbf{Question ouverte :} Relevez dans le texte les deux occurrences du préfixe honorifique \zh{贵} et précisez à quel nom elles se rapportent.
\item \textbf{Question ouverte :} Pourquoi Wang Wei utilise-t-il \zh{敬启} en appel et \zh{启} dans la signature ? Quelle est la différence de registre ?
\item \textbf{Traduction :} Traduisez en français : \zh{烦请贵方予以考虑。} Expliquez le rôle de \zh{予以}.
\end{enumerate}
\end{exercice}

\begin{corrige}[Exercice 5.1]
\begin{enumerate}
\item \textbf{B.} \zh{高二(1)班} = Première (classe de 1ère). Au Cameroun, 6ème-3ème = Collège (\zh{初中}), 2nde-Tle = Lycée (\zh{高中}: \zh{高一}, \zh{高二}, \zh{高三}).
\item \textbf{C.} \zh{通过了HSK四级考试} (a réussi l'examen HSK niveau 4).
\item \zh{贵公司} (votre honorable société) et \zh{贵方} (votre honorable partie/votre entreprise). Les deux renvoient à l'entreprise destinataire (\zh{深圳科技有限公司}).
\item \zh{敬启} est une formule d'\textbf{appel} (début de lettre), solennelle, adressée à une institution. \zh{启} (seul) en \textbf{signature} (fin de lettre) marque la soumission/respect de l'expéditeur (inférieur) vers le destinataire (supérieur).
\item « Je vous prie respectueusement de bien vouloir examiner [ma candidature]. » \zh{予以} (\emph{yǔyǐ}) est un verbe auxiliaire formel signifiant « accorder / donner » ; il précède un verbe d'action mentale ou administrative (\zh{考虑}, \zh{回复}, \zh{批准}, \zh{协助}) pour former une structure transitive : « accorder de la considération » = « considérer ».
\end{enumerate}
\end{corrige}

\courssub{Entraînement aux collocations (Associations lexicales)}

\begin{exercice}[Collocations : Reliez les colonnes pour former des expressions correctes]
\begin{center}
\begin{tabularx}{\linewidth}{|X|X|X|}
\hline
\rowcolor{popblueL} Colonne A (Mots-outils / Verbes) & Colonne B (Noms / Compléments) & Expression formée (à écrire en caractères) \\
\hline
\zh{关于} & \zh{贵公司 / 贵校 / 贵国} & \zh{关于贵公司一事} \\
\hline
\zh{烦请} & \zh{考虑 / 回复 / 批准 / 协助} & \zh{烦请予以考虑} \\
\hline
\zh{予以} & \zh{实习生 / 志愿者 / 奖学金} & \zh{招聘实习生} \\
\hline
\zh{招聘} & \zh{商祺 / 时祺 / 教祺 / 公祺} & \zh{顺颂商祺} \\
\hline
\zh{顺颂} & \zh{贵方 / 您 / 贵处} & \zh{烦请贵方} \\
\hline
\end{tabularx}
\end{center}
\end{exercice}

\begin{corrige}[Exercice 5.2]
\zh{关于贵公司一事} (ou \zh{贵校}, \zh{贵国}) ; \zh{烦请贵方} (ou \zh{您}, \zh{贵处}) ; \zh{予以考虑} (ou \zh{回复}, \zh{批准}, \zh{协助}) ; \zh{招聘实习生} (ou \zh{志愿者}...) ; \zh{顺颂商祺} (contexte commercial). Autres formules finales : \zh{顺颂时祺} (général), \zh{顺颂教祺} (éducation), \zh{此致敬礼} (administratif général).
\end{corrige}

\coursec{Production guidée : Atelier d'écriture « Ma première lettre officielle »}

\begin{methode}[Rédiger une lettre formelle en 5 étapes]
\begin{enumerate}
\item \textbf{Analyser la consigne :} Qui écrit à qui ? Quel objet ? Quel registre (commercial, scolaire, administratif) ?
\item \textbf{Poser la structure (squelette) :} Écrire les 6 blocs vides avec l'alignement correct (appel gauche, corps indenté, signature/date droite).
\item \textbf{Remplir l'en-tête et l'appel :} Choisir \zh{敬启} (institution) ou \zh{尊敬的[X]经理，您好！} (personne). Ajouter \zh{贵} devant l'entité destinataire.
\item \textbf{Rédiger le corps (3 paragraphes types) :}
      \begin{itemize}
      \item \zh{身份} : \zh{我是[Pays/École/Classe]的[Nom]。}
      \item \zh{动机/优势} : \zh{关于...一事，我...，且...。} (Utiliser \zh{关于}, \zh{且}, vocabulaire positif : \zh{刻苦}, \zh{优异}, \zh{熟练}).
      \item \zh{请求} : \zh{烦请贵方/您予以[考虑/回复/批准]。}
      \end{itemize}
\item \textbf{Choisir la formule finale et signer :} Commercial \zh{顺颂商祺} / Scolaire \zh{顺颂教祺} / Général \zh{此致敬礼}. Signature : \zh{[Nom] 启}. Date : \zh{XXXX年X月X日}.
\end{enumerate}
\end{methode}

\begin{experience}[Atelier d'écriture guidée : « Candidature spontanée »]
\textbf{Consigne :} Tu es \zh{李华} (\emph{Lǐ Huá}), élève de \zh{高二(2)班} au \zh{喀麦隆多拉第二中学} (Lycée de Douala 2). Tu souhaites écrire au \zh{中国驻喀麦隆大使馆文化处} (Service culturel de l'Ambassade de Chine au Cameroun) pour demander des informations sur les bourses d'études en Chine (\zh{中国政府奖学金}) pour l'année 2026.
\begin{enumerate}
\item Choisis la formule d'appel adaptée (institution).
\item Écris le corps de la lettre en utilisant : \zh{关于...一事}, \zh{我是...}, \zh{烦请...予以...}.
\item Choisis la formule de politesse finale adaptée (contexte diplomatique/éducatif).
\item Écris la lettre complète sur ta copie (caractères + pinyin). Respecte la mise en page (indentation, alignement).
\end{enumerate}
\end{experience}

\begin{exemplebox}[Modèle de réponse attendue (Corrigé type)]
\textbf{Chinois :}\\
中国驻喀麦隆大使馆文化处：\\
\u3000\u3000敬启！\\
\u3000\u3000我是喀麦隆多拉第二中学高二(2)班的学生李华。关于贵处发放2026年中国政府奖学金一事，我非常关注。我品学兼优，热爱中华文化，且通过了HSK三级考试。烦请贵方予以告知申请流程与材料清单。\\
\u3000\u3000顺颂\\
\u3000\u3000公祺！\\
\u3000\u3000\u3000\u3000\u3000\u3000李华 启\\
\u3000\u3000\u3000\u3000\u3000\u30002025年11月20日\\[4pt]
\textbf{Pinyin :} Zhōngguó zhù Kāmàilóng Dàshǐguǎn Wénhuà Chù : Jìng qǐ ! Wǒ shì Kāmàilóng Duōlā Dì-èr Zhōngxué gāo èr (èr) bān de xuésheng Lǐ Huá. Guānyú guì chù fāfàng 2026 nián Zhōngguó Zhèngfǔ Jiǎngxuéjīn yī shì, wǒ fēicháng guānzhù. Wǒ pǐnxué jiānyōu, rè'ài Zhōnghuá wénhuà, qiě tōngguò le HSK sān jí kǎoshì. Fán qǐng guìfāng yǔyǐ gàozhī shēnqǐng liúchéng yǔ cáiliào qīngdān. Shùn sòng gōng qí ! Lǐ Huá qǐ. 2025 nián shíyī yuè èrshí rì.
\end{exemplebox}

\coursec{Extension Bac \& Évaluation formative (Complément examen)}

\begin{savaistu}[Formules de politesse finales : le « menu » selon le destinataire]
Le chinois codifie la formule finale (\zh{结语}) selon la relation et le secteur. Choisis la bonne « paire » (souvent 2 lignes : \zh{顺颂} / \zh{XXX祺} ou \zh{此致} / \zh{敬礼}) :
\begin{center}
\begin{tabularx}{\linewidth}{|X|X|X|}
\hline
\rowcolor{popblueL} Contexte / Destinataire & Formule (2 lignes) & Pinyin / Sens \\
\hline
Commercial / Entreprise (\zh{公司}, \zh{行}) & \zh{顺颂} \newline \zh{商祺} & \emph{shùn sòng shāng qí} : Souhaiter la prospérité des affaires \\
\hline
Éducation / École / Professeur (\zh{校}, \zh{老师}) & \zh{顺颂} \newline \zh{教祺} & \emph{shùn sòng jiào qí} : Souhaiter la prospérité de l'éducation \\
\hline
Administration / Supérieur hiérarchique / Général & \zh{此致} \newline \zh{敬礼} & \emph{cǐ zhì jìng lǐ} : Par la présente, salutations respectueuses \\
\hline
Aîné / Parent / Ami proche (informel) & \zh{祝} \newline \zh{安好 / 身体健康} & \emph{zhù ān hǎo / shēntǐ jiànkāng} : Souhaiter le bien / la santé \\
\hline
\end{tabularx}
\end{center}
\textbf{Règle d'or :} \zh{顺颂} s'aligne à gauche, \zh{商祺/教祺/公祺} aussi, sur la ligne d'après. \zh{此致} s'aligne à gauche (souvent avec retrait de 2 caractères), \zh{敬礼} sur la ligne d'après, aligné sur \zh{此致} ou retrait de 2 caractères.
\end{savaistu}

\begin{exercice}[Sujet type Bac - Expression écrite (Durée : 30 min)]
\textbf{Sujet :} Tu t'appelles \zh{张明} (\emph{Zhāng Míng}), élève de \zh{高二(3)班} au \zh{雅温得第三中学}. Tu as vu une annonce sur le site de l'\zh{喀麦隆中国文化中心} (Centre Culturel Chinois du Cameroun) recrutant des bénévoles (\zh{志愿者}) pour la « Fête du Printemps 2026 » (\zh{2026年春节文化节}). Écris une lettre formelle de candidature (environ 120-150 caractères chinois).

\textbf{Contraintes obligatoires :}
\begin{enumerate}
\item Structure complète (6 blocs) et mise en page correcte.
\item Utilisation correcte de \zh{贵} (au moins 2 fois) et \zh{您} (si personne physique visée).
\item Utilisation des structures : \zh{关于...一事}, \zh{我是...}, \zh{烦请...予以...}.
\item Vocabulaire ciblé : \zh{志愿者}, \zh{春节}, \zh{文化节}, \zh{品学兼优} (excellence morale et scolaire), \zh{服务} (servir).
\item Formule de politesse finale adaptée au contexte culturel (\zh{顺颂教祺} ou \zh{顺颂公祺}).
\end{enumerate}
\end{exercice}

\begin{corrige}[Sujet type Bac - Proposition de corrigé]
\textbf{Chinois :}\\
喀麦隆中国文化中心：\\
\u3000\u3000敬启！\\
\u3000\u3000我是雅温得第三中学高二(3)班的学生张明。关于贵中心招募2026年春节文化节志愿者一事，我踊跃报名。我品学兼优，性格开朗，热爱中喀文化交流，且有一定中文沟通能力。烦请贵方予以考虑。\\
\u3000\u3000顺颂\\
\u3000\u3000公祺！\\
\u3000\u3000\u3000\u3000\u3000\u3000张明 启\\
\u3000\u3000\u3000\u3000\u3000\u30002025年11月1日\\[4pt]
\textbf{Grille d'évaluation (20 pts) :} Structure/Mise en page (4 pts) | Vocabulaire honorifique \zh{贵}/\zh{您} (3 pts) | Structures grammaticales ciblées (5 pts) | Cohérence lexicale/ton (4 pts) | Caractères/Pinyin corrects (4 pts).
\end{corrige}

\begin{aretenir}[L'essentiel de la leçon ZHA39 — Fiche de révision]
\begin{itemize}
\item \textbf{Format :} 6 blocs immuables (En-tête, Appel, Corps indenté \u3000\u3000, Formule finale 2 lignes, Signature \zh{启} à droite, Date \zh{年/月/日} à droite).
\item \textbf{Déférence (敬语) :} \zh{贵} + bien de l'autre (\zh{贵公司}, \zh{贵校}, \zh{贵方}) ; \zh{您} pour personne ; \zh{尊敬的} + titre.
\item \textbf{Appels :} \zh{敬启} (institution) ; \zh{尊敬的[X]，您好！} (personne).
\item \textbf{Corps :} \zh{关于...一事} (objet) ; \zh{我是...} (identité) ; \zh{且} (argument additionnel) ; \zh{烦请...予以+V} (demande formelle).
\item \textbf{Formules finales :} \zh{顺颂商祺} (commerce) ; \zh{顺颂教祺/公祺} (éduca./officiel) ; \zh{此致敬礼} (admin. général).
\item \textbf{Écriture :} Ordre traits \zh{敬} (haut→bas, \zh{攵} dernier), \zh{启} (ext→int, \zh{户} cadre premier).
\item \textbf{Interdits :} \zh{亲爱的} (amoureux/amis), \zh{你公司} (familier), \zh{号} (date formelle), signature sans \zh{启} (ou \zh{上} si supérieur→inférieur).
\end{itemize}
\end{aretenir}

\coursec{Glossaire thématique : Vocabulaire de la correspondance administrative}

Dans la section précédente, nous avons lu la lettre de demande de stage de Wang Lin adressée à une entreprise de Douala. Vous avez remarqué que certaines expressions ne ressemblent pas au chinois de la conversation quotidienne : on ne dit pas \zh{你的公司} mais \zh{贵公司}, on ne dit pas \zh{看} mais \zh{拜读}. C'est tout l'objet de cette section : construire, pas à pas, le lexique spécialisé de la \cle{correspondance administrative et formelle}, en distinguant les registres de langue, en maîtrisant la prononciation et en classant les mots selon leur fonction dans la lettre.

\courssub{Le tableau de vocabulaire : vingt entrées essentielles}

Observez d'abord le tableau ci-dessous. Lisez chaque mot à voix haute en suivant le pinyin, puis vérifiez la nature grammaticale : en chinois, connaître la nature d'un mot (nom, verbe, préposition) est indispensable pour le placer correctement dans la phrase.

\begin{center}
\begin{tabularx}{\linewidth}{|X|X|X|X|}
\hline
\rowcolor{popblueL} Caractères & Pinyin & Nature & Traduction \\
\hline
贵 & guì & Adj. honorifique & votre (poli) \\
\hline
函 & hán & N & lettre formelle, missive \\
\hline
信 & xìn & N & lettre (courante) \\
\hline
称呼 & chēnghu & N & formule d'appel, titre \\
\hline
落款 & luòkuǎn & N & signature (date et nom en bas) \\
\hline
正文 & zhèngwén & N & corps du texte \\
\hline
附件 & fùjiàn & N & pièce jointe \\
\hline
申请 & shēnqǐng & V / N & postuler, demander ; candidature \\
\hline
寄 & jì & V & envoyer (par la poste) \\
\hline
发 & fā & V & envoyer (général, email) \\
\hline
拜读 & bàidú & V honorifique & lire respectueusement \\
\hline
批示 & pīshì & V / N honorifique & donner des instructions, valider \\
\hline
回复 & huífù & V / N & répondre ; réponse \\
\hline
等候 & děnghòu & V & attendre (poli) \\
\hline
烦请 & fánqǐng & V de politesse & je vous prie de \\
\hline
关于 & guānyú & Prép & au sujet de, concernant \\
\hline
此外 & cǐwài & Connecteur & de plus, en outre \\
\hline
专此 & zhuāncǐ & Connecteur & par la présente (spécialement pour cela) \\
\hline
盼复 & pànfù & Formule & dans l'attente de votre réponse \\
\hline
顺颂商祺 & shùn sòng shāng qí & Formule figée & salutations commerciales distinguées \\
\hline
\end{tabularx}
\end{center}

\begin{definition}[贵 : l'adjectif honorifique]
\zh{贵} (guì) est un adjectif honorifique signifiant « votre / votre honorable ». Il se place \textbf{toujours devant un nom} et jamais seul : \zh{贵公司} (guì gōngsī, « votre entreprise »), \zh{贵姓} (guì xìng, « votre nom », pour demander poliment le nom de famille), \zh{贵国} (guì guó, « votre pays »), \zh{贵校} (guì xiào, « votre école »). Il remplace \zh{你的} dans tout contexte formel.
\end{definition}

\begin{propriete}[顺颂商祺 : la formule de fermeture des affaires]
\zh{顺颂商祺} (shùn sòng shāng qí) est la formule de politesse de clôture standard des lettres d'affaires. Décomposition : \zh{顺} « avec », \zh{颂} « souhaiter », \zh{商} « commerce », \zh{祺} « bonheur ». Elle est \textbf{invariable} et se place seule, sur une ligne, avant la signature. Variantes selon le destinataire : \zh{顺颂公祺} pour un supérieur administratif, \zh{顺颂安祺} pour une personne âgée. Équivalent français : « Veuillez agréer l'expression de mes salutations distinguées ».
\end{propriete}

\courssub{Paires honorifiques et neutres : choisir le bon registre}

Le français oppose « tu » et « vous » ; le chinois administratif oppose des \cle{paires de mots} : un terme neutre pour la vie quotidienne, un terme honorifique pour la correspondance officielle. Voici le tableau comparatif à mémoriser.

\begin{center}
\begin{tabularx}{\linewidth}{|X|X|X|}
\hline
\rowcolor{popgreenL} Registre neutre (quotidien) & Registre honorifique (formel) & Emploi \\
\hline
你 / 您 nǐ / nín (tu / vous) & 贵 guì + nom & devant un nom : 贵公司 \\
\hline
发 fā (envoyer, email) & 寄 jì (envoyer par courrier officiel) & 寄材料 envoyer un dossier \\
\hline
信 xìn (lettre courante) & 函 hán (missive officielle) & 公函 lettre officielle \\
\hline
看 kàn (regarder, lire) & 拜读 bàidú (lire respectueusement) & 拜读您的大作 \\
\hline
回复 huífù (répondre) & 批示 pīshì (valider, instruire) & 烦请批示 \\
\hline
\end{tabularx}
\end{center}

\begin{exemplebox}[Exemples traduits]
1. \zh{关于贵公司招聘实习生一事，我愿申请。}\\
\emph{Guānyú guì gōngsī zhāopìn shíxíshēng yī shì, wǒ yuàn shēnqǐng.}\\
« Au sujet du recrutement de stagiaires par votre entreprise, je souhaite postuler. »\\[4pt]
2. \zh{烦请批示。}\\
\emph{Fánqǐng pīshì.}\\
« Je vous prie de bien vouloir donner vos instructions. »\\[4pt]
3. \zh{专此函达，顺颂商祺。}\\
\emph{Zhuāncǐ hándá, shùn sòng shāng qí.}\\
« Je vous informe par la présente, en vous souhaitant prospérité. »
\end{exemplebox}

\begin{attention}[Piège francophone : 烦请 ne veut pas dire « ennuyer »]
\zh{烦请} (fánqǐng) contient \zh{烦} « ennuyer, déranger », mais la formule entière est \textbf{figée et très polie} : elle signifie « je vous saurais gré de… », « je vous prie de… ». Ne traduisez jamais \zh{烦请贵方尽快回复} par « je vous ennuie pour répondre vite » mais par « je vous prie de bien vouloir répondre rapidement ». De même, \zh{贵} ne se dit jamais seul : on ne peut pas dire \zh{贵好} pour « bonjour à vous ».
\end{attention}

\courssub{Travail phonétique : discrimination des tons}

Dans le vocabulaire administratif, une erreur de ton peut transformer un mot poli en mot étrange. Entraînez-vous sur ces paires minimales, à lire à voix haute :

\begin{itemize}
\item \zh{贵} guì (4\up{e} ton, « honorable ») contre \zh{诡} guǐ (3\up{e} ton, « rusé ») : dire \zh{guǐ gōngsī} serait une catastrophe diplomatique !
\item \zh{批} pī (1\ier{} ton, « approuver ») contre \zh{皮} pí (2\up{e} ton, « peau »).
\item \zh{函} hán (2\up{e} ton, « missive ») contre \zh{喊} hǎn (3\up{e} ton, « crier »).
\item \zh{寄} jì (4\up{e} ton, « envoyer ») contre \zh{急} jí (2\up{e} ton, « pressé »).
\end{itemize}

\begin{methode}[Mémoriser les tons du glossaire]
\begin{enumerate}
\item Lire le mot isolément en exagérant le ton (main qui monte, descend ou tombe).
\item Lire le mot dans sa collocation : \zh{贵公司}, \zh{烦请批示}, \zh{专此函达}.
\item Fermer le tableau, écouter le professeur et écrire le pinyin avec les tons sur l'ardoise (dictée muette).
\item Se corriger en binôme : un élève lit, l'autre pointe le caractère correspondant sur la flashcard.
\end{enumerate}
\end{methode}

\courssub{Classification par fonction dans la lettre}

Un vocabulaire bien organisé se mémorise deux fois plus vite. La carte mentale ci-dessous classe nos vingt mots en quatre familles fonctionnelles : les honorifiques, les mots de structure de la lettre, les verbes d'action et les connecteurs.

\begin{popfigure}
\begin{tikzpicture}[mindmap, concept color=popblue!50, grow cyclic, level 1/.style={level distance=3.2cm, sibling angle=90}, level 2/.style={level distance=2.6cm, sibling angle=45}]
\node[concept] {Lettre formelle 正式信函}
[clockwise from=0]
child[concept color=popgreen!50] { node[concept] {Honorifiques 敬语}
child { node {贵公司} }
child { node {拜读} }
child { node {批示} }
child { node {顺颂商祺} } }
child[concept color=poporange!50] { node[concept] {Structure 结构}
child { node {称呼} }
child { node {正文} }
child { node {落款} }
child { node {附件} } }
child[concept color=poppink!50] { node[concept] {Actions 动词}
child { node {申请} }
child { node {烦请} }
child { node {寄} }
child { node {盼复} } }
child[concept color=poppurple!50] { node[concept] {Connecteurs 连接词}
child { node {关于} }
child { node {此外} }
child { node {专此} }
child { node {函} } };
\end{tikzpicture}
\legende{Carte mentale du vocabulaire de la correspondance administrative : quatre familles fonctionnelles.}
\end{popfigure}

Retenez la logique : les mots de \cle{structure} (\zh{称呼} l'appel, \zh{正文} le corps, \zh{落款} la signature, \zh{附件} la pièce jointe) nomment les parties de la lettre ; les mots d'\cle{action} (\zh{申请}, \zh{寄}, \zh{等候}) expriment ce que fait l'expéditeur ; les \cle{connecteurs} (\zh{关于}, \zh{此外}, \zh{专此}) articulent le propos ; les \cle{honorifiques} marquent le respect dû au destinataire.

\begin{savaistu}[函, la lettre officielle]
Le caractère \zh{函} (hán) signifie « lettre formelle, missive ». Il implique un document officiel, souvent accompagné d'une pièce jointe (\zh{附件} fùjiàn). Un courriel formel se dit \zh{正式邮件} (zhèngshì yóujiàn) ou \zh{电子函} (diànzǐ hán). Au Cameroun comme en Chine, l'administration distingue soigneusement la lettre ordinaire de la correspondance officielle : en chinois, on parle de \zh{公函} (gōnghán, « lettre officielle ») pour les échanges entre administrations.
\end{savaistu}

\courssub{Texte d'application : un extrait de correspondance}

\begin{textebox}[Extrait d'une lettre de demande de stage]
尊敬的张经理：您好！\\
\emph{Zūnjìng de Zhāng jīnglǐ : Nín hǎo !}\\
« Monsieur le Directeur Zhang, bonjour ! »\\[4pt]
关于贵公司招聘实习生一事，我在学校的公告栏上看到了招聘函，非常感兴趣，特此写信申请。\\
\emph{Guānyú guì gōngsī zhāopìn shíxíshēng yī shì, wǒ zài xuéxiào de gōnggàolán shang kàndào le zhāopìn hán, fēicháng gǎn xìngqù, tècǐ xiěxìn shēnqǐng.}\\
« Concernant le recrutement de stagiaires par votre entreprise, j'ai vu l'avis de recrutement sur le panneau d'affichage de mon école ; très intéressé, je vous écris spécialement pour postuler. »\\[4pt]
此外，我的成绩单和推荐信在附件里，烦请查收。\\
\emph{Cǐwài, wǒ de chéngjìdān hé tuījiànxìn zài fùjiàn lǐ, fánqǐng cháshōu.}\\
« De plus, mon relevé de notes et ma lettre de recommandation se trouvent en pièce jointe ; je vous prie de bien vouloir vérifier leur réception. »\\[4pt]
等候您的回复。专此函达，顺颂商祺！\\
\emph{Děnghòu nín de huífù. Zhuāncǐ hándá, shùn sòng shāng qí !}\\
« Dans l'attente de votre réponse. Je vous informe par la présente et vous adresse mes salutations distinguées. »
\end{textebox}

Questions de compréhension rapide : 1) Où l'élève a-t-il vu l'annonce ? 2) Quels documents sont joints à la lettre ? 3) Quelle formule de fermeture est utilisée et pourquoi ?

\courssub{Exercices de reconnaissance rapide}

\begin{exoresolu}[Flashcards : choisir le bon registre]
Énoncé. Pour chaque situation, choisissez le mot correct entre les deux propositions :\\
1. Vous écrivez à une entreprise : « votre entreprise » = \zh{你的公司} ou \zh{贵公司} ?\\
2. Vous demandez poliment une validation : \zh{请你看} ou \zh{烦请批示} ?\\
3. Vous parlez d'une lettre officielle entre ministères : \zh{信} ou \zh{公函} ?\\
4. Vous envoyez un dossier par la poste : \zh{发材料} ou \zh{寄材料} ?
\tcblower
Solution détaillée.\\
1. \zh{贵公司} : dans une lettre formelle, l'adjectif honorifique \zh{贵} remplace obligatoirement \zh{你的}.\\
2. \zh{烦请批示} : \zh{批示} est le verbe honorifique réservé au supérieur qui valide ; \zh{看} est trop familier.\\
3. \zh{公函} : \zh{函} marque le caractère officiel de la missive, contrairement à \zh{信}, lettre courante.\\
4. \zh{寄材料} : \zh{寄} convient à l'envoi postal d'un dossier ; \zh{发} s'emploie plutôt pour les emails (\zh{发邮件}).
\end{exoresolu}

\begin{exercice}[Dictée muette sur ardoise]
Le professeur prononce à voix haute, sans montrer les caractères : \zh{guì gōngsī}, \zh{fùjiàn}, \zh{pīshì}, \zh{zhuāncǐ}, \zh{pànfù}, \zh{luòkuǎn}. Écrivez sur votre ardoise : a) le pinyin avec les marques de tons ; b) les caractères ; c) la traduction française. Puis, en binôme, classez ces six mots dans les quatre familles de la carte mentale (honorifiques, structure, actions, connecteurs).
\end{exercice}

\begin{corrige}[Dictée muette]
\zh{贵公司} guì gōngsī : votre entreprise (honorifique) ; \zh{附件} fùjiàn : pièce jointe (structure) ; \zh{批示} pīshì : valider, instructions (honorifique/action) ; \zh{专此} zhuāncǐ : par la présente (connecteur) ; \zh{盼复} pànfù : dans l'attente de votre réponse (formule/action) ; \zh{落款} luòkuǎn : signature en bas de page (structure). Attention aux tons : \zh{guì} 4\up{e} ton, \zh{fù} 4\up{e} ton, \zh{pī} 1\ier{} ton.
\end{corrige}

\begin{aretenir}[L'essentiel du glossaire]
\begin{itemize}
\item \zh{贵} + nom = « votre » honorifique ; jamais seul.
\item Paires de registre : \zh{发}/\zh{寄}, \zh{信}/\zh{函}, \zh{看}/\zh{拜读}, \zh{回复}/\zh{批示}.
\item \zh{烦请} = « je vous prie de », formule figée, sans idée d'ennui.
\item \zh{顺颂商祺} clôt toute lettre d'affaires ; \zh{专此} introduit la formule finale.
\item Quatre familles : honorifiques, structure, actions, connecteurs.
\end{itemize}
\end{aretenir}

\coursec{Atelier Caractères : Radicaux, Ordre des traits \& Mémorisation}

Dans la section précédente, nous avons constitué le glossaire de la correspondance administrative : \zh{贵公司} (guì gōngsī, « votre honorable société »), \zh{敬启} (jìng qǐ, « respectueusement ouvrir »), \zh{批示} (pī shì, « instructions écrites »), \zh{函件} (hán jiàn, « courrier officiel »). Mais un mot appris uniquement en pinyin reste fragile : il s'oublie en quelques jours. Pour que ce vocabulaire devienne \emph{actif} — c'est-à-dire que vous sachiez l'écrire à la main dans votre lettre de demande de stage — il faut entrer \emph{dans} les caractères : comprendre leur architecture (radical + phonétique), maîtriser l'ordre des traits, et s'appuyer sur des images mnémotechniques. C'est l'objet de cet atelier.

\courssub{Pourquoi analyser les caractères ? Observation d'abord}

Observez ces deux séries de caractères tirés de notre lettre formelle :

\begin{exemplebox}[Observation guidée]
Série A : \zh{礼} (lǐ, rite) — \zh{祝} (zhù, souhaiter) — \zh{福} (fú, bonheur) — \zh{祺} (qí, auspice)\\
Série B : \zh{批} (pī, annoter) — \zh{打} (dǎ, frapper) — \zh{报} (bào, annoncer) — \zh{推} (tuī, pousser)
\end{exemplebox}

Que remarquez-vous ? Tous les caractères de la série A contiennent à gauche le même élément \zh{礻} ; tous ceux de la série B contiennent \zh{扌}. Ce n'est pas un hasard : \zh{礻} est la clé de l'\emph{autel} (domaine du sacré, des rites, des vœux), et \zh{扌} est la clé de la \emph{main} (domaine des actions manuelles). Le radical donne le \cle{champ sémantique} ; l'autre partie donne souvent un indice de \cle{prononciation} (le composant phonétique).

\begin{definition}[Radical et composant phonétique]
Un caractère chinois composé se décompose généralement en deux parties : le \cle{radical} (ou clé, \zh{部首} bùshǒu), qui indique la famille de sens, et le \cle{composant phonétique} (\zh{声旁} shēngpáng), qui suggère la prononciation. Exemple : \zh{祺} qí = \zh{礻} (autel, sens : rite, vœu) + \zh{其} qí (phonétique). Apprendre un caractère, c'est donc apprendre une \emph{équation} : sens + son.
\end{definition}

\begin{definition}[Radical 礻 (shì), la clé de l'autel]
\zh{礻} est la forme comprimée à gauche du caractère \zh{示} (shì, « montrer, autel »). Il indique le domaine du \emph{sacré}, des \emph{rites} et des \emph{vœux}. On le retrouve dans : \zh{礼} (lǐ, rite, politesse), \zh{祝} (zhù, souhaiter), \zh{福} (fú, bonheur), \zh{祥} (xiáng, auspice), \zh{祺} (qí, auspice), \zh{祖} (zǔ, ancêtre). Piège : ne pas confondre avec \zh{衤} (yī, clé du vêtement, avec un trait de plus) : \zh{补} (bǔ, réparer), \zh{被} (bèi, couverture).
\end{definition}

\courssub{Les huit caractères stratégiques de la lettre formelle}

Voici le tableau d'analyse complet des huit caractères cibles de cette leçon. Pour chacun : le radical, le composant phonétique, et une piste de mémorisation.

\begin{center}
\begin{tabularx}{\linewidth}{|X|X|X|X|}
\hline
\rowcolor{popblueL} Caractère & Radical (clé) & Phonétique & Analyse et mémo \\
\hline
\zh{贵} guì (précieux, honorifique) & \zh{贝} bèi (coquillage, monnaie) & partie haute & Ce qui contient des coquillages-monnaie est \emph{précieux} : \zh{贵公司} « votre honorable société » \\
\hline
\zh{敬} jìng (respecter) & \zh{攵} pū (frapper doucement, action) & \zh{苟} gǒu & Agir avec retenue = respect : \zh{敬启} « respectueusement ouvrir » \\
\hline
\zh{启} qǐ (ouvrir, commencer) & \zh{口} kǒu (bouche) & \zh{户} hù (porte) & Ouvrir la porte avec la parole : \zh{启事} « avis public » \\
\hline
\zh{烦} fán (ennuyer, prier de) & \zh{火} huǒ (feu) & \zh{页} yè (tête, page) & Le feu sur la tête : cela chauffe dans la tête, d'où « tracas » ; en lettre : \zh{烦请您…} « je vous prie de… » \\
\hline
\zh{批} pī (annoter, approuver) & \zh{扌} shǒu (main) & \zh{比} bǐ & Action de la main qui compare et juge : \zh{批示} « instructions écrites » \\
\hline
\zh{示} shì (montrer, indiquer) & \zh{示} shì (autel) — radical entier & — & L'autel où le ciel \emph{montre} les signes ; donne la clé \zh{礻} \\
\hline
\zh{函} hán (lettre, courrier) & \zh{凵} kǎn (récipient ouvert) & intérieur & Une enveloppe-boîte qui contient un message : \zh{函件} « courrier officiel ». \emph{Note : phonétique 咸 xián masquée en simplifié.} \\
\hline
\zh{祺} qí (auspice, bonheur) & \zh{礻} shì (autel, rite) & \zh{其} qí & Le bonheur obtenu par le rite : formule de politesse finale \zh{顺颂时祺} \\
\hline
\end{tabularx}
\end{center}

\begin{exemplebox}[Deux collocations à retenir]
\zh{批示} (pī shì) : verbe / nom. Instructions écrites d'un supérieur hiérarchique. \zh{请领导批示。} \emph{Qǐng lǐng dǎo pī shì.} « Je prie la direction de bien vouloir annoter et approuver. »\\
\zh{函件} (hán jiàn) : nom. Courrier officiel, lettre formelle. \zh{贵公司的函件我们已经收到。} \emph{Guì gōng sī de hán jiàn wǒ men yǐ jīng shōu dào le.} « Nous avons bien reçu le courrier de votre honorable société. »
\end{exemplebox}

\courssub{L'ordre des traits (笔顺 bǐ shùn) : les règles d'or}

Écrire un caractère n'est pas un dessin libre : c'est une \emph{chorégraphie} codifiée. Respecter l'ordre des traits rend l'écriture plus rapide, plus régulière, et permet de consulter un dictionnaire par nombre de traits.

\begin{propriete}[Règle d'or de l'ordre des traits]
\begin{enumerate}
\item \cle{Horizontal} (\zh{横} héng) \emph{avant} \cle{vertical} (\zh{竖} shù) : on écrit \zh{十} en commençant par la barre horizontale.
\item \cle{Gauche} (\zh{左}) \emph{avant} \cle{droite} (\zh{右}) : dans \zh{批}, on écrit d'abord \zh{扌}, puis \zh{比}.
\item \cle{Haut} (\zh{上}) \emph{avant} \cle{bas} (\zh{下}) : dans \zh{贵}, on écrit la partie supérieure avant \zh{贝}.
\item \cle{Extérieur} (\zh{外}) \emph{avant} \cle{intérieur} (\zh{内}) \emph{avant} fermeture (\zh{封}) : dans \zh{函}, on trace d'abord le contenu, puis le récipient \zh{凵} qui ferme.
\item Le \cle{point} (\zh{点} diǎn) se trace en dernier s'il est en haut à droite (comme dans \zh{我}), sinon en premier.
\end{enumerate}
\end{propriete}

Application au caractère \zh{贵} (guì), \textbf{8 traits} (forme simplifiée), de haut en bas :
1. Trait vertical central haut (\zh{丨}) ; 2. Trait horizontal long (\zh{一}) ; 3. Trait vertical gauche (\zh{丨}) ; 4. Pli horizontal-vertical (\zh{㇇}, haut de \zh{贝}) ; 5. Horizontal moyen (\zh{一}) ; 6. Horizontal bas (\zh{一}) ; 7. Horizontal inférieur (\zh{一}) ; 8. Trait vertical de fermeture (\zh{丨}).

\begin{popfigure}
\begin{tikzpicture}[scale=0.85]
% Grille de fond 田字格
\draw[step=1cm, gray!30] (0,0) grid (4,4);
\draw[thick, gray] (2,0) -- (2,4); % médiane verticale
\draw[thick, gray] (0,2) -- (4,2); % médiane horizontale

% Traits du caractère 贵 (guì) - 8 traits
% 1. Vertical central top (bleu)
\draw[very thick, popblue, ->] (2,3.5) -- (2,2.5);
\node[popblue, right] at (2.2,3) {\small 1. 丨};

% 2. Horizontal top (orange)
\draw[very thick, poporange, ->] (0.5,3.5) -- (3.5,3.5);
\node[poporange, above] at (2,3.6) {\small 2. 一};

% 3. Vertical left top (vert)
\draw[very thick, popgreen, ->] (0.5,3.5) -- (0.5,2.5);
\node[popgreen, left] at (0.4,3) {\small 3. 丨};

% 4. Horizontal fold top-right (pourpre) - une seule flèche courbe ou deux segments
\draw[very thick, poppurple, ->] (0.5,2.5) -- (3.5,2.5);
\draw[very thick, poppurple, ->] (3.5,2.5) -- (3.5,1.5);
\node[poppurple, right] at (3.6,2) {\small 4. ㇇};

% 5. Horizontal middle 1 (rose)
\draw[very thick, poppink, ->] (1,2.0) -- (3,2.0);
\node[poppink, left] at (0.9,2.0) {\small 5. 一};

% 6. Horizontal middle 2 (bleu)
\draw[very thick, popblue, ->] (1,1.5) -- (3,1.5);
\node[popblue, left] at (0.9,1.5) {\small 6. 一};

% 7. Horizontal lower (orange)
\draw[very thick, poporange, ->] (1,1.0) -- (3,1.0);
\node[poporange, left] at (0.9,1.0) {\small 7. 一};

% 8. Vertical closing right (vert)
\draw[very thick, popgreen, ->] (3,2.5) -- (3,0.5);
\node[popgreen, right] at (3.1,1.5) {\small 8. 丨};

\node[align=center, font=\small] at (2,-0.4) {Ordre des traits : \zh{贵} (guì) — Radical \zh{贝} (bèi)};
\end{tikzpicture}
\legende{Tracé guidé sur quadrillage 田字格 : horizontal (2) avant vertical (1,3) ; haut vers bas ; gauche vers droite ; fermeture finale (8). Couleurs = ordre séquentiel.}
\end{popfigure}

\courssub{Étymologie vivante : le cas de 祺 (qí)}

\begin{savaistu}[Le « bonheur rituel »]
\zh{祺} (qí) se décompose en \zh{礻} (l'autel, le rite) + \zh{其} (qí, phonétique parfait). Son sens premier est le « bonheur obtenu par le rite », l'\emph{auspice}. Aujourd'hui, il survit presque exclusivement dans les formules de politesse finales des lettres formelles : \zh{顺颂时祺} (shùn sòng shí qí, « je vous souhaite, au fil des jours, toute félicité »), équivalent chinois de « Veuillez agréer, Madame, Monsieur, l'expression de mes salutations distinguées ». Apprendre ce caractère, c'est donc apprendre directement une \emph{formule de lettre} entière.
\end{savaistu}

\begin{methode}[Mnémotechnique pour 烦 (fán, ennuyer / prier de)]
\begin{enumerate}
\item Décomposer : \zh{烦} = \zh{火} (le feu, radical à gauche) + \zh{页} (la tête / la page, à droite).
\item Construire l'image : « le feu sur la tête » — quand quelque chose vous tracasse, cela « chauffe » dans la tête.
\item Relier au sens courant : ce qui fait chauffer la tête est \emph{ennuyeux, fatigant}.
\item Relier à l'emploi épistolaire : demander un service à quelqu'un, c'est lui « faire chauffer la tête » ; d'où la tournure hyper-polie \zh{烦请您…} (fán qǐng nín…, « je vous prie de bien vouloir faire l'effort de… »), très fréquente dans une demande de stage.
\end{enumerate}
\end{methode}

\courssub{Pièges visuels : ne pas confondre}

\begin{attention}[Confusion visuelle : 贵 (guì) et 遗 (yí)]
\zh{贵} (guì, précieux) et \zh{遗} (yí, perdre, laisser) se ressemblent car ils partagent la même partie supérieure. La différence est en bas : \zh{贵} se termine par \zh{贝} (le coquillage-monnaie, structure fermée, 8 traits au total), tandis que \zh{遗} est enveloppé par la clé de la marche \zh{辶} (chuò). Erreur fréquente des francophones : écrire \zh{遗公司} au lieu de \zh{贵公司} dans l'en-tête de la lettre — ce qui transformerait « votre honorable société » en « société perdue » ! Vérifiez toujours le bas du caractère.
\end{attention}

\begin{attention}[礻 ou 衤 ? Un trait qui change tout]
La clé de l'autel \zh{礻} a \emph{quatre} traits (丶 一 丨 丶) ; la clé du vêtement \zh{衤} en a \emph{cinq} (丶 ㇇ 丨 丶 一 — un horizontal final en plus). Écrire \zh{祺} avec la clé du vêtement produit un caractère inexistant. Astuce : les vœux et les rites (\zh{礼}, \zh{祝}, \zh{福}, \zh{祺}) n'ont rien à voir avec les habits — un seul point en haut à droite du radical, pas deux.
\end{attention}

\courssub{Exercices d'écriture guidée sur quadrillage (田字格)}

\begin{experience}[Atelier d'écriture en trois temps]
\begin{enumerate}
\item \cle{Tracé en l'air} : le professeur trace \zh{贵} au tableau en annonçant chaque trait à voix haute ; les élèves « écrivent » en l'air avec l'index, en comptant les traits en chinois : \zh{一、二、三、四、五、六、七、八} (yī, èr, sān, sì, wǔ, liù, qī, bā).
\item \cle{Tracé sur tablette ou cahier} : chaque élève recopie les huit caractères cibles dans un quadrillage 田字格, trois fois chacun, en respectant l'ordre des traits. Le caractère doit occuper le centre du carré, équilibré entre les quatre cases.
\item \cle{Dictée éclair} : le professeur dicte \zh{批示}, \zh{函件}, \zh{敬启}, \zh{贵公司} ; les élèves écrivent en caractères, puis vérifient en binôme.
\end{enumerate}
\end{experience}

\begin{exoresolu}[Jeu du caractère intrus]
Énoncé : dans la série suivante, un caractère ne partage pas le même radical que les autres. Lequel ? Justifiez.\\
\zh{示} — \zh{礼} — \zh{祺} — \zh{福} — \zh{批}
\tcblower
Solution détaillée : \zh{示} est le radical « autel » sous sa forme indépendante (clé 113) ; \zh{礼}, \zh{祺} et \zh{福} contiennent tous la forme composée \zh{礻} (variante de gauche de la clé 113), marquant le champ des rites et vœux. En revanche, \zh{批} (pī, annoter) contient la clé de la main \zh{扌} (clé 64) : il appartient à la famille des actions manuelles. L'intrus est donc \zh{批}. Ce jeu entraîne l'œil à repérer instantanément le radical, compétence indispensable pour utiliser un dictionnaire chinois classé par clés.
\end{exoresolu}

\begin{exercice}[À vous de jouer]
1. Décomposez les caractères \zh{敬} et \zh{函} en radical + composant, et proposez une image mnémotechnique pour chacun.\\
2. Trouvez l'intrus : \zh{祝} — \zh{祥} — \zh{祖} — \zh{打} — \zh{福}.\\
3. Écrivez \zh{贵} en respectant l'ordre des traits (8 traits), puis utilisez-le dans la formule d'appel d'une lettre formelle.
\end{exercice}

\begin{corrige}[Exercice 3.6]
1. \zh{敬} = \zh{苟} (phonétique gǒu, indice de son) + \zh{攵} (radical action) : « agir avec retenue, c'est respecter ». \zh{函} = \zh{凵} (radical récipient ouvert) + contenu interne : « une boîte-enveloppe qui contient un message » (phonétique historique 咸 xián masquée).\\
2. L'intrus est \zh{打} (dǎ, frapper) : clé de la main \zh{扌}, alors que \zh{祝}, \zh{祥}, \zh{祖}, \zh{福} portent la clé de l'autel \zh{礻}.\\
3. Ordre des traits : 丨 → 一 → 丨 → ㇇ → 一 → 一 → 一 → 丨. Formule attendue : \zh{贵公司：} (Guì gōng sī :) « À l'attention de votre honorable société », placée en haut de la lettre, suivie de deux points.
\end{corrige}

\begin{aretenir}[L'essentiel de l'atelier]
\begin{itemize}
\item Huit caractères clés de la lettre formelle : \zh{贵}, \zh{敬}, \zh{启}, \zh{烦}, \zh{批}, \zh{示}, \zh{函}, \zh{祺}.
\item Chaque caractère = radical (sens) + phonétique (son) : l'analyser, c'est déjà le mémoriser.
\item Ordre des traits : horizontal avant vertical, gauche avant droite, haut avant bas, extérieur avant intérieur avant fermeture.
\item Vigilance visuelle : \zh{贵} (8 traits, clé \zh{贝}) / \zh{遗} (clé \zh{辶}) ; \zh{礻} (4 traits, autel) / \zh{衤} (5 traits, vêtement).
\end{itemize}
\end{aretenir}

Vos caractères sont désormais solidement ancrés : vous savez les lire, les décomposer et les écrire. La prochaine étape consistera à les assembler dans l'architecture rhétorique de la lettre : connecteurs et structures types du corps de la lettre.

\coursec{Architecture rhétorique : Connecteurs \& Structures types du corps de lettre}

Après avoir appris à écrire et mémoriser les caractères clés de la correspondance administrative (\zh{信}, \zh{函}, \zh{复}, \zh{请}…), nous allons maintenant assembler ces briques : comment construit-on, en chinois formel, le \emph{corps} d'une lettre officielle ? Le chinois épistolaire possède une architecture très codifiée, presque mathématique : chaque étape du raisonnement correspond à une structure figée (\zh{句式}, \emph{jùshì}) qu'il suffit d'apprendre comme un moule. C'est une excellente nouvelle pour vous : une fois les patrons maîtrisés, rédiger une lettre formelle devient un exercice de précision, pas d'invention.

\courssub{Observation : trois modèles de paragraphes de corps de lettre}

Lisons d'abord trois extraits de lettres authentiques de style administratif. Observez les expressions en gras : elles reviendront dans toutes les lettres formelles.

\begin{textebox}[Trois corps de lettre : demande, réclamation, envoi de documents]
\textbf{1. Demande (申请)} :\\
关于\textbf{申请暑期实习}一事，本人系雅温得第一中学高二年级学生，特此致函，恳请贵公司给予实习机会。\\
\emph{Guānyú shēnqǐng shǔqī shíxí yī shì, běnrén xì Yǎwēndé Dìyī Zhōngxué gāoèr niánjí xuésheng, tècǐ zhìhán, kěnqǐng guì gōngsī jǐyǔ shíxí jīhuì.}\\
« Au sujet de la demande de stage d'été, je soussigné, élève de Première au Lycée de Yaoundé I, écris la présente lettre pour solliciter de votre honorable société une opportunité de stage. »\\[4pt]
\textbf{2. Réclamation (投诉/索赔)} :\\
鉴于贵方延迟交货，给我方造成严重损失，特此提出索赔，烦请尽快处理为盼。\\
\emph{Jiànyú guìfāng yánchí jiāohuò, gěi wǒfāng zàochéng yánzhòng sǔnshī, tècǐ tíchū suǒpéi, fánqǐng jǐnkuài chǔlǐ wéi pàn.}\\
« Vu que votre partie a retardé la livraison, causant à notre partie de graves pertes, nous présentons par la présente une réclamation et vous prions de bien vouloir traiter ce dossier au plus vite. »\\[4pt]
\textbf{3. Envoi de documents (寄送材料)} :\\
兹随函附上本人成绩单及推荐信各一份，请查收。静候佳音。\\
\emph{Zī suíhán fùshàng běnrén chéngjìdān jí tuījiànxìn gè yī fèn, qǐng cháshōu. Jìnghòu jiāyīn.}\\
« Nous joignons à la présente un relevé de notes et une lettre de recommandation ; veuillez en accuser réception. Dans l'attente de votre favorable réponse. »
\end{textebox}

\textbf{Observons ensemble.} Trois constantes apparaissent :
\begin{itemize}
\item La lettre s'ouvre sur l'\cle{objet} introduit par \zh{关于} (au sujet de) suivi de \zh{一事} (l'affaire de…) ou \zh{问题} (la question de…).
\item Le motif est introduit par \zh{鉴于} ou \zh{由于} (vu que, étant donné que) — jamais par le \zh{因为} de la conversation.
\item L'action est formulée avec \zh{特此} / \zh{专此} (par la présente) et la demande polie avec \zh{烦请…为盼} (je vous prie de…, dans l'espoir de votre réponse).
\end{itemize}

\courssub{Les patrons syntaxiques (句式) du corps de lettre}

\begin{propriete}[Patron n°1 : 关于 + sujet + 一事/问题，…]
\zh{关于} (guānyú, « au sujet de ») ouvre la lettre en annonçant l'objet, suivi de \zh{一事} (« l'affaire de ») ou \zh{问题} (« la question de »). C'est l'équivalent du français « Objet : … » ou « J'ai l'honneur de … au sujet de … ».\\
Structure : \textbf{关于 + [objet] + 一事，+ [phrase principale]}\\
Exemple : \zh{关于贵公司招聘翻译一事，本人特此申请。} \emph{Guānyú guì gōngsī zhāopìn fānyì yī shì, běnrén tècǐ shēnqǐng.} « Au sujet du recrutement d'un traducteur par votre société, je pose par la présente ma candidature. »
\end{propriete}

\begin{propriete}[Patron n°2 : 烦请 + sujet + verbe + 为盼 / 见复]
\zh{烦请} (fánqǐng, « je vous prie de bien vouloir ») introduit la demande polie. La clôture \zh{为盼} (wéi pàn, « dans l'espoir ») ou \zh{见复} (jiàn fù, « en attendant votre réponse ») termine la phrase.\\
Structure : \textbf{烦请 + [destinataire] + [action] + 为盼/见复}\\
Exemple : \zh{烦请贵处核实为盼。} \emph{Fánqǐng guìchù héshí wéi pàn.} « Je vous prie de bien vouloir vérifier, dans l'espoir d'une suite favorable. »\\
Variante : \zh{盼复} (pàn fù) « dans l'attente de votre réponse ».
\end{propriete}

\begin{propriete}[Structure « 鉴于…特此… »]
\zh{鉴于} (jiànyú, « donné que / vu que ») introduit le motif légal ou factuel. \zh{特此} (tècǐ, « par la présente / spécialement pour cela ») introduit l'action formelle : \zh{函达} (communiquer par lettre), \zh{通知} (notifier), \zh{申请} (solliciter), \zh{致歉} (présenter ses excuses).\\
Structure : \textbf{鉴于 + [motif]，特此 + [action formelle]}\\
Exemple : \zh{鉴于贵方延迟交货，特此提出索赔。} \emph{Jiànyú guìfāng yánchí jiāohuò, tècǐ tíchū suǒpéi.} « Vu le retard de livraison de votre part, nous présentons par la présente une demande d'indemnisation. »
\end{propriete}

\begin{definition}[专此 (zhuān cǐ)]
« Spécialement pour cette affaire / par la présente ». Sert à clore le corps de la lettre juste avant la formule de politesse : \zh{专此函达。} « Par la présente, nous vous communiquons ceci. » Synonyme formel : \zh{特此} (tè cǐ). \textbf{Ne pas confondre} avec \zh{专门} (zhuānmén, « spécialement / dédié à »), qui est un adverbe courant : \zh{我专门为你做了这道菜。} \emph{Wǒ zhuānmén wèi nǐ zuòle zhè dào cài.} « J'ai préparé ce plat spécialement pour toi. »
\end{definition}

\courssub{Comparaison français / chinois : thème, rhème et la voix passive}

\begin{propriete}[Ordre thème–rhème et rareté du passif 被]
En français, la lettre formelle aime le passif : « Votre demande \emph{a été reçue} », « Il vous \emph{est demandé} de… ». En chinois, le passif avec \zh{被} (bèi) est \textbf{rare} dans la correspondance administrative, car il porte souvent une nuance négative (subir). On préfère :
\begin{enumerate}
\item \textbf{L'actif impersonnel} : \zh{已收到贵方来函。} \emph{Yǐ shōudào guìfāng láihán.} « (Nous) avons bien reçu votre lettre » — le sujet est sous-entendu.
\item \textbf{La mise en avant du thème} : le sujet connu (thème) se place en tête, l'information nouvelle (rhème) après : \zh{材料已随函附上。} \emph{Cáiliào yǐ suíhán fùshàng.} « Les documents sont joints à la présente » (litt. « documents — déjà joints à la lettre »).
\item \textbf{Le passif « noble » avec 承/获} : \zh{承蒙贵方支持，特此致谢。} \emph{Chéngméng guìfāng zhīchí, tècǐ zhìxiè.} « Ayant bénéficié de votre soutien, nous vous remercions par la présente. »
\end{enumerate}
\end{propriete}

\begin{attention}[Piège : le verbe 致 (zhì)]
\zh{致} signifie « présenter / adresser » (hommages, salutations, respect). Structure : \textbf{致 + [nom abstrait] + 于/给 + [destinataire]}.\\
Exemple : \zh{向贵公司致以崇高的敬意。} \emph{Xiàng guì gōngsī zhìyǐ chónggāo de jìngyì.} « Nous présentons à votre honorable société nos respectueuses salutations. »\\
\textbf{Erreur fréquente des francophones} : traduire « présenter à » par \zh{给} en style très formel. Préférez \zh{于} ou la tournure avec \zh{向…致以…} ; \zh{给} reste correct mais marque un registre plus courant.
\end{attention}

\courssub{Connecteurs logiques : langue écrite (书面语) vs langue parlée (口语)}

Le passage du style oral au style épistolaire repose sur le remplacement systématique des connecteurs. Retenez ce tableau par cœur :

\begin{center}
\begin{tabularx}{\linewidth}{|X|X|X|X|}
\hline
\rowcolor{popblueL} Sens logique & Oral (口语) & Formel (书面语) & Exemple formel \\
\hline
Addition & 而且 érqiě & 此外 / 另外 cǐwài / lìngwài & 此外，本人精通法语。 \\
\hline
Cause & 因为 yīnwèi & 由于 / 鉴于 yóuyú / jiànyú & 由于工作调整，… \\
\hline
Conséquence & 所以 suǒyǐ & 因此 / 故 yīncǐ / gù & 故取消原定行程。 \\
\hline
Opposition & 但是 dànshì & 然而 / 但 rán'ér / dàn & 然而条件尚不成熟。 \\
\hline
Condition & 如果 rúguǒ & 若 / 如 ruò / rú & 若有疑问，请联系。 \\
\hline
But & 为了 wèile & 为 / 以便 wèi / yǐbiàn & 为促进合作，… \\
\hline
\end{tabularx}
\end{center}

\begin{exemplebox}[Transformations oral → épistolaire]
1. Oral : \zh{因为下雨，所以我不去了。} → Écrit : \zh{鉴于天气恶劣，故取消行程。} \emph{Jiànyú tiānqì èliè, gù qǔxiāo xíngchéng.} « Vu les intempéries, le déplacement est annulé. »\\
2. Oral : \zh{请你帮我寄一下。} → Écrit : \zh{烦请惠寄。} \emph{Fánqǐng huìjì.} « Je vous prie de bien vouloir l'envoyer (avec bienveillance). »\\
3. Oral : \zh{等你回复。} → Écrit : \zh{静候佳音。/ 盼复。} \emph{Jìnghòu jiāyīn. / Pànfù.} « Dans l'attente de votre favorable réponse. »
\end{exemplebox}

\begin{attention}[Erreurs fréquentes des francophones]
\begin{itemize}
\item Ne dites jamais \zh{因为…所以…} dans une lettre formelle : c'est la marque immédiate du style parlé. Utilisez \zh{鉴于…故…} ou \zh{由于…因此…}.
\item \zh{故} (gù) seul suffit ; n'écrivez pas \zh{故所以} (pléonasme).
\item Attention à \zh{另外} : il ajoute un point (« de plus »), il ne signifie pas « autrement ».
\end{itemize}
\end{attention}

\courssub{Schéma de l'architecture d'une lettre formelle}

\begin{popfigure}
\begin{tikzpicture}[node distance=1.1cm, auto]
\node[draw, rectangle, fill=popblueL, align=center, rounded corners, minimum width=7cm] (start) {\textbf{Objet} : 关于…一事};
\node[draw, rectangle, fill=popgreenL, below of=start, align=center, rounded corners, minimum width=7cm] (context) {\textbf{Contexte / Motif} : 鉴于 / 由于…};
\node[draw, rectangle, fill=popgoldL, below of=context, align=center, rounded corners, minimum width=7cm] (action) {\textbf{Demande / Action} : 烦请 / 希 / 拟…};
\node[draw, rectangle, fill=poporangeL, below of=action, align=center, rounded corners, minimum width=7cm] (closing) {\textbf{Clôture du corps} : 专此 / 特此…};
\node[draw, rectangle, fill=poppinkL, below of=closing, align=center, rounded corners, minimum width=7cm] (polite) {\textbf{Formule de politesse} : 顺颂…};
\draw[->, thick, popdark] (start) -- (context);
\draw[->, thick, popdark] (context) -- (action);
\draw[->, thick, popdark] (action) -- (closing);
\draw[->, thick, popdark] (closing) -- (polite);
\end{tikzpicture}
\legende{Les cinq étages obligatoires du corps d'une lettre formelle chinoise : chaque étage correspond à un patron syntaxique figé.}
\end{popfigure}

\begin{methode}[Rédiger le corps d'une lettre formelle en 5 étapes]
\begin{enumerate}
\item \textbf{Annoncer l'objet} avec \zh{关于…一事} (jamais de phrase longue d'introduction comme en français).
\item \textbf{Exposer le motif} avec \zh{鉴于} ou \zh{由于} + fait précis (date, référence, événement).
\item \textbf{Formuler la demande} avec \zh{烦请…为盼} ou \zh{特此申请}.
\item \textbf{Clore le corps} avec \zh{专此函达} ou \zh{特此告知}.
\item \textbf{Terminer par la formule de politesse} : \zh{顺颂商祺} (à une entreprise), \zh{此致敬礼} (formule générale).
\end{enumerate}
\end{methode}

\begin{savaistu}[Le petit mot 即 (jí)]
En fin de phrase ou de proposition, \zh{即} signifie « c'est-à-dire / à savoir / soit ». Il est très fréquent dans les courriers administratifs pour préciser une date ou un détail : \zh{拟定下周一即10月20日面试。} \emph{Nǐdìng xià zhōuyī jí shí yuè èrshí rì miànshì.} « Entretien prévu lundi prochain, soit le 20 octobre. » En chinois classique, \zh{即} signifiait déjà « c'est précisément » — un héritage de plus de deux mille ans de correspondance officielle !
\end{savaistu}

\courssub{Exercice résolu : du style oral au style épistolaire}

\begin{exoresolu}[Réécrire en style formel]
Réécrivez ces cinq phrases orales en style épistolaire formel :\\
1. \zh{因为你们发货晚了，所以我们很不高兴。}\\
2. \zh{我想申请你们公司的实习。}\\
3. \zh{请你快点给我回信。}\\
4. \zh{我把成绩单放在信里了，你看看。}\\
5. \zh{而且我还会说英语。}
\tcblower
\textbf{Solutions détaillées :}\\
1. \zh{鉴于贵方延迟发货，特此表示不满。} — \zh{因为…所以…} devient \zh{鉴于…特此…} ; « pas contents » devient la formule neutre \zh{表示不满}.\\
2. \zh{关于申请贵公司实习一事，特此致函。} — ouverture par le patron \zh{关于…一事} + \zh{特此}.\\
3. \zh{烦请尽早见复为盼。} — \zh{快点} (vite !) disparaît : la politesse formelle dit \zh{尽早} (au plus tôt) + \zh{见复/为盼}.\\
4. \zh{兹随函附上成绩单一份，请查收。} — « j'ai mis dans la lettre » devient \zh{随函附上} ; « regarde » devient \zh{请查收} (veuillez vérifier la réception).\\
5. \zh{此外，本人亦通英语。} — \zh{而且} devient \zh{此外} ; « je » devient \zh{本人} ; \zh{亦} remplace \zh{也} en style écrit.
\end{exoresolu}

\begin{exercice}[À vous de jouer]
Transformez en style épistolaire : 1. \zh{因为生病了，所以我不能来面试。} 2. \zh{谢谢你上次帮我。} 3. \zh{但是你们的条件不太好。} (Indices : \zh{鉴于}, \zh{特此致谢}, \zh{然而}…)
\end{exercice}

\begin{corrige}[Exercice : À vous de jouer]
1. \zh{鉴于本人身体不适，故无法参加面试，特此致歉。} « Vu mon indisposition, je ne puis me présenter à l'entretien et vous présente mes excuses. »\\
2. \zh{承蒙贵方此前鼎力相助，特此致谢。} « Reconnaissant de votre précieuse aide passée, nous vous remercions par la présente. »\\
3. \zh{然而贵方所提条件尚欠妥当。} « Toutefois, les conditions proposées par votre partie ne semblent pas encore appropriées. »
\end{corrige}

\begin{aretenir}[L'essentiel de la section]
\begin{itemize}
\item Le corps de lettre suit cinq étages : objet (\zh{关于…一事}) → motif (\zh{鉴于/由于}) → demande (\zh{烦请…为盼}) → clôture (\zh{专此/特此}) → politesse (\zh{顺颂}).
\item Tout connecteur oral a un équivalent écrit : \zh{因为}→\zh{鉴于}, \zh{所以}→\zh{因此/故}, \zh{但是}→\zh{然而}, \zh{而且}→\zh{此外}.
\item Le passif \zh{被} est à éviter ; préférez l'impersonnel ou la mise en tête du thème.
\item \zh{专此} ≠ \zh{专门} ; \zh{致} exige un nom abstrait (\zh{致以敬意}).
\end{itemize}
\end{aretenir}

\coursec{Compréhension de texte \& Collocations : Du passif à l'actif}

Dans la section précédente, nous avons démonté le « squelette » de la lettre formelle : ses connecteurs (\zh{关于}、\zh{鉴于}、\zh{烦请}、\zh{专此}、\zh{顺颂}) et ses structures types. Une belle architecture ne suffit cependant pas : il faut maintenant vérifier que nous \emph{comprenons} réellement la lettre de demande de stage de \zh{王伟} (Wáng Wěi) étudiée en section 1, et que nous savons \emph{réemployer} activement ses blocs lexicaux. C'est exactement le double mouvement de cette section : d'abord la lecture (compréhension progressive, du global à l'inférentiel), ensuite la production (collocations, cloze, traduction). Passer du passif à l'actif, c'est transformer des mots reconnus en mots maîtrisés.

\courssub{Lire en profondeur : la compréhension progressive}

\begin{methode}[La lecture en trois passes]
Face à une lettre formelle chinoise à l'examen, procédez toujours ainsi :
\begin{enumerate}
\item \textbf{Lecture globale} : identifier l'objet de la lettre (\zh{关于} + nom en début de texte) et les acteurs (qui écrit ? à qui ?).
\item \textbf{Lecture détaillée} : repérer les informations explicites (chiffres, niveaux, diplômes, dates) et l'argumentaire (les atouts présentés, souvent introduits par \zh{首先} / \zh{其次} / \zh{此外}).
\item \textbf{Lecture inférentielle} : dégager le ton et le registre (choix des honorifiques, formules de politesse) et répondre à la question implicite : \emph{pourquoi l'auteur écrit-il ainsi ?}
\end{enumerate}
\end{methode}

\begin{popfigure}
\begin{tikzpicture}
\matrix (m) [matrix of nodes,
    nodes={draw=popline, fill=popblueL, align=center, minimum height=1cm, text width=3.5cm, inner sep=4pt, font=\small},
    row 1/.style={nodes={fill=popblue, font=\bfseries\small, text=white}},
    row sep=-\pgflinewidth,
    column sep=-\pgflinewidth
] {
\textbf{Niveau} & \textbf{Type de question} & \textbf{Exemple} \\
Global & Objet principal & De quoi parle la lettre ? \\
Global & Acteurs & Qui écrit à qui ? \\
Détailé & Info explicite & Quel niveau HSK ? \\
Détailé & Argumentaire & Quels deux atouts ? \\
Inférentiel & Registre / Ton & Pourquoi \zh{烦请} ? \\
};
\end{tikzpicture}
\legende{Grille de compréhension progressive : trois niveaux de lecture, cinq questions types.}
\end{popfigure}

\courssub{Les cinq questions de compréhension (modèle Bac)}

\begin{exemplebox}[Questions types Bac — avec réponses attendues]
\textbf{1. (Global — Acteurs)} \zh{信的发件人是谁？收件人是谁？}\\
\emph{Xìn de fājiànrén shì shéi ? Shōujiànrén shì shéi ?}\\
« Qui est l'expéditeur de la lettre ? Qui est le destinataire ? »\\
Réponse attendue : \zh{发件人是王伟，收件人是贵公司的人力资源部经理。} « L'expéditeur est Wang Wei ; le destinataire est le directeur des ressources humaines de votre honorable société. »

\textbf{2. (Global — Objet)} \zh{王伟想申请什么职位？}\\
\emph{Wáng Wěi xiǎng shēnqǐng shénme zhíwèi ?}\\
« À quel poste Wang Wei souhaite-t-il postuler ? »\\
Réponse : \zh{他想申请实习生的职位。} « Il souhaite postuler à un poste de stagiaire. »

\textbf{3. (Détaillé — Argumentaire)} \zh{文中提到的王伟的两个优势是什么？}\\
\emph{Wénzhōng tídào de Wáng Wěi de liǎng ge yōushì shì shénme ?}\\
« Quels sont les deux atouts de Wang Wei mentionnés dans le texte ? »\\
Réponse : \zh{第一，他通过了HSK四级考试；第二，他的学习成绩很好。} « Premièrement, il a réussi l'examen HSK niveau 4 ; deuxièmement, il a d'excellents résultats scolaires. »

\textbf{4. (Détaillé — Vocabulaire)} \zh{请用中文解释“烦请贵方予以考虑”的意思。}\\
\emph{Qǐng yòng Zhōngwén jiěshì “fánqǐng guìfāng yǔyǐ kǎolǜ” de yìsi.}\\
« Expliquez en chinois le sens de “fánqǐng guìfāng yǔyǐ kǎolǜ”. »\\
Réponse : \zh{意思是：请您们认真地考虑我的申请。} « Cela signifie : je vous prie de bien vouloir examiner ma candidature. »

\textbf{5. (Inférentiel — Opinion)} \zh{如果你是招聘经理，你会录取他吗？为什么？}\\
\emph{Rúguǒ nǐ shì zhāopìn jīnglǐ, nǐ huì lùqǔ tā ma ? Wèishénme ?}\\
« Si tu étais le directeur du recrutement, l'embaucherais-tu ? Pourquoi ? »\\
Réponse possible : \zh{我会录取他，因为他的汉语水平很高，而且学习很认真。} « Je l'embaucherais, car son niveau de chinois est élevé et il est sérieux dans ses études. »
\end{exemplebox}

\begin{attention}[Question inférentielle : justifier le registre]
À la question « Pourquoi \zh{烦请} et pas \zh{请} ? », la réponse doit mobiliser la notion de \cle{registre soutenu} : \zh{烦请} (fánqǐng, litt. « je vous ennuie en vous priant ») marque la déférence envers un supérieur hiérarchique inconnu ; \zh{请} seul serait correct mais trop neutre, presque sec, dans une demande de stage. Au Bac, citez toujours le mot du texte + sa fonction sociale.
\end{attention}

\courssub{Les collocations (\zh{搭配}) : les blocs de construction de la lettre}

Observer d'abord : dans la lettre de \zh{王伟}, on ne trouve jamais un verbe seul, mais des \cle{collocations}, c'est-à-dire des associations stables verbe + complément. En chinois formel, employer la mauvaise combinaison trahit immédiatement le débutant.

\begin{definition}[La collocation \zh{予以} (yǔyǐ)]
\zh{予以} est un verbe auxiliaire \textbf{formel} signifiant « accorder / donner ». Il ne se combine qu'avec des \textbf{verbes d'action abstraite} :
\zh{予以} + \zh{考虑} (examiner), \zh{批示} (approuver), \zh{回复} (répondre), \zh{支持} (soutenir), \zh{关注} (suivre), \zh{重视} (accorder de l'importance).\\
Exemple : \zh{请予以批示。} \emph{Qǐng yǔyǐ pīshì.} « Prière d'approuver (de donner vos instructions). »\\
\textbf{Exception / interdit} : jamais avec des verbes concrets — on ne dit pas \zh{予以吃饭} ni \zh{予以去}. Registre réservé à l'écrit administratif.
\end{definition}

\begin{center}
\begin{tabularx}{\linewidth}{|X|X|X|}
\hline
\rowcolor{popblueL} Collocation & Exemple (caractères + pinyin) & Traduction \\
\hline
\zh{通过} + 考试 & \zh{我通过了HSK四级考试。} \emph{Wǒ tōngguò le HSK sìjí kǎoshì.} & J'ai réussi l'examen HSK 4. \\
\hline
\zh{进行} + 交流 & \zh{我们进行了友好的交流。} \emph{Wǒmen jìnxíng le yǒuhǎo de jiāoliú.} & Nous avons eu un échange amical. \\
\hline
\zh{予以} + 考虑 & \zh{烦请贵方予以考虑。} \emph{Fánqǐng guìfāng yǔyǐ kǎolǜ.} & Je vous prie d'examiner (ma demande). \\
\hline
\zh{招聘} + 实习生 & \zh{贵公司正在招聘实习生吗？} \emph{Guì gōngsī zhèngzài zhāopìn shíxíshēng ma ?} & Votre société recrute-t-elle des stagiaires ? \\
\hline
\zh{感兴趣} + 于 & \zh{我对国际贸易很感兴趣。} \emph{Wǒ duì guójì màoyì hěn gǎnxìngqù.} & Je m'intéresse beaucoup au commerce international. \\
\hline
\end{tabularx}
\end{center}

\begin{attention}[Piège de traduction : \zh{关于……一事}]
Ne traduisez jamais \zh{一事} par « une affaire » (connotation policière ou judiciaire en français !). Traduisez par « la question de / le dossier de » ou omettez purement et simplement.\\
\zh{关于会议一事} → « Concernant la réunion » / « Au sujet de la réunion » — et non « À propos de l'affaire de la réunion ».
\end{attention}

\begin{savaistu}[Le cachet rouge, âme du document officiel chinois]
Dans les lettres très formelles (administration, diplomatie), la Chine utilise encore le papier à en-tête rouge (\zh{红头文件}, \emph{hóngtóu wénjiàn}) et le cachet officiel rouge (\zh{公章}, \emph{gōngzhāng}), apposé à cheval sur la signature et le papier. La signature manuscrite seule n'a pas valeur légale sans le cachet : c'est l'empreinte de l'institution qui authentifie le document, et non celle de la personne.
\end{savaistu}

\courssub{Exercice de complétion (cloze) : reconnaissance des caractères}

\begin{exoresolu}[Cloze — complétez sans pinyin]
Énoncé : complétez avec \zh{通过}、\zh{予以}、\zh{进行}、\zh{招聘}、\zh{感兴趣}.\\
1. 我\_\_\_\_了汉语水平考试。\\
2. 烦请贵方\_\_\_\_考虑。\\
3. 我们希望能与贵公司\_\_\_\_合作交流。\\
4. 贵公司正在\_\_\_\_实习生。\\
5. 我对中国文化非常\_\_\_\_。
\tcblower
Solution détaillée :\\
1. \zh{通过} — « réussir » un examen : collocation obligatoire \zh{通过考试}.\\
2. \zh{予以} — devant un verbe abstrait (\zh{考虑}), registre formel : \zh{予以考虑}.\\
3. \zh{进行} — verbe support des actions abstraites : \zh{进行交流} « mener un échange ».\\
4. \zh{招聘} — \zh{招聘实习生} « recruter des stagiaires ».\\
5. \zh{感兴趣} — structure \zh{对……感兴趣} « s'intéresser à ». Attention : \zh{感兴趣} se place après \zh{非常}, jamais \zh{非常感兴趣我}.
\end{exoresolu}

\begin{exercice}[Cloze d'entraînement — sans aide]
Complétez : 1. 关于实习\_\_\_\_事，我想…… 2. 鉴于我的成绩，我\_\_\_\_信自己适合这个职位。 3. 专此\_\_\_\_达我的敬意。 4. 顺\_\_\_\_商祺！
\end{exercice}

\begin{corrige}[Exercice — Cloze d'entraînement]
1. \zh{一} (\zh{关于……一事}) ; 2. \zh{相} (\zh{相信} « croire, être convaincu ») ; 3. \zh{表} (\zh{表达} « exprimer ») ; 4. \zh{颂} (\zh{顺颂商祺} formule finale commerciale).
\end{corrige}

\courssub{Traduction inversée (version) : du français vers le chinois formel}

\begin{methode}[Méthode de version FR vers CN en 4 étapes]
\begin{enumerate}
\item \textbf{Identifier le registre} : ici, toujours formel (lettre administrative).
\item \textbf{Choisir l'honorifique} : \zh{贵公司} (votre honorable société), \zh{尊敬的} (honorable, respecté).
\item \textbf{Structurer avec les connecteurs de la leçon} : \zh{关于}… \zh{鉴于}… \zh{烦请}… \zh{专此}… \zh{顺颂}…
\item \textbf{Vérifier la forme} : indentation de deux espaces en début de paragraphe, ponctuation pleine chasse (\zh{。}、\zh{，}、\zh{、}), aucun caractère traditionnel.
\end{enumerate}
\end{methode}

\begin{exoresolu}[Version — 5 phrases]
Énoncé : traduisez en chinois formel.\\
1. Je vous écris au sujet du stage dans votre société.\\
2. J'ai réussi l'examen HSK niveau 4.\\
3. Vu mes résultats scolaires, je pense convenir à ce poste.\\
4. Je vous prie de bien vouloir examiner ma candidature.\\
5. Je vous adresse mes salutations distinguées.
\tcblower
Solutions détaillées :\\
1. \zh{关于在贵公司实习一事，我特写此信。} \emph{Guānyú zài guì gōngsī shíxí yīshì, wǒ tè xiě cǐ xìn.} — connecteur \zh{关于} + \zh{一事} (non traduit en français).\\
2. \zh{我已经通过了HSK四级考试。} \emph{Wǒ yǐjīng tōngguò le HSK sìjí kǎoshì.} — collocation \zh{通过考试} + \zh{了} de l'accompli.\\
3. \zh{鉴于我的学习成绩，我认为自己适合这个职位。} \emph{Jiànyú wǒ de xuéxí chéngjì, wǒ rènwéi zìjǐ shìhé zhège zhíwèi.} — \zh{鉴于} en tête de phrase, jamais en fin.\\
4. \zh{烦请贵方予以考虑。} \emph{Fánqǐng guìfāng yǔyǐ kǎolǜ.} — formule figée à mémoriser telle quelle.\\
5. \zh{专此致敬！} \emph{Zhuāncǐ zhìjìng !} — salutation finale administrative standard.
\end{exoresolu}

\begin{attention}[Erreurs fréquentes des francophones en version]
\begin{itemize}
\item Traduire « au sujet de » par \zh{对于} au lieu de \zh{关于} en tête de lettre : \zh{关于} introduit le thème, \zh{对于} marque l'attitude.
\item Oublier le \zh{了} après \zh{通过} : \zh{我通过了考试} (action accomplie).
\item Placer le complément de lieu après le verbe : on dit \zh{在贵公司实习}, jamais \zh{实习在贵公司}.
\item Ponctuation française dans le texte chinois : utilisez \zh{，} et \zh{。}, pas « , » et « . ».
\end{itemize}
\end{attention}

\courssub{Auto-évaluation : la grille de critères}

Après chaque production (cloze ou version), évaluez-vous honnêtement avec cette grille, identique à celle de la section 6 :

\begin{center}
\begin{tabularx}{\linewidth}{|X|X|X|}
\hline
\rowcolor{popblueL} Critère & Ce qui est vérifié & Ma note (0 à 2) \\
\hline
\cle{Vocabulaire} & Collocations correctes (\zh{通过考试}, \zh{予以考虑}…) & \\
\hline
\cle{Structure} & Connecteurs bien placés (\zh{关于}…\zh{鉴于}…\zh{烦请}…) & \\
\hline
\cle{Caractères} & Simplifiés exacts, traits et clés respectés & \\
\hline
\cle{Ponctuation formelle} & Pleine chasse \zh{。，、}, indentation de 2 espaces & \\
\hline
\end{tabularx}
\end{center}

\begin{experience}[Relecture croisée en binôme]
Échangez votre version des 5 phrases avec votre voisin. Chacun corrige la copie de l'autre en utilisant uniquement la grille ci-dessus : un point perdu doit être justifié à voix haute en français (« tu as mis \zh{对于} au lieu de \zh{关于} »). Puis la classe entière compare ses choix pour la phrase 3 : combien de structures correctes différentes peut-on produire ?
\end{experience}

Vous voilà désormais capable de \emph{lire} finement une lettre formelle et de \emph{réemployer} ses blocs lexicaux dans vos propres phrases. Il ne reste plus qu'à franchir le dernier pas : rédiger votre propre lettre officielle complète — ce sera l'objet de la section 6, l'atelier d'écriture guidée.

\coursec{Production guidée : Atelier d'écriture « Ma première lettre officielle »}

Dans la section précédente, nous avons appris à \emph{lire} activement une lettre formelle : repérer les collocations (\zh{鉴于}, \zh{烦请}, \zh{专此函达}), comprendre l'intention de l'auteur et répondre à des questions de compréhension. Il est temps de passer \textbf{du passif à l'actif} : à votre tour de \emph{rédiger}. Cet atelier vous conduit, étape par étape, de l'idée à la lettre officielle prête à être envoyée — ou affichée au mur de la salle de chinois du lycée de Yaoundé.

\courssub{Choix du scénario : le tirage au sort}

L'enseignant prépare quatre fiches pliées. Chaque élève (ou binôme) tire un scénario au sort ; un élève rapide peut choisir librement. Chaque scénario correspond à une situation réelle de la vie d'un lycéen camerounais sinisant.

\begin{center}
\begin{tabularx}{\linewidth}{|X|X|X|X|}
\hline
\rowcolor{popblueL} N° & Scénario & Destinataire & Registre \\
\hline
1 & Demande de stage dans une entreprise chinoise à Douala & \zh{人事部经理} (responsable du personnel) & Très formel \\
\hline
2 & Demande d'information sur une bourse d'études & \zh{大学招生办公室} (bureau des admissions) & Formel \\
\hline
3 & Réclamation : produit défectueux acheté en ligne & \zh{客服经理} (responsable du service client) & Formel mais ferme \\
\hline
4 & Invitation officielle d'une personnalité à la fête du lycée & \zh{尊敬的校长} ou une autorité & Solennel \\
\hline
\end{tabularx}
\end{center}

\begin{attention}[Adapter le registre au destinataire]
En français, on écrit « Monsieur le Directeur » ou « Madame ». En chinois, on ne commence \textbf{jamais} une lettre formelle par \zh{亲爱的} (qīn'ài de, « cher/chère »), réservé aux proches. L'appel standard est \zh{尊敬的} + titre (\emph{zūnjìng de}, « honorable ») ou le nom du service + \zh{敬启} (formule d'appel d'ouverture). Une erreur de registre est plus grave qu'une faute de caractère : elle montre qu'on ne maîtrise pas la culture de la correspondance.
\end{attention}

\courssub{Étape 1 : Planification — la carte mentale}

Avant d'écrire le moindre caractère, on \textbf{planifie}. La lettre formelle chinoise repose sur cinq blocs obligatoires : le destinataire, l'objet, deux arguments, l'action demandée, la formule de politesse. Voici la carte mentale d'une demande de stage à Shenzhen.

\begin{popfigure}
\begin{tikzpicture}[mindmap, concept color=popblue!30,
  level 1/.style={sibling angle=72, level distance=3.6cm},
  level 2/.style={sibling angle=60, level distance=2.4cm}]
\node[concept] {Mon sujet : stage à Shenzhen}
  [clockwise from=90]
  child[concept color=popgreen!30] { node[concept] {Appel : \zh{人事部} / \zh{贵公司}} }
  child[concept color=poporange!30] { node[concept] {Corps}
      child { node[concept] {Argument 1 : \zh{HSK 4}} }
      child { node[concept] {Argument 2 : \zh{专业对口}} }
      child { node[concept] {Demande : \zh{烦请面试}} }
  }
  child[concept color=poppink!40] { node[concept] {Clôture : \zh{专此函达}} }
  child[concept color=poppurple!30] { node[concept] {Formule : \zh{顺颂商祺}} }
  child[concept color=popgold!40] { node[concept] {Signature : nom + \zh{启} + date} };
\end{tikzpicture}
\legende{Carte mentale de planification : les cinq blocs de la lettre formelle.}
\end{popfigure}

\begin{methode}[Planifier en cinq questions]
\begin{enumerate}
\item \textbf{À qui ?} Notez le destinataire exact : service (\zh{人事部}), titre (\zh{经理}), ou personne (\zh{尊敬的王校长}).
\item \textbf{Pourquoi ?} Formulez l'objet en une phrase : \zh{关于……一事} (« au sujet de … »).
\item \textbf{Quels arguments ?} Deux arguments maximum, reliés par \zh{鉴于} (« vu que ») ou \zh{首先……其次……} (« d'abord… ensuite… »).
\item \textbf{Quelle action ?} La demande précise introduite par \zh{烦请} + verbe (« je vous prie de … »).
\item \textbf{Quelle formule ?} Choisissez selon le destinataire : \zh{顺颂商祺} (entreprise), \zh{顺颂公祺} (administration), \zh{顺颂教祺} (enseignant), \zh{顺颂安祺} (aîné), \zh{顺颂拜上} (très respectueux).
\end{enumerate}
\end{methode}

\courssub{Étape 2 : Le brouillon sur papier quadrillé \zh{田字格}}

Le brouillon se rédige sur une feuille à cases \zh{田字格} (tiánzìgé, « cases en forme de champ »), un caractère par case. C'est l'occasion de soigner la \textbf{mise en page}, qui obéit à des règles strictes :

\begin{itemize}
\item L'\textbf{appel} (\zh{尊敬的人事部经理：}) s'écrit en haut, \textbf{à la marge gauche}, sans indentation, suivi des deux-points chinois \zh{：}.
\item La formule d'ouverture \zh{敬启！} (ou \zh{启！}) occupe la ligne suivante, alignée à gauche.
\item Chaque paragraphe du corps commence avec une \textbf{indentation de deux caractères} (deux cases vides).
\item La formule de clôture \zh{专此函达} termine le corps. Le vœu final (\zh{顺颂商祺}) se place sur la ligne suivante, aligné à gauche (sans indentation).
\item La \textbf{signature} (nom + \zh{启} / \zh{拜上} / \zh{敬上}) et la \textbf{date} (\zh{2025年10月16日}) se placent en bas \textbf{à droite}, la date sous la signature.
\end{itemize}

\begin{definition}[\zh{启} (qǐ), \zh{拜上} (bài shàng) et \zh{敬上} (jìng shàng) : trois signatures, trois degrés]
\zh{启} (litt. « ouvrir, présenter ») suit le nom de l'expéditeur : \zh{李华启} « présenté par Li Hua ». C'est la norme standard pour les lettres adressées à des supérieurs, des administrations ou des entreprises (égaux ou supérieurs).\\
\zh{拜上} (« présenter respectueusement ») est plus solennel que \zh{启} : il s'adresse aux aînés, aux professeurs, à un supérieur hiérarchique direct ou à une institution prestigieuse (université, ambassade).\\
\zh{敬上} (« respectueusement présenté ») est le plus solennel, réservé aux aînés très proches ou contextes très traditionnels. \textbf{Attention :} on évite \zh{敬上} si l'ouverture contient déjà \zh{敬启} (répétition de \zh{敬}).
\end{definition}

\begin{attention}[Cohérence appel / signature — Règle d'or]
L'ouverture standard est \zh{尊敬的XXX：\newline 敬启！}. La signature doit \textbf{éviter la répétition du morphème \zh{敬}}.
\begin{itemize}
\item Ouverture \zh{敬启} → Signature \zh{王伟启} ou \zh{王伟拜上} (recommandé pour université/autorité).
\item Ouverture \zh{启} (rare seul) → Signature \zh{王伟敬上} possible.
\item Ouverture \zh{尊敬的张经理：\newline 启！} → Signature \zh{王伟敬上} ou \zh{王伟拜上}.
\end{itemize}
Un seul marqueur de respect maximal par zone (ouverture / clôture) suffit.
\end{attention}

\courssub{Un modèle complet : la lettre de réclamation}

\begin{exemplebox}[Modèle rédigé — Scénario 3 : réclamation e-commerce]
\zh{尊敬的客服经理：}\\
\emph{Zūnjìng de kèfú jīnglǐ :} « Honorable responsable du service client : »\\
\zh{敬启！}\\
\emph{Jìng qǐ !} « Salutations respectueuses ! »\\
\zh{关于订单号20251015的手机屏幕破损一事，鉴于收货时包装完好，判定为物流责任。}\\
\emph{Guānyú dìngdān hào 20251015 de shǒujī píngmù pòsǔn yī shì, jiànyú shōuhuò shí bāozhuāng wánhǎo, pàndìng wéi wùliú zérèn.} « Au sujet de l'écran cassé du téléphone de la commande n° 20251015, vu que l'emballage était intact à la réception, la responsabilité incombe au transporteur. »\\
\zh{烦请贵方核实并予以退换货处理。}\\
\emph{Fán qǐng guìfāng héshí bìng yǔ yǐ tuìhuàn huò chǔlǐ.} « Je vous prie de vérifier et de procéder à un échange ou un remboursement. »\\
\zh{专此函达，}\\
\zh{顺颂商祺！}\\
\emph{Zhuāncǐ hándá, shùn sòng shāngqí !} « C'est l'objet de la présente ; je vous prie d'agréer mes salutations commerciales distinguées. »\\
\zh{李华启}\\
\emph{Lǐ Huá qǐ} « Li Hua »\\
\zh{2025年10月16日}\\
\emph{2025 nián 10 yuè 16 rì} « le 16 octobre 2025 »
\end{exemplebox}

Observez l'économie du texte : quatre phrases suffisent. La lettre formelle chinoise est \textbf{brève, factuelle, sans détour} — contrairement à la lettre administrative française qui multiplie les périphrases (« J'ai l'honneur de solliciter de votre haute bienveillance… »).

\begin{center}
\begin{tabularx}{\linewidth}{|X|X|X|X|}
\hline
\rowcolor{popblueL} Bloc & Chinois & Pinyin & Équivalent français \\
\hline
Appel & \zh{尊敬的客服经理} & zūnjìng de kèfú jīnglǐ & Monsieur le responsable \\
\hline
Ouverture & \zh{敬启} & jìng qǐ & Salutations respectueuses \\
\hline
Objet & \zh{关于……一事} & guānyú … yī shì & Objet : … \\
\hline
Argument & \zh{鉴于……} & jiànyú … & Vu que / considérant que \\
\hline
Demande & \zh{烦请贵方核实} & fán qǐng guìfāng héshí & Je vous prie de vérifier \\
\hline
Clôture & \zh{专此函达} & zhuāncǐ hándá & C'est l'objet de la présente \\
\hline
Formule & \zh{顺颂商祺} & shùn sòng shāngqí & Salutations distinguées (commerce) \\
\hline
Signature & \zh{李华启} & Lǐ Huá qǐ & Li Hua \\
\hline
Date & \zh{2025年10月16日} & 2025 nián 10 yuè 16 rì & Le 16/10/2025 \\
\hline
\end{tabularx}
\end{center}

\courssub{Étape 3 : Relecture par les pairs}

Les binômes échangent leurs brouillons et évaluent la lettre de l'autre à l'aide de la grille ci-dessous. Chaque critère vaut un point ; l'objectif est le « 5/5 ».

\begin{center}
\begin{tabularx}{\linewidth}{|X|X|X|}
\hline
\rowcolor{popgreenL} Critère & Question de contrôle & Oui / Non \\
\hline
Structure & Les 5 blocs (Appel, Ouverture, Corps, Formule, Signature/Date) sont-ils présents et dans l'ordre ? & \\
\hline
Vocabulaire honorifique & \zh{贵公司}, \zh{贵方}, \zh{贵校}, \zh{尊敬的} sont-ils bien employés selon le destinataire ? & \\
\hline
Caractères & Chaque caractère est-il correct, lisible et centré dans sa case ? & \\
\hline
Ponctuation & Deux-points \zh{：} après l'appel, \zh{。} en fin de phrase, pas d'espaces parasites ? & \\
\hline
Cohérence & Appel (\zh{敬启}) et signature (\zh{启} / \zh{拜上}) sont-ils assortis sans répétition de \zh{敬} ? & \\
\hline
\end{tabularx}
\end{center}

\begin{methode}[Check-list « Zéro Faute » avant de rendre]
\begin{itemize}
\item[$\square$] En-tête : service + \zh{贵公司}/\zh{贵校} ou \zh{尊敬的} + titre ?
\item[$\square$] Ouverture : \zh{敬启！} ou \zh{启！} sur ligne séparée, alignée à gauche ? (Jamais \zh{亲爱的} !)
\item[$\square$] Corps : indentation de 2 caractères ? Connecteurs formels (\cle{鉴于} / \cle{烦请} / \cle{专此函达}) ?
\item[$\square$] Formule : sur 2 lignes (\zh{专此函达，} + \zh{顺颂XX祺！}) alignées à gauche ? Adaptée au destinataire (\zh{商祺} / \zh{公祺} / \zh{教祺} / \zh{安祺}) ?
\item[$\square$] Signature : nom + \zh{启} (standard) ou \zh{拜上} (solennel) ? Date au format \zh{AAAA年MM月DD日} à droite sous la signature ?
\end{itemize}
\end{methode}

\begin{exoresolu}[Corriger la lettre de Rodrigue]
Énoncé. Rodrigue (scénario 2, bourse universitaire) a écrit : \zh{亲爱的老师：你好！我想申请奖学金。请给我信息。谢谢！王伟敬上。} Relevez quatre erreurs et réécrivez la lettre.
\tcblower
Solution détaillée. Erreur 1 : l'appel \zh{亲爱的老师} est familier ; il faut \zh{尊敬的招生办公室：} (service) ou \zh{尊敬的老师：} (personne). Erreur 2 : \zh{你好} est une salutation orale, absente d'une lettre formelle ; on écrit \zh{敬启！} sur la ligne suivante. Erreur 3 : \zh{请给我信息} est direct et impoli ; il faut \zh{烦请贵校提供奖学金申请的相关信息}. Erreur 4 : incohérence de signature — avec une ouverture \zh{敬启}, la signature \zh{敬上} répète \zh{敬} ; de plus, manquent la formule finale (\zh{顺颂教祺}) et la date. Version corrigée :\\
\zh{尊敬的招生办公室：}\\
\zh{敬启！}\\
\zh{关于贵校奖学金申请一事，鉴于本人已通过HSK四级，且成绩优良，烦请贵校提供申请材料及截止日期。}\\
\zh{专此函达，}\\
\zh{顺颂教祺！}\\
\zh{王伟拜上}\\
\zh{2025年10月16日}\\
\emph{Zūnjìng de zhāoshēng bàngōngshì : jìng qǐ ! Guānyú guìxiào jiǎngxuéjīn shēnqǐng yī shì, jiànyú běnrén yǐ tōngguò HSK sìjí, qiě chéngjì yōuliáng, fán qǐng guìxiào tígōng shēnqǐng cáiliào jí jiézhǐ rìqī. Zhuāncǐ hándá, shùn sòng jiàoqí ! Wáng Wěi bài shàng. 2025 nián 10 yuè 16 rì.} « Honorable bureau des admissions : salutations respectueuses ! Au sujet d'une demande de bourse de votre université, vu que j'ai réussi le HSK 4 avec d'excellents résultats, je vous prie de me communiquer le dossier et la date limite. C'est l'objet de la présente ; salutations distinguées. Wang Wei, le 16 octobre 2025. »
\end{exoresolu}

\courssub{Étapes 4 et 5 : Version finale et simulation « Enveloppe »}

\textbf{Étape 4 — La copie propre.} La lettre corrigée est recopiée avec soin : traits nets, caractères centrés, aucune rature. Elle sera soit affichée au mur de la classe, soit scannée pour le portfolio numérique OnBuch+.

\textbf{Étape 5 — L'enveloppe.} Dernier geste professionnel : rédiger l'adresse en chinois. Contrairement au français, l'adresse chinoise va \textbf{du plus grand au plus petit} : pays → province → ville → district/arrondissement → rue/parc → numéro/entité → nom. Sur la \textbf{face} (devant) : l'adresse du destinataire, suivie de \zh{收} (shōu, « à remettre à »). Au \textbf{dos} (derrière) : l'adresse de l'expéditeur, suivie de \zh{寄} (jì, « envoyé par »).

\begin{exemplebox}[L'enveloppe de la lettre de stage]
Face : \zh{广东省深圳市南山区科技园 华为公司人事部 王经理 收}\\
\emph{Guǎngdōng shěng Shēnzhèn shì Nánshān qū Kējìyuán, Huáwéi gōngsī rénshìbù, Wáng jīnglǐ shōu.} « Province du Guangdong, ville de Shenzhen, district de Nanshan, parc technologique — à remettre à M. Wang, directeur du personnel, société Huawei. »\\
Dos : \zh{喀麦隆雅温得第一高级中学 李华 寄}\\
\emph{Kāmàilóng Yǎwēndé dì-yī gāojí zhōngxué, Lǐ Huá jì.} « Envoyé par Li Hua, lycée classique n° 1 de Yaoundé, Cameroun. »
\end{exemplebox}

\begin{savaistu}[Et l'e-mail formel chinois ?]
Le courriel formel chinois conserve toute la structure de la lettre, avec trois ajouts. D'abord, un \textbf{objet} (\zh{主题}) très clair entre crochets : \zh{【实习申请】王伟-喀麦隆雅温得一中-HSK4} (« [Demande de stage] Wang Wei — Lycée de Yaoundé — HSK 4 »). Ensuite, le corps du mail reprend mot pour mot le corps de la lettre (appel \zh{敬启}, paragraphes indentés, formule, signature). Enfin, on ajoute une \textbf{signature automatique} (\zh{签名档}) avec téléphone, WeChat et parfois LinkedIn. En Chine, beaucoup de démarches se font d'ailleurs via WeChat : la politesse écrite y reste la même !
\end{savaistu}

\begin{experience}[Atelier en classe : de la carte mentale à l'enveloppe]
\begin{enumerate}
\item \textbf{Tirage au sort} (5 min) : chaque binôme tire un scénario et annonce son destinataire à la classe.
\item \textbf{Carte mentale} (10 min) : sur une feuille A4, reproduisez le modèle de la figure avec vos cinq blocs.
\item \textbf{Brouillon} (20 min) : rédaction sur papier \zh{田字格}, dictionnaire de caractères autorisé.
\item \textbf{Relecture croisée} (10 min) : échange des brouillons, grille à cinq critères, annotation au crayon.
\item \textbf{Copie propre + enveloppe} (15 min) : version finale et adresses face/dos ; les meilleures lettres rejoignent le mur d'exposition et le portfolio numérique.
\end{enumerate}
\end{experience}

\begin{exercice}[À vous d'écrire]
Choisissez le scénario 4 : votre lycée invite une personnalité (par exemple le directeur de l'Institut Confucius) à la fête de fin d'année. Rédigez la lettre complète (appel, ouverture \zh{敬启}, objet, deux arguments, demande, formule \zh{顺颂教祺}, signature, date) et l'adresse de l'enveloppe (face et dos).
\end{exercice}

\begin{corrige}[Exercice — éléments de réponse]
Appel : \zh{尊敬的杜主任：} ; ouverture : \zh{敬启！} ; corps (indentation 2 cases) : \zh{关于我校年终庆典一事，鉴于贵中心长期支持我校中文教学，且同学们渴望展示学习成果，烦请您拨冗出席。} ; clôture : \zh{专此函达，} (ligne suivante) \zh{顺颂教祺！} ; signature (alignée à droite) : \zh{雅温得一中学生会 拜上} + date \zh{2025年12月20日}. Vérifiez la cohérence appel (\zh{敬启}) / signature (\zh{拜上}), l'indentation de deux caractères pour les paragraphes du corps, et l'ordre des adresses sur l'enveloppe (du grand au petit).
\end{corrige}

\coursec{Extension Bac \& Évaluation formative (Complément examen)}

Après l'atelier d'écriture de la section 6, où vous avez produit votre première lettre officielle guidée, il est temps de passer à la vitesse supérieure : celle de l'\textbf{examen}. Au Baccalauréat camerounais (série A), l'épreuve écrite de chinois LV2 évalue exactement ce que nous venons de travailler : la \cle{compréhension écrite} d'un courrier administratif et l'\cle{expression écrite} d'une lettre formelle. Cette section vous donne les sujets types, les stratégies de gain de temps, la grille de notation officielle et des outils de remédiation. Objectif : transformer le format de la lettre formelle en points « faciles ».

\courssub{Analyse de sujets types Bac}

\begin{savaistu}[L'épreuve de chinois au Bac Cameroun (1ère A)]
L'épreuve de chinois au Baccalauréat comporte souvent une \textbf{compréhension de texte} portant sur un courrier administratif (lettre de demande, réponse officielle, avis) suivi de questions rédigées \emph{en français}, puis une \textbf{production écrite} : lettre ou courriel formel d'environ 80 à 120 caractères. La maîtrise du format rapporte des points « faciles » : la structure (6 blocs) et les formules de politesse représentent environ 8 points sur 20. Un élève qui connaît son squelette de lettre part donc avec une longueur d'avance.
\end{savaistu}

\begin{popfigure}
\begin{tikzpicture}
\node[align=left, draw=popink, fill=popink!10, rounded corners, text width=10cm] { \textbf{Sujet Type Bac Blanc -- Expression Écrite (20 pts)} \\ \textit{Vous êtes Paul BIYA (élève), vous écrivez au Directeur de l'Institut Confucius de Yaoundé pour demander une attestation de niveau HSK 4 perdue. Rédigez la lettre formelle complète (env. 100 caractères).} \\ \textbf{Critères de réussite :} \\ 1. Format impeccable (6 blocs). \\ 2. Vocabulaire : 补办, 遗失, 烦请, 补发, 顺颂教祺. \\ 3. Cohérence appel/signature. \\ 4. Caractères lisibles. } ;
\end{tikzpicture}
\legende{Sujet type Bac Blanc : la demande d'attestation, grand classique de l'épreuve.}
\end{popfigure}

\begin{methode}[Stratégie « Gain de temps » en examen]
\begin{enumerate}
\item \textbf{Recopier le squelette pré-mémorisé} : les 6 blocs (appel, salutation, corps, formule finale, signature, date) sont posés sur la copie en moins de deux minutes.
\item \textbf{Remplir les variables} : Destinataire (\zh{贵校} / \zh{贵中心} / \zh{贵馆}), Objet (\zh{关于……一事}), Arguments (\zh{遗失了……，需要补办……}).
\item \textbf{Insérer 3 formules magiques} : \zh{关于……一事}, \zh{烦请……}, \zh{专此函达} ou \zh{顺颂……祺}.
\item \textbf{Relire les caractères} à haut risque : \zh{贵}, \zh{敬}, \zh{启}, \zh{祺} (traits manquants ou confusions).
\end{enumerate}
\end{methode}

\begin{attention}[Sujet piège : la lettre amicale déguisée]
Si le sujet demande une « lettre à un ami chinois pour lui donner des nouvelles », ce n'est \textbf{pas} une lettre formelle ! N'utilisez alors \emph{jamais} \zh{贵公司}, \zh{顺颂商祺} ni \zh{烦请}. Employez : \zh{亲爱的……} (qīn'ài de..., « cher/chère... »), \zh{祝好} ou \zh{祝安好} en formule finale, et \zh{请} au lieu de \zh{烦请}. Savoir distinguer les registres de langue est explicitement noté au Bac : confondre lettre amicale et lettre administrative fait perdre jusqu'à 5 points de « Structure/Format ».
\end{attention}

\begin{exemplebox}[Phrases « passe-partout » à mémoriser par cœur]
\zh{关于贵方来函所述事宜，现函复如下。} \emph{(Guānyú guìfāng láihán suǒ shù shìyí, xiàn hánfù rúxià.)} « Au sujet des points mentionnés dans votre courrier, je réponds par la présente. »\\[4pt]
\zh{烦请查收附件，并予以批示为盼。} \emph{(Fánqǐng cháshōu fùjiàn, bìng yǔyǐ pīshì wéipàn.)} « Prière de vérifier la pièce jointe et de bien vouloir l'approuver. »\\[4pt]
\zh{即此颂达，顺颂工作顺利！} \emph{(Jícǐ sòngdá, shùnsòng gōngzuò shùnlì!)} « Je m'arrête là, en vous souhaitant bon travail ! » (moins formel que \zh{顺颂商祺}, utile pour un destinataire non commercial).
\end{exemplebox}

\courssub{Exercice type « QCM Caractères » : chasser les homophones}

Au Bac, un caractère faux dans une formule de politesse coûte cher. Les correcteurs ciblent les \cle{homophones} et les caractères graphiquement proches. Entraînons-nous sur trois paires classiques.

\begin{center}
\begin{tabularx}{\linewidth}{|X|X|X|X|}
\hline
\rowcolor{popblueL} Caractères & Pinyin & Nature & Traduction \\
\hline
贵 & guì & adjectif & honorable, cher (votre...) \\
\hline
桂 & guì & nom & osmanthe (plante) ; Guangxi \\
\hline
批 & pī & verbe & approuver, annoter \\
\hline
启 & qǐ & verbe & ouvrir, adresser (formel) \\
\hline
绮 & qǐ & adjectif & beau, orné (littéraire) \\
\hline
祺 & qí & nom & bonheur, auspice \\
\hline
\end{tabularx}
\end{center}

\begin{exoresolu}[QCM Caractères : choisir le bon sinogramme]
Énoncé. Complétez avec le bon caractère : (1) \zh{＿公司} (votre honorable société) : a. 贵 b. 桂 ; (2) \zh{＿示} (approbation) : a. 批 b. 比 ; (3) \zh{＿事} (ouvrir la lettre, formel) : a. 绮 b. 启 ; (4) \zh{顺颂商＿} : a. 旗 b. 祺.
\tcblower
Solution détaillée. (1) \textbf{a. 贵} : \zh{贵公司} (guì gōngsī) est la formule de déférence ; \zh{桂} désigne une plante et la province du Guangxi — écrire \zh{桂公司} serait comique et fautif. (2) \textbf{a. 批} : \zh{批示} (pīshì) « approbation, instruction » ; la clé de la main (扌) indique l'action. (3) \textbf{b. 启} : \zh{启事} / appel \zh{……启} ; \zh{绮} (qǐ, « orné ») n'apparaît jamais dans la correspondance. (4) \textbf{b. 祺} : \zh{顺颂商祺} (shùnsòng shāngqí) « salutations commerciales distinguées » ; \zh{旗} (qí, « drapeau ») est un homophone parfait — piège classique du Bac.
\end{exoresolu}

\begin{exercice}[QCM Caractères -- à vous]
Choisissez le bon caractère : (1) \zh{＿请} (je vous prie de...) : a. 烦 b. 繁 ; (2) \zh{补＿} (rééditer un document) : a. 发 b. 法 ; (3) \zh{＿失} (égarer) : a. 遗 b. 遣 ; (4) \zh{＿件} (pièce jointe) : a. 附 b. 符.
\end{exercice}

\begin{corrige}[Exercice QCM Caractères]
(1) \textbf{a. 烦} : \zh{烦请} (fánqǐng) « j'ose vous prier de... » ; \zh{繁} (fán) signifie « nombreux, compliqué ». (2) \textbf{a. 发} : \zh{补发} (bǔfā) « délivrer à nouveau » ; \zh{法} (fǎ) = « loi ». (3) \textbf{a. 遗} : \zh{遗失} (yíshī) « perdre, égarer » (registre formel) ; \zh{遣} (qiǎn) n'a ni le même son ni le même sens. (4) \textbf{a. 附} : \zh{附件} (fùjiàn) « pièce jointe » ; \zh{符} (fú) = « symbole ».
\end{corrige}

\courssub{Exercice type « Remise en ordre » : reconstruire la lettre}

L'épreuve de compréhension propose parfois des phrases mélangées à réorganiser. La logique est toujours la même : \textbf{appel → objet → faits → demande → formule finale}. Observons d'abord une lettre modèle complète, puis entraînons-nous.

\begin{textebox}[Modèle : demande de visa d'études X1]
尊敬的使馆签证处负责人：您好！\\
\emph{Zūnjìng de shǐguǎn qiānzhèngchù fùzérén : nín hǎo !}\\
« Monsieur le Responsable du service des visas de l'Ambassade, bonjour ! »\\[4pt]
关于申请X1学习签证一事，现致函如下。\\
\emph{Guānyú shēnqǐng X1 xuéxí qiānzhèng yī shì, xiàn zhìhán rúxià.}\\
« Concernant ma demande de visa d'études X1, je vous écris la présente lettre. »\\[4pt]
我叫玛丽，是雅温得第一高中的学生，已被北京语言大学录取。\\
\emph{Wǒ jiào Mǎlì, shì Yǎwēndé dì-yī gāozhōng de xuésheng, yǐ bèi Běijīng Yǔyán Dàxué lùqǔ.}\\
« Je m'appelle Marie, élève au lycée de Yaoundé, déjà admise à l'Université des langues de Pékin. »\\[4pt]
烦请告知办理该签证所需的全部材料，以便及时准备。\\
\emph{Fánqǐng gàozhī bànlǐ gāi qiānzhèng suǒ xū de quánbù cáiliào, yǐbiàn jíshí zhǔnbèi.}\\
« Je vous prie de m'indiquer toutes les pièces nécessaires à ce visa, afin de les préparer à temps. »\\[4pt]
专此函达，顺颂工作顺利！\\
\emph{Zhuāncǐ hándá, shùnsòng gōngzuò shùnlì !}\\
« Par la présente, salutations distinguées et bon travail ! »\\[4pt]
此致\\
敬礼！\\
学生：玛丽\\
2025年3月10日\\
\emph{Cǐzhì jìnglǐ ! Xuésheng : Mǎlì, 2025 nián 3 yuè 10 rì.} « Veuillez agréer mes salutations respectueuses ! L'élève : Marie, le 10 mars 2025. »
\end{textebox}

\begin{exoresolu}[Remise en ordre : 5 phrases mélangées]
Énoncé. Remettez dans l'ordre : (A) \zh{烦请补发该证明，不胜感激。} (B) \zh{尊敬的校长：您好！} (C) \zh{此致敬礼！学生：保罗} (D) \zh{关于遗失成绩单一事，现致函如下。} (E) \zh{我不慎遗失了去年的成绩单，现急需该文件申请奖学金。}
\tcblower
Solution détaillée. Ordre correct : \textbf{B → D → E → A → C}. Logique : (B) l'appel ouvre toujours la lettre ; (D) la formule \zh{关于……一事} annonce l'objet juste après ; (E) expose les faits (perte du bulletin, besoin urgent) ; (A) formule la demande avec \zh{烦请} ; (C) clôt avec la salutation et la signature. Astuce d'examen : repérez les « marqueurs de position » — \zh{尊敬的} = début, \zh{关于……一事} = deuxième phrase, \zh{烦请} = avant-dernière, \zh{此致敬礼} = fin.
\end{exoresolu}

\courssub{Exercice type « Rédaction courte » (80-100 caractères)}

\begin{exercice}[Rédaction : lettre à l'Ambassade de Chine]
Écrivez une lettre à l'Ambassade de Chine au Cameroun pour demander les documents nécessaires au visa X1 (études longues). 80 à 100 caractères. Utilisez au moins : \zh{关于……一事}, \zh{烦请}, \zh{材料}, \zh{此致敬礼}.
\end{exercice}

\begin{corrige}[Exercice de rédaction -- modèle de copie (environ 95 caractères)]
\zh{尊敬的签证处负责人：您好！关于申请X1学习签证一事，现致函咨询。我已被上海师范大学录取，计划今年九月入学。烦请告知办理该签证所需的全部材料及费用，以便及时准备。专此函达，此致敬礼！申请人：恩东，2025年4月2日。}\\[4pt]
\emph{Zūnjìng de qiānzhèngchù fùzérén : nín hǎo ! Guānyú shēnqǐng X1 xuéxí qiānzhèng yī shì, xiàn zhìhán zīxún. Wǒ yǐ bèi Shànghǎi Shīfàn Dàxué lùqǔ, jìhuà jīnnián jiǔyuè rùxué. Fánqǐng gàozhī bànlǐ gāi qiānzhèng suǒ xū de quánbù cáiliào jí fèiyòng, yǐbiàn jíshí zhǔnbèi. Zhuāncǐ hándá, cǐzhì jìnglǐ ! Shēnqǐngrén : Ēndōng, 2025 nián 4 yuè 2 rì.}\\
« Monsieur le Responsable des visas, bonjour ! Concernant ma demande de visa d'études X1, je vous écris pour me renseigner. Admis à l'Université normale de Shanghai, je prévois ma rentrée en septembre. Je vous prie de m'indiquer toutes les pièces et les frais nécessaires, afin de les préparer à temps. Par la présente, salutations respectueuses ! Le demandeur : Ndong, le 2 avril 2025. »
\end{corrige}

\courssub{Grille d'évaluation officielle MINESEC adaptée}

\begin{propriete}[Grille de notation de la lettre formelle (Total : 20 pts)]
\begin{center}
\begin{tabularx}{\linewidth}{|X|X|X|}
\hline
\rowcolor{popblueL} Critère & Points & Ce que le correcteur vérifie \\
\hline
Contenu & 5 & Objet clair, faits précis, demande explicite \\
\hline
Structure / Format & 5 & Les 6 blocs présents et bien placés, indentation \\
\hline
Langue / Vocabulaire & 5 & Formules (\zh{烦请}, \zh{予以}, \zh{鉴于}), connecteurs, registre soutenu \\
\hline
Caractères / Calligraphie & 5 & Sinogrammes corrects, lisibles, proportions respectées \\
\hline
\end{tabularx}
\end{center}
Conséquence stratégique : même avec un contenu moyen, un format parfait et des caractères propres rapportent déjà 10 points.
\end{propriete}

\begin{propriete}[Grille d'auto-évaluation élève (à coller dans le cahier)]
Cochez avant de rendre toute copie :\\
{[}\ \ {]} Je connais les 6 parties de la lettre formelle.\\
{[}\ \ {]} J'écris \zh{贵公司} sans hésiter.\\
{[}\ \ {]} Je distingue \zh{启} / \zh{敬上} / \zh{拜上}.\\
{[}\ \ {]} J'utilise \zh{烦请} / \zh{予以} / \zh{鉴于}.\\
{[}\ \ {]} Mes caractères tiennent dans le carré.\\
{[}\ \ {]} Mon indentation est correcte (deux caractères en retrait au début du corps).
\end{propriete}

\courssub{Remédiation différenciée}

\textbf{Pour les élèves faibles en caractères : la fiche « Radical Power ».} Chaque semaine, 5 caractères du thème sont décomposés par leur clé : \zh{贵} (clé 贝 bèi, le coquillage → la valeur), \zh{烦} (clé 火 huǒ, le feu → l'irritation), \zh{遗} (clé 辶 chuò, la marche → ce qui est laissé en chemin), \zh{附} (clé 阝 fù, la colline), \zh{祺} (clé 礻 shì, l'autel → le bonheur auspicieux). Méthode : écrire chaque caractère 5 fois en nommant la clé à voix haute, puis le réécrire de mémoire dans un carré.

\textbf{Pour les élèves faibles en structure : les gabarits « Phrases toutes faites ».} L'élève ne rédige pas : il \emph{remplit} un squelette : \zh{尊敬的＿＿：您好！关于＿＿一事，现致函如下。……烦请＿＿，不胜感激。此致敬礼！＿＿（签名）＿＿年＿＿月＿＿日。} Trois trous à compléter (destinataire, objet, demande) suffisent pour obtenir une lettre correcte et notée.

\begin{attention}[Dernières consignes avant l'épreuve]
Ne traduisez jamais mot à mot depuis le français : « Je vous serais reconnaissant » devient \zh{不胜感激} (bùshèng gǎnjī), pas une construction calquée. Vérifiez que la date suit l'ordre chinois année-mois-jour (\zh{2025年4月2日}), jamais l'ordre français. Enfin, un caractère illisible vaut un caractère faux : écrivez lentement, dans des carrés imaginaires.
\end{attention}

Vous disposez désormais de tout l'arsenal : squelette, formules, stratégies et grille de notation. La lettre formelle n'a plus de secret pour vous — au Bac, elle deviendra votre source de points la plus sûre. 加油！(Jiāyóu ! Courage !)

\newpage

\coursec{Activité d'intégration}

\courssub{Objectifs de l'activité}

À l'issue de cette activité d'intégration, l'élève sera capable de :
\begin{itemize}
\item rédiger de manière autonome une lettre formelle complète en chinois (formule d'appel, objet, corps structuré, formule de politesse, date et signature) ;
\item mobiliser le vocabulaire de la correspondance administrative étudié dans la leçon (\zh{申请}, \zh{证明}, \zh{办理}, \zh{回复}…) avec l'orthographe correcte des caractères ;
\item employer les connecteurs et structures types du corps de lettre (\zh{因为……所以……}, \zh{我想……}, \zh{请您……}) ;
\item lire, évaluer et « traiter » la demande d'un pair en adoptant le point de vue de l'administrateur destinataire ;
\item rédiger une réponse administrative courte (accord, refus poli ou demande de complément) dans un registre formel ;
\item présenter oralement un échange épistolaire et évaluer la production d'autrui à l'aide d'une grille critériée.
\end{itemize}

\courssub{Situation}

La classe de Première A se transforme pour une heure en « Centre de Services Administratifs » du lycée de Yaoundé. Au guichet unique, on peut obtenir tous les documents officiels : relevés de notes, certificats de scolarité, attestations, réponses aux réclamations. Mais attention : ici, \textbf{tout se passe en chinois}, car le centre fonctionne en partenariat avec l'Institut Confucius, qui exige des dossiers rédigés en chinois simplifié.

Chaque élève tire au sort un \textbf{Ticket Mission}. Il devra rédiger la lettre formelle correspondante, puis la transmettre à un camarade qui jouera le rôle de l'administrateur du guichet. Voici les quatre tickets possibles :

\begin{textebox}[Les quatre « Tickets Mission » du Guichet Unique]
\textbf{Ticket 1 — Demande de relevé de notes}\\
\zh{任务：给校长写一封信，申请你的成绩单。}\\
\emph{Rènwù : gěi xiàozhǎng xiě yì fēng xìn, shēnqǐng nǐ de chéngjìdān.}\\
« Mission : écris une lettre au proviseur pour demander ton relevé de notes. »\\[4pt]
\textbf{Ticket 2 — Réclamation d'un colis}\\
\zh{任务：给邮局写一封信，你的包裹还没到，请他们帮你找。}\\
\emph{Rènwù : gěi yóujú xiě yì fēng xìn, nǐ de bāoguǒ hái méi dào, qǐng tāmen bāng nǐ zhǎo.}\\
« Mission : écris une lettre à la poste : ton colis n'est pas encore arrivé, demande-leur de t'aider à le retrouver. »\\[4pt]
\textbf{Ticket 3 — Candidature bénévole}\\
\zh{任务：给孔子学院写一封信，你想当志愿者。}\\
\emph{Rènwù : gěi Kǒngzǐ Xuéyuàn xiě yì fēng xìn, nǐ xiǎng dāng zhìyuànzhě.}\\
« Mission : écris une lettre à l'Institut Confucius : tu veux devenir bénévole. »\\[4pt]
\textbf{Ticket 4 — Certificat de scolarité pour visa}\\
\zh{任务：给学校写一封信，申请在学证明，因为你要办签证。}\\
\emph{Rènwù : gěi xuéxiào xiě yì fēng xìn, shēnqǐng zàixué zhèngmíng, yīnwèi nǐ yào bàn qiānzhèng.}\\
« Mission : écris une lettre à l'école pour demander un certificat de scolarité, car tu dois faire une demande de visa. »
\end{textebox}

\courssub{Consignes}

\textbf{Étape 1 — Rédaction de la lettre (20 minutes, travail individuel).}
\begin{enumerate}
\item Lis attentivement ton Ticket Mission et souligne : le \cle{destinataire}, l'\cle{objet} de la demande et la \cle{raison}.
\item Rédige ta lettre formelle en chinois, de 8 à 12 phrases, en respectant l'architecture étudiée : formule d'appel (\zh{尊敬的……：}), présentation de soi, objet de la demande, justification, formule de politesse finale (\zh{谢谢您的帮助！}), date et signature.
\item Utilise \textbf{au moins cinq mots} du glossaire de la correspondance administrative et \textbf{au moins deux connecteurs} (\zh{因为……所以……}, \zh{首先}, \zh{然后}, \zh{另外}).
\item Soigne l'écriture des caractères : ordre des traits correct, proportions équilibrées. Les caractères-clés de la leçon (\zh{请}, \zh{谢}, \zh{信}, \zh{证}) doivent être irréprochables.
\end{enumerate}

\textbf{Étape 2 — Traitement du dossier par l'administrateur (10 minutes, travail en binôme).}
\begin{enumerate}
\item Échange ta lettre avec ton voisin. Tu deviens l'\cle{administrateur} du guichet.
\item Vérifie le dossier reçu à l'aide de la grille « Zéro Faute » : la formule d'appel est-elle présente ? La demande est-elle claire ? Le registre est-il poli ? Y a-t-il des fautes de caractères ou de tons dans le pinyin ?
\item Rédige une \cle{réponse administrative} courte (4 à 6 phrases) : accord (\zh{我们同意您的申请。}), refus poli (\zh{很抱歉，……}) ou demande de complément (\zh{请您再给我们……}).
\end{enumerate}

\textbf{Étape 3 — Restitution orale et évaluation collective (15 minutes).}
\begin{enumerate}
\item Trois binômes présentent leur échange au tableau : lecture de la lettre à voix haute (attention aux tons !), puis lecture de la réponse.
\item La classe évalue chaque production avec la grille « Zéro Faute » affichée au tableau et propose des corrections.
\end{enumerate}

\begin{attention}[La grille « Zéro Faute »]
\begin{itemize}
\item \textbf{Forme} : appel \zh{尊敬的} + destinataire + deux-points chinois \zh{：} ; date et signature en bas à droite.
\item \textbf{Clarté} : on comprend la demande dès la deuxième phrase.
\item \textbf{Politesse} : emploi de \zh{您} (et non \zh{你}) pour s'adresser au destinataire ; \zh{请} avant chaque requête.
\item \textbf{Langue} : ordre des mots correct (Sujet + Verbe + Complément), tons corrects, caractères bien formés.
\end{itemize}
\end{attention}

\courssub{Ressources à mobiliser}

\begin{center}
\begin{tabularx}{\linewidth}{|X|X|X|X|}
\hline
\rowcolor{popblueL} Caractères & Pinyin & Nature & Traduction \\
\hline
\zh{尊敬的} & zūnjìng de & adj. & honorable, respecté (appel) \\
\hline
\zh{申请} & shēnqǐng & v. / n. & demander, faire une demande \\
\hline
\zh{证明} & zhèngmíng & n. / v. & certificat, attestation ; prouver \\
\hline
\zh{成绩单} & chéngjìdān & n. & relevé de notes \\
\hline
\zh{办理} & bànlǐ & v. & effectuer (des démarches) \\
\hline
\zh{签证} & qiānzhèng & n. & visa \\
\hline
\zh{志愿者} & zhìyuànzhě & n. & bénévole, volontaire \\
\hline
\zh{包裹} & bāoguǒ & n. & colis, paquet \\
\hline
\zh{回复} & huífù & v. / n. & répondre ; réponse \\
\hline
\zh{很抱歉} & hěn bàoqiàn & expr. & je suis désolé (registre formel) \\
\hline
\end{tabularx}
\end{center}

\begin{aretenir}[Structures types à réemployer]
\begin{itemize}
\item Objet de la demande : \zh{我想申请……} \emph{(Wǒ xiǎng shēnqǐng……)} « Je voudrais demander… »
\item Justification : \zh{因为……，所以……} \emph{(yīnwèi……, suǒyǐ……)} « parce que…, donc… »
\item Requête polie : \zh{请您 + verbe} \emph{(qǐng nín……)} « je vous prie de… »
\item Clôture : \zh{谢谢您的帮助！} \emph{(Xièxie nín de bāngzhù !)} « Merci de votre aide ! »
\item Réponse administrative : \zh{我们同意您的申请。} \emph{(Wǒmen tóngyì nín de shēnqǐng.)} « Nous acceptons votre demande. »
\end{itemize}
\end{aretenir}

\courssub{Proposition de production et corrigé commenté}

\begin{exoresolu}[Modèle complet : Ticket 4 et sa réponse administrative]
\textbf{Étape 1 — La lettre de l'élève (demande de certificat de scolarité).}
\tcblower
\zh{尊敬的校长先生：}\\
\emph{Zūnjìng de xiàozhǎng xiānsheng :}\\
« Monsieur le Proviseur, »\\[3pt]
\zh{您好！我叫玛丽，是贵校一年级A班的学生。因为我要办理签证，所以想申请一份在学证明。}\\
\emph{Nín hǎo ! Wǒ jiào Mǎlì, shì guìxiào yì niánjí A bān de xuésheng. Yīnwèi wǒ yào bànlǐ qiānzhèng, suǒyǐ xiǎng shēnqǐng yí fèn zàixué zhèngmíng.}\\
« Bonjour ! Je m'appelle Marie, je suis élève de la classe de Première A de votre établissement. Comme je dois faire une demande de visa, je voudrais demander un certificat de scolarité. »\\[3pt]
\zh{请您帮我办理这个证明。谢谢您的帮助！}\\
\emph{Qǐng nín bāng wǒ bànlǐ zhège zhèngmíng. Xièxie nín de bāngzhù !}\\
« Je vous prie de bien vouloir m'établir ce certificat. Merci de votre aide ! »\\[3pt]
\zh{学生：玛丽　　二〇二五年五月十日}\\
\emph{Xuésheng : Mǎlì　　èr líng èr wǔ nián wǔ yuè shí rì}\\
« L'élève : Marie — le 10 mai 2025 »\\[6pt]
\textbf{Commentaire des choix de langue :}
\begin{itemize}
\item L'appel \zh{尊敬的校长先生：} respecte la norme : adjectif honorifique + fonction + deux-points chinois pleine chasse.
\item \zh{贵校} (« votre honorable établissement ») est le terme de politesse attendu ; dire \zh{你的学校} serait trop familier.
\item La justification utilise la structure \zh{因为……所以……} : en chinois, les deux éléments peuvent coexister dans la même phrase, contrairement au français où « parce que… donc… » serait incorrect.
\item \zh{一份} : classificateur \zh{份} obligatoire devant \zh{在学证明} (document administratif).
\item La requête \zh{请您帮我……} place \zh{请} avant \zh{您} : registre formel irréprochable.
\end{itemize}
\textbf{Étape 2 — La réponse de l'administrateur (accord).}\\[3pt]
\zh{玛丽同学：}\\
\emph{Mǎlì tóngxué :}\\
« Chère élève Marie, »\\[3pt]
\zh{我们收到了您的信。我们同意您的申请。请您星期五上午来学校办公室拿在学证明。}\\
\emph{Wǒmen shōudào le nín de xìn. Wǒmen tóngyì nín de shēnqǐng. Qǐng nín xīngqīwǔ shàngwǔ lái xuéxiào bàngōngshì ná zàixué zhèngmíng.}\\
« Nous avons bien reçu votre lettre. Nous acceptons votre demande. Venez chercher votre certificat de scolarité vendredi matin au bureau de l'école. »\\[3pt]
\zh{学校办公室　　二〇二五年五月十二日}\\
\emph{Xuéxiào bàngōngshì　　èr líng èr wǔ nián wǔ yuè shí'èr rì}\\
« Le bureau de l'école — le 12 mai 2025 »\\[6pt]
\textbf{Commentaire :} la réponse reprend les marques du registre administratif : \zh{收到} (« accuser réception »), \zh{同意您的申请} (formule d'accord type), et conserve le vouvoiement \zh{您} envers l'élève, ce qui est la marque du style officiel chinois.
\end{exoresolu}

\courssub{Bilan de l'activité}

Cette simulation du « Guichet Unique » a permis de réunir en une seule tâche authentique l'ensemble des savoir-faire de la leçon : \textbf{comprendre} une consigne en chinois (le Ticket Mission), \textbf{mobiliser} le lexique de la correspondance administrative, \textbf{structurer} une lettre formelle selon l'architecture rhétorique chinoise, \textbf{écrire} les caractères avec soin, puis \textbf{changer de rôle} pour lire, évaluer et répondre à la demande d'autrui. L'élève a également découvert que la lettre formelle n'est pas un exercice scolaire artificiel : obtenir un certificat, un visa ou un stage suppose réellement de maîtriser ce genre de texte, au Cameroun comme en Chine. La grille « Zéro Faute », réutilisable à chaque production écrite, constitue désormais un outil d'auto-évaluation permanent. Pour le baccalauréat, l'élève retiendra trois réflexes : un appel honorifique (\zh{尊敬的}), une demande annoncée clairement dès la deuxième phrase, et une clôture polie (\zh{谢谢您的帮助！}) suivie de la date et de la signature.

\newpage

\coursec{Méthodes et exercices résolus}

\coursec{Leçon 39 — Vocabulaire et compréhension de texte : rédaction d'une lettre formelle}

\courssub{Mise en situation \& Texte support : Lettre de demande de stage}

\begin{textebox}[Lettre de demande de stage — Contexte : élève de Première A, Lycée de Yaoundé]
尊敬的王经理：您好！\\
Zūnjìng de Wáng jīnglǐ : Nín hǎo !\\
Honorable Directeur Wang, bonjour !\\
我是雅温得第二中学高二年级的学生，我叫李华。\\
Wǒ shì Yǎwēndé Dì-èr Zhōngxué gāo èr niánjí de xuésheng, wǒ jiào Lǐ Huá.\\
Je suis élève de Première (Seconde année du lycée) au Lycée n°2 de Yaoundé, je m'appelle Li Hua.\\
我写信是为了申请贵公司的暑期实习。\\
Wǒ xiěxìn shì wèile shēnqǐng guì gōngsī de shǔqī shíxí.\\
Je vous écris pour solliciter un stage d'été au sein de votre honorable entreprise.\\
我学习汉语已经两年了，我的汉语说得比较好，也很喜欢电脑和媒体工作。\\
Wǒ xuéxí Hànyǔ yǐjīng liǎng nián le, wǒ de Hànyǔ shuō de bǐjiào hǎo, yě hěn xǐhuan diànnǎo hé méitǐ gōngzuò.\\
J'étudie le chinois depuis deux ans déjà, je parle chinois assez bien, et j'aime beaucoup l'informatique et le travail dans les médias.\\
谢谢您的考虑！此致敬礼！\\
Xièxie nín de kǎolǜ ! Cǐ zhì jìnglǐ !\\
Merci de l'attention que vous porterez à ma demande ! Veuillez agréer, Monsieur le Directeur, mes salutations respectueuses !\\
李华，六月十五日\\
Lǐ Huá, liù yuè shíwǔ rì\\
Li Hua, le 15 juin
\end{textebox}

\begin{aretenir}[Repères socioculturels : Cameroun / Chine]
\begin{itemize}
\item \textbf{Appellation} : En Chine, on utilise le nom de famille + titre/fonction (ex. \zh{王经理} Wáng jīnglǐ, « Directeur Wang »). Au Cameroun, on dit « Monsieur le Directeur » ou « Monsieur [Nom] ».
\item \textbf{Formule de clôture} : \zh{此致敬礼} (Cǐ zhì jìnglǐ) est l'équivalent standard de « Veuillez agréer... salutations distinguées ». Elle s'écrit sur deux lignes : \zh{此致} en fin de paragraphe, \zh{敬礼} en début de ligne suivante, alignée à gauche.
\item \textbf{Date} : Format chinois \zh{年月日} (année-mois-jour) ou simplement \zh{月日} dans les lettres courantes. Ici : \zh{六月十五日} (15 juin).
\item \textbf{Politesse} : L'usage systématique de \zh{您} (nín) et du préfixe \zh{贵} (guì) devant le nom de l'institution du destinataire (\zh{贵公司}, \zh{贵校}) est obligatoire en contexte formel.
\end{itemize}
\end{aretenir}

\courssub{Glossaire thématique : Vocabulaire de la correspondance administrative}

\begin{center}
\begin{tabularx}{\linewidth}{|X|X|X|X|}
\hline
\rowcolor{popblueL} \textbf{Caractères} & \textbf{Pinyin} & \textbf{Nature} & \textbf{Traduction} \\ \hline
尊敬的 & zūnjìng de & adj. + part. & Honorable, respecté (appel) \\ \hline
经理 & jīnglǐ & n. & Directeur, gérant, manager \\ \hline
申请 & shēnqǐng & v. / n. & Solliciter, demander / une demande \\ \hline
实习 & shíxí & v. / n. & Faire un stage / un stage \\ \hline
暑期 & shǔqī & n. & Période estivale, vacances d'été \\ \hline
贵公司 & guì gōngsī & n. phr. & Votre entreprise (honorifique) \\ \hline
考虑 & kǎolǜ & v. / n. & Considérer, examiner / considération \\ \hline
此致敬礼 & cǐ zhì jìnglǐ & expr. figée & Veuillez agréer mes salutations respectueuses \\ \hline
封 & fēng & class. & Classificateur pour lettres, enveloppes \\ \hline
寄 & jì & v. & Envoyer (par la poste) \\ \hline
回信 & huíxìn & v. / n. & Répondre à une lettre / lettre de réponse \\ \hline
媒体 & méitǐ & n. & Médias \\ \hline
比较 & bǐjiào & adv. & Assez, relativement \\ \hline
已经 & yǐjīng & adv. & Déjà (marque l'accompli / durée) \\ \hline
希望 & xīwàng & v. / n. & Espérer, souhaiter / espoir \\ \hline
能 & néng & v. aux. & Pouvoir (capacité, possibilité) \\ \hline
\end{tabularx}
\end{center}

\courssub{Atelier Caractères : Radicaux, Ordre des traits \& Mémorisation}

\begin{propriete}[Analyse de 8 caractères clés de la leçon]
\begin{center}
\begin{tabularx}{\linewidth}{|X|X|X|X|X|}
\hline
\rowcolor{popgreenL} \textbf{Car.} & \textbf{Pinyin} & \textbf{Clé (Radical)} & \textbf{Nb traits} & \textbf{Aide-mémoire / Ordre principal} \\ \hline
信 & xìn & \zh{亻} (rén, homme) & 9 & Homme (亻) + Parole (言) = La parole de l'homme = confiance / lettre. Ordre : 亻(2) → 言(7). \\ \hline
请 & qǐng & \zh{讠} (yán, parole) & 10 & Parole (讠) + 青 (qīng, phonétique) = Parler pour demander. Ordre : 讠(2) → 青(8). \\ \hline
谢 & xiè & \zh{讠} (yán, parole) & 12 & Parole (讠) + 射 (shè, tirer) = Paroles qui partent comme une flèche = remercier/s'excuser. Ordre : 讠(2) → 身(7) → 寸(3). \\ \hline
申 & shēn & \zh{田} (tián, champ) & 5 & 田 + 丨 = Étendre une demande sur le champ. Ordre : 田(5) → 丨(1) — \emph{attention : 5 traits total, 田 est 5 traits, mais ici 申 s'écrit en 5 traits : 丨 一 一 一 丨 (vertical, 3 horizontales, verticale finale)}. \\ \hline
实 & shí & \zh{宀} (mián, toit) & 8 & Toit (宀) + 贞 (zhēn, vrai) = La vérité sous le toit = réel, solide. Ordre : 宀(3) → 贞(5). \\ \hline
习 & xí & \zh{羽} (yǔ, plume) & 11 & Plumes (羽) + 白 (blanc) = Répéter jusqu'à blanchir les plumes = s'exercer. Ordre : 羽(6) → 白(5). \\ \hline
封 & fēng & \zh{寸} (cùn, pouce) & 9 & 寸 (mesure) + 丰 (fēng, abondant) = Mesure abondante = sceller une lettre. Ordre : 丰(4) → 寸(3) — \emph{structure haut-bas}. \\ \hline
贵 & guì & \zh{贝} (bèi, cauri) & 12 & Cauris (贝) + 朱 (zhū, vermillon) + 贝 = Double richesse = honorable, cher. Ordre : 贝(4) → 朱(6) → 贝(4) — \emph{structure haut-milieu-bas}. \\ \hline
\end{tabularx}
\end{center}
\end{propriete}

\begin{attention}[Pièges graphiques : Paires de caractères à ne pas confondre]
\begin{itemize}
\item \zh{请} (qǐng, demander, clé \zh{讠}) $\neq$ \zh{情} (qíng, sentiment, clé \zh{忄}) $\neq$ \zh{清} (qīng, clair, clé \zh{氵}). \textbf{Astuce} : Pour « demander », on utilise la \textbf{parole} (\zh{讠}).
\item \zh{信} (xìn, lettre, clé \zh{亻}) $\neq$ \zh{住} (zhù, habiter, clé \zh{亻}). \textbf{Astuce} : \zh{信} a \zh{言} (parole) à droite ; \zh{住} a \zh{主} (maître) à droite.
\item \zh{实} (shí, réel) $\neq$ \zh{寔} (forme traditionnelle de 实) — en simplifié, seul \zh{实} existe.
\item \zh{申} (shēn, 5 traits) $\neq$ \zh{由} (yóu, 5 traits, raison) $\neq$ \zh{甲} (jiǎ, 5 traits, carapace). Différence : trait vertical final pour \zh{申}.
\end{itemize}
\end{attention}

\courssub{Architecture rhétorique : Connecteurs \& Structures types du corps de lettre}

\begin{propriete}[Connecteurs logiques (连词) pour lettre formelle]
\begin{center}
\begin{tabularx}{\linewidth}{|X|X|X|}
\hline
\rowcolor{poppurpleL} \textbf{Fonction} & \textbf{Connecteur chinois} & \textbf{Exemple en contexte} \\ \hline
Addition / Aussi & \zh{也} (yě), \zh{而且} (érqiě) & \zh{我会说汉语，而且会写汉字。} \\ \hline
Cause / Parce que & \zh{因为} (yīnwèi) ... \zh{所以} (suǒyǐ) & \zh{因为我喜欢媒体，所以我申请实习。} \\ \hline
But / Afin que & \zh{为了} (wèile) + [verbe] & \zh{为了学习技术，我来贵公司。} \\ \hline
Condition / Si & \zh{如果} (rúguǒ) ... \zh{就} (jiù) & \zh{如果有机会，我就很高兴。} \\ \hline
Concession / Bien que & \zh{虽然} (suīrán) ... \zh{但是} (dànshì) & \zh{虽然我年轻，但是我努力。} \\ \hline
Séquence / Ensuite & \zh{然后} (ránhòu), \zh{此外} (cǐwài) & \zh{首先我学习，然后我工作。} \\ \hline
\end{tabularx}
\end{center}
\end{propriete}

\begin{propriete}[5 Structures types du corps de lettre (Schémas syntaxiques)]
\begin{enumerate}
\item \textbf{Présentation} : \zh{我是} + [Lieu/École] + \zh{的学生}，\zh{我叫} + [Prénom]. \\
      \textit{Ex.} \zh{我是杜阿拉第一中学高二年级的学生，我叫玛丽。}
\item \textbf{Motif (Objet)} : \zh{我写信是为了} + [Verbe/Action] + [Objet]. \\
      \textit{Ex.} \zh{我写信是为了申请贵公司的暑期实习。}
\item \textbf{Durée d'apprentissage / Expérience} : \zh{我} + [Verbe] + [Objet] + \zh{已经} + [Durée] + \zh{了}. \\
      \textit{Ex.} \zh{我学习汉语已经两年了。} (\textbf{Attention} : Durée APRÈS le verbe/objet).
\item \textbf{Compétences / Atouts} : \zh{我} + \zh{会} / \zh{能} + [Verbe], \zh{也} + [Verbe/Adj]. \\
      \textit{Ex.} \zh{我会说法语和汉语，也很擅长电脑。}
\item \textbf{Souhait final / Appel à l'action} : \zh{我希望能} + [Verbe] + [Lieu/Contexte]. \\
      \textit{Ex.} \zh{我希望能在贵公司学习。}
\end{enumerate}
\end{propriete}

\courssub{Compréhension de texte \& Collocations : Du passif à l'actif}

\begin{methode}[Décoder une lettre officielle chinoise — 5 étapes]
\begin{enumerate}
\item \textbf{Repérer l'en-tête} : L'appel poli \zh{尊敬的…} (zūnjìng de...) suivi du titre/fonction, puis \zh{您好！}.
\item \textbf{Identifier l'objet} : La structure \zh{我写信是为了…} annonce le motif principal.
\item \textbf{Extraire le profil du candidat} : École (\zh{...的学生}), Nom (\zh{我叫...}), Durée d'étude (\zh{已经...了}), Compétences (\zh{会/能/喜欢}).
\item \textbf{Repérer la clôture codifiée} : \zh{谢谢您的考虑！} (Remerciement) → \zh{此致敬礼！} (Formule de politesse finale) → Signature (Prénom) + Date (\zh{月日}) à droite.
\item \textbf{Stratégie questions} : Une question = un bloc. Répondre par une phrase complète reprenant les mots-clés du texte.
\end{enumerate}
\end{methode}

\begin{exoresolu}[Compréhension détaillée : Analyse de la lettre de Li Hua]
Lis le texte support (Mise en situation) et réponds en chinois (caractères + pinyin).

\textbf{Questions :}
1) 李华在哪里上学？ 
2) 他学习汉语多长时间了？ 
3) 他有什么特长 (spécialités / points forts) ? 
4) Quelle est la formule de politesse finale utilisée ?

\tcblower
\textbf{Solution détaillée :}
\begin{enumerate}
\item \textbf{Q1 (Localisation)} : \zh{李华在雅温得第二中学上学。} (Lǐ Huá zài Yǎwēndé Dì-èr Zhōngxué shàngxué.) — Repéré dans le bloc présentation : \zh{雅温得第二中学高二年级的学生}.
\item \textbf{Q2 (Durée)} : \zh{他学习汉语已经两年了。} (Tā xuéxí Hànyǔ yǐjīng liǎng nián le.) — Repéré : \zh{学习汉语已经两年了}. \textbf{Note} : \zh{两} (liǎng) et non \zh{二} devant classifieur/année.
\item \textbf{Q3 (Compétences)} : \zh{他会说汉语，也很喜欢电脑和媒体工作。} (Tā huì shuō Hànyǔ, yě hěn xǐhuan diànnǎo hé méitǐ gōngzuò.) — Synthèse du bloc arguments : \zh{汉语说得比较好} + \zh{喜欢电脑和媒体工作}.
\item \textbf{Q4 (Clôture)} : \zh{此致敬礼！} (Cǐ zhì jìnglǐ !) — Formule figée obligatoire en bas à gauche.
\end{enumerate}
\textbf{Conseil méthodologique} : Ne recopie pas le texte entier. Construis ta réponse avec le sujet (他/李华) + le verbe/structure cible.
\end{exoresolu}

\begin{exoresolu}[Collocations \& Classificateurs : Complétion guidée]
Complète les phrases avec les mots du glossaire : \zh{封、贵、您、申请、实习、寄、回信、考虑} .

1) 我想写一\_\_\_信给经理。 
2) 我\_\_\_\_\_\_贵公司的暑期\_\_\_。 
3) 谢谢\_\_\_的\_\_\_！ 
4) 请\_\_\_信给我。 
5) \_\_\_公司很有名。

\tcblower
\textbf{Solution détaillée :}
\begin{enumerate}
\item \zh{我想写一\cle{封}信给经理。} (Wǒ xiǎng xiě \textbf{yì fēng} xìn gěi jīnglǐ.) — Classificateur obligatoire pour « lettre ».
\item \zh{我\cle{申请}\cle{贵}公司的暑期\cle{实习}。} (Wǒ \textbf{shēnqǐng guì gōngsī de shǔqī shíxí}.) — Verbe \zh{申请} + \zh{贵} (honorifique) + Nom \zh{实习}.
\item \zh{谢谢\cle{您}的\cle{考虑}！} (Xièxie \textbf{nín de kǎolǜ}!) — \zh{您} (poli) + \zh{考虑} (nom « considération/attention »).
\item \zh{请\cle{寄}\cle{回信}给我。} (Qǐng \textbf{jì huíxìn} gěi wǒ.) — Verbe \zh{寄} (envoyer) + Objet \zh{回信} (lettre de réponse).
\item \zh{\cle{贵}公司很有名。} (\textbf{Guì gōngsī} hěn yǒumíng.) — Préfixe honorifique \zh{贵} devant l'institution du destinataire.
\end{enumerate}
\textbf{Règle d'or} : Mémorise les \textbf{blocs indissociables} : \zh{写一封信}、\zh{申请实习}、\zh{贵公司}、\zh{谢谢您}、\zh{寄回信}.
\end{exoresolu}

\courssub{Production guidée : Atelier d'écriture « Ma première lettre officielle »}

\begin{methode}[Rédiger une lettre formelle en 6 étapes (Check-list Bac)]
\begin{enumerate}
\item \textbf{En-tête (Gauche)} : \zh{尊敬的} + [Titre/Nom] + \zh{：} \quad \zh{您好！}
\item \textbf{Identité (Qui ?)} : \zh{我是} + [École/Ville] + \zh{高二年级的学生，我叫} + [Prénom] + \zh{。}
\item \textbf{Objet (Pourquoi ?)} : \zh{我写信是为了} + \zh{申请} + \zh{贵公司/贵校} + [Poste/Stage] + \zh{。}
\item \textbf{Arguments (Preuves)} : 2 phrases min. \zh{我学习汉语已经[Durée]了，我会[Compétence 1]，也[Compétence 2]。} + \zh{因为}[Raison] + \zh{所以}[Motivation].
\item \textbf{Clôture (Formules fixes)} : \zh{谢谢您的考虑！此致敬礼！} (Sur 2 lignes).
\item \textbf{Signature (Droite)} : [Prénom] \quad [Mois]月[Jour]日 (ex. \zh{七月二日}).
\end{enumerate}
\end{methode}

\begin{exoresolu}[Composition guidée : Ta lettre pour un stage à Douala]
\textbf{Consigne} : Rédige une lettre formelle (6-8 phrases) à \zh{张经理} (Directeur Zhang) d'une entreprise de médias à Douala (\zh{杜阿拉媒体公司}). Tu es \zh{保罗} (Bǎoluó), élève de \zh{杜阿拉第一中学} (Lycée n°1 de Douala), \zh{高二年级}. Tu apprends le chinois \zh{已经三年了}, tu \zh{会说} chinois, tu \zh{喜欢} journalisme (\zh{新闻}) et \zh{电脑}. Demande un \zh{暑期实习}.

\tcblower
\textbf{Modèle de rédaction (Version élève attendue) :}

\zh{尊敬的张经理：您好！} \\
\zh{我是杜阿拉第一中学高二年级的学生，我叫保罗。} \\
\zh{我写信是为了申请贵公司的暑期实习。} \\
\zh{我学习汉语已经三年了，我会说法语和汉语，也很喜欢新闻和电脑。} \\
\zh{因为我对媒体工作很感兴趣，所以我希望能在贵公司学习。} \\
\zh{谢谢您的考虑！此致敬礼！} \\
\zh{保罗，七月五日}

\textbf{Analyse linguistique du modèle :}
\begin{itemize}
\item \textbf{Registre} : \zh{尊敬的张经理} + \zh{您} + \zh{贵公司} = 3 marqueurs de politesse (Cible Bac : 3/3).
\item \textbf{Grammaire} : \zh{已经三年了} (durée post-verbale) ; \zh{会说法语和汉语} (capacité) ; \zh{因为…所以…} (cause-conséquence correcte) ; \zh{在贵公司学习} (lieu avant verbe).
\item \textbf{Vocabulaire précis} : \zh{新闻} (journalisme/info), \zh{感兴趣} (intéressé par), \zh{暑期实习} (stage d'été).
\item \textbf{Mise en page} : \zh{此致} fin de ligne, \zh{敬礼} début ligne suivante ; Signature/Date à droite.
\end{itemize}
\end{exoresolu}

\courssub{Extension Bac \& Évaluation formative (Complément examen)}

\begin{exercice}[Type Bac — Compréhension écrite : Lettre de réponse]
Lis la lettre de réponse du Directeur Zhang à Paul, puis réponds aux questions en français.

\begin{textebox}[Réponse de l'entreprise — 杜阿拉媒体公司]
保罗同学：您好！\\
Bǎoluó tóngxué : Nín hǎo !\\
Cher étudiant Paul, bonjour !\\
收到您的来信，我们很高兴。\\
Shōudào nín de láixìn, wǒmen hěn gāoxìng.\\
Nous avons bien reçu votre lettre, nous sommes très contents.\\
贵校的汉语教学质量很高，我们接受您的实习申请。\\
Guì xiào de Hànyǔ jiàoxué zhìliàng hěn gāo, wǒmen jiēshòu nín de shíxí shēnqǐng.\\
La qualité de l'enseignement du chinois dans votre école est élevée, nous acceptons votre demande de stage.\\
请您七月十日来公司面试。\\
Qǐng nín qī yuè shí rì lái gōngsī miànshì.\\
Merci de venir passer un entretien à l'entreprise le 10 juillet.\\
此致敬礼！\\
Cǐ zhì jìnglǐ !\\
杜阿拉媒体公司，七月一日\\
Dù'ālā Méitǐ Gōngsī, qī yuè yī rì\\
Entreprise Média de Douala, le 1er juillet
\end{textebox}

\textbf{Questions :}
1) Qui écrit cette lettre et à qui ? 
2) Quelle est la décision prise par l'entreprise ? 
3) Quelle action est demandée à Paul et pour quand ? 
4) Relevez deux marqueurs de politesse formelle utilisés par l'expéditeur.
\end{exercice}

\begin{corrige}[Exercice Compréhension écrite]
\begin{enumerate}
\item \textbf{Expéditeur} : \zh{杜阿拉媒体公司} (Entreprise Média de Douala) — signature en bas. \textbf{Destinataire} : \zh{保罗同学} (Étudiant Paul) — appel en tête. \textit{Note} : \zh{同学} (tóngxué) est utilisé par un supérieur/entreprise pour s'adresser poliment à un élève/étudiant.
\item \textbf{Décision} : \zh{我们接受您的实习申请} (Nous acceptons votre demande de stage).
\item \textbf{Action / Date} : \zh{请您七月十日来公司面试} (Veuillez venir à l'entreprise pour un entretien le 10 juillet). \zh{面试} (miànshì) = entretien d'embauche / stage.
\item \textbf{Marqueurs de politesse} : 
      \begin{itemize}
      \item \zh{您} (nín) — pronom « vous » poli (utilisé 2 fois).
      \item \zh{贵校} (guì xiào) — « votre honorable école » (préfixe \zh{贵}).
      \item \zh{请} (qǐng) — « veuillez / merci de » (impératif poli) devant \zh{您七月十日来...}.
      \item \zh{此致敬礼} — formule de clôture.
      \end{itemize}
\end{enumerate}
\end{corrige}

\begin{exercice}[Type Bac — Traduction Thème (Français → Chinois)]
Traduis en chinois (caractères + pinyin) les phrases suivantes :
1) Je vous écris pour demander un stage d'été dans votre entreprise.
2) J'étudie le chinois depuis trois ans, je peux parler et écrire.
3) Parce que j'aime les médias, j'espère pouvoir travailler chez vous.
4) Merci de votre attention. Veuillez agréer mes salutations respectueuses.
5) Paul, le 12 août.
\end{exercice}

\begin{corrige}[Exercice Thème]
\begin{enumerate}
\item \zh{我写信是为了申请贵公司的暑期实习。} (Wǒ xiěxìn shì wèile shēnqǐng guì gōngsī de shǔqī shíxí.)
\item \zh{我学习汉语已经三年了，我会说也会写。} (Wǒ xuéxí Hànyǔ yǐjīng sān nián le, wǒ huì shuō yě huì xiě.) \textit{Note : \zh{三年} (sān nián) pas de \zh{两} car 3.}
\item \zh{因为我喜欢媒体，所以我希望能在贵公司工作。} (Yīnwèi wǒ xǐhuan méitǐ, suǒyǐ wǒ xīwàng néng zài guì gōngsī gōngzuò.)
\item \zh{谢谢您的考虑！此致敬礼！} (Xièxie nín de kǎolǜ ! Cǐ zhì jìnglǐ !)
\item \zh{保罗，八月十二日} (Bǎoluó, bā yuè shí'èr rì.)
\end{enumerate}
\end{corrige}

\begin{exercice}[Type Bac — Correction d'erreurs (Relevés d'examinateur)]
Corrige les 5 erreurs dans cet extrait de lettre d'élève. Réécris la phrase correcte et justifie la correction grammaticale ou lexicale.

\textbf{Extrait fautif :}
\zh{尊敬的经理：你好！我是学生，我叫马力。我写信为了申请你公司的实习。我学习中文已经二年了，我可以说中文。希望在你公司工作。谢谢你们的考虑。此致敬礼 马力，八月十日}
\end{exercice}

\begin{corrige}[Exercice Correction d'erreurs]
\begin{enumerate}
\item \zh{尊敬的经理：\cle{您}好！} — \textbf{Erreur} : \zh{你} (familier) → \textbf{Correction} : \zh{您} (poli) obligatoire après \zh{尊敬的}.
\item \zh{我写信\cle{是为了}申请\cle{贵}公司的实习。} — \textbf{Erreur 1} : Manque \zh{是} dans la structure \zh{写信是为了}. \textbf{Erreur 2} : \zh{你公司} → \textbf{Correction} : \zh{贵公司} (honorifique).
\item \zh{我学习中文已经\cle{两}年了。} — \textbf{Erreur} : \zh{二年} → \textbf{Correction} : \zh{两年} (liǎng nián) devant classifieur/durée.
\item \zh{我\cle{会说}中文。} — \textbf{Erreur} : \zh{可以说} (peut parler - possibilité) → \textbf{Correction} : \zh{会说} (sait parler - compétence acquise) ou \zh{能说}. \textit{Préférence Bac} : \zh{会说} pour compétence.
\item \zh{希望\cle{能}在\cle{贵}公司工作。} — \textbf{Erreur 1} : Sujet \zh{我} omis (recommandé en début de phrase). \textbf{Erreur 2} : \zh{你公司} → \textbf{Correction} : \zh{贵公司}. \textbf{Erreur 3} : Manque \zh{能} (souhait de possibilité) → \zh{我希望能在贵公司工作}.
\item \zh{谢谢\cle{您}的考虑。} — \textbf{Erreur} : \zh{你们} (pluriel) → \textbf{Correction} : \zh{您} (singulier poli, s'adresse au Directeur).
\item \zh{此致敬礼！\par 马力，八月十日} — \textbf{Mise en page} : \zh{此致敬礼} sur deux lignes ; Signature/Date à droite.
\end{enumerate}
\textbf{Version corrigée complète :}
\zh{尊敬的经理：您好！我是学生，我叫马力。我写信是为了申请贵公司的实习。我学习中文已经两年了，我会说法语和中文。我希望能在贵公司工作。谢谢您的考虑！此致\par 敬礼！\hfill 马力，八月十日}
\end{corrige}

\begin{attention}[Les 5 erreurs fatales au Bac (Synthèse finale)]
\begin{enumerate}
\item \textbf{Tons \& Prononciation} : \zh{申请} (\cle{shēnqǐng} 1-3), \zh{实习} (\cle{shíxí} 2-2), \zh{经理} (\cle{jīnglǐ} 1-3), \zh{贵} (\cle{guì} 4). Un ton faux = mot incompréhensible = 0 pt au mot.
\item \textbf{Caractères « jumeaux »} : \zh{请/情/清}, \zh{信/住}, \zh{申/由/甲}, \zh{实/寔}. \textbf{Règle} : Identifie la \textbf{clé} (radical) avant d'écrire. \zh{讠}= parole (demander, remercier) ; \zh{忄}= cœur (sentiment) ; \zh{氵}= eau (clair).
\item \textbf{Ordre des mots (Syntaxique)} : 
      \begin{itemize}
      \item Lieu \zh{在贵公司} AVANT verbe \zh{工作/学习} (Jamais \zh{工作在贵公司}).
      \item Durée \zh{两年} APRÈS verbe/objet \zh{学习中文两年了} (Jamais \zh{两年学习中文}).
      \item Structure \zh{写信是为了} + V (Pas \zh{写信为了}).
      \end{itemize}
\item \textbf{Registre de politesse (Sociolinguistique)} : 
      \begin{itemize}
      \item \zh{您} (pas \zh{你}), \zh{贵公司/贵校} (pas \zh{你公司}), \zh{尊敬的} + Titre (pas \zh{亲爱的} qui est affectif).
      \item \zh{此致敬礼} sur 2 lignes alignées à gauche.
      \end{itemize}
\item \textbf{Classificateurs \& Nombres} : 
      \begin{itemize}
      \item 1 lettre = \zh{一\cle{封}信} (pas \zh{个}, \zh{张}).
      \item 2 ans = \zh{\cle{两}年} (pas \zh{二年}).
      \item Étages/Numéro = \zh{二} (ex. \zh{二楼}, \zh{二号}).
      \end{itemize}
\end{enumerate}
\end{attention}

\begin{savaistu}[Le saviez-vous ? L'évolution de la correspondance chinoise]
Autrefois, les lettres officielles en Chine impériale utilisaient le \zh{文言文} (wényánwén, chinois classique) avec des formules extrêmement codifiées selon la hiérarchie (ex. \zh{启者} qǐ zhě pour ouvrir, \zh{顿首} dùnshǒu pour saluer). Aujourd'hui, le \zh{白话文} (báihuàwén, chinois moderne) a simplifié les formules (\zh{尊敬的...您好}, \zh{此致敬礼}), mais la structure rigide (En-tête → Corps → Clôture → Signature) reste une marque de professionnalisme. Au Cameroun, l'administration utilise le français avec des codes similaires (Objet, Vu, Arrête). La maîtrise de ces codes dans les deux langues est un atout majeur pour l'employabilité dans la zone CEMAC et avec les partenaires chinois.
\end{savaistu}

\newpage

\coursec{Exercices}

\courssub{Je vérifie mes connaissances}

\begin{exercice}[QCM : la lettre formelle]
Pour chaque question, entoure la bonne réponse.
\begin{enumerate}
\item Pour commencer une lettre formelle à un directeur que l'on ne connaît pas, on écrit :
\begin{enumerate}
\item \zh{你好！} \item \zh{尊敬的经理，您好！} \item \zh{喂，经理！} \item \zh{老朋友，你好！}
\end{enumerate}
\item La formule de politesse finale correcte est :
\begin{enumerate}
\item \zh{再见！} \item \zh{谢谢你的信。} \item \zh{此致敬礼！} \item \zh{晚安！}
\end{enumerate}
\item \zh{申请} (shēnqǐng) signifie :
\begin{enumerate}
\item inviter \item demander officiellement \item refuser \item téléphoner
\end{enumerate}
\item Sur une enveloppe chinoise, le caractère placé après le nom du destinataire est :
\begin{enumerate}
\item \zh{寄} \item \zh{收} \item \zh{写} \item \zh{信}
\end{enumerate}
\end{enumerate}
\end{exercice}

\begin{exercice}[Vrai ou faux justifié]
Indique si chaque affirmation est vraie ou fausse, puis justifie ta réponse en une phrase en français.
\begin{enumerate}
\item Dans une lettre formelle chinoise, on peut tutoyer le destinataire avec \zh{你} (nǐ).
\item La date s'écrit en bas à droite de la lettre, avec la signature.
\item \zh{此致敬礼} (cǐ zhì jìnglǐ) est une formule de salutation d'ouverture.
\item En chinois, l'adresse sur l'enveloppe va du plus grand (le pays) au plus petit (le numéro de rue).
\end{enumerate}
\end{exercice}

\begin{exercice}[Vocabulaire : complète le tableau]
Recopie et complète le tableau suivant (nature et traduction française) :

\begin{center}
\begin{tabularx}{\linewidth}{|X|X|X|X|}
\hline
\rowcolor{popblueL} Caractères & Pinyin & Nature & Traduction \\
\hline
\zh{尊敬} & zūnjìng & \ldots & \ldots \\
\hline
\zh{实习} & shíxí & \ldots & \ldots \\
\hline
\zh{机会} & jīhuì & \ldots & \ldots \\
\hline
\zh{感谢} & gǎnxiè & \ldots & \ldots \\
\hline
\zh{附件} & fùjiàn & \ldots & \ldots \\
\hline
\zh{回复} & huífù & \ldots & \ldots \\
\hline
\end{tabularx}
\end{center}
\end{exercice}

\begin{exercice}[Associe le pinyin aux caractères]
Relie chaque caractère ou mot à son pinyin :
\begin{enumerate}
\item \zh{信} \hspace{1cm} a. xiě
\item \zh{写} \hspace{1cm} b. jì
\item \zh{寄} \hspace{1cm} c. xìn
\item \zh{收} \hspace{1cm} d. zhí
\item \zh{职} \hspace{1cm} e. shōu
\end{enumerate}
\end{exercice}

\courssub{Je m'entraîne}

\begin{exercice}[Traduction : du français au chinois]
Traduis les phrases suivantes en chinois (caractères et pinyin) :
\begin{enumerate}
\item Je vous écris pour demander un stage dans votre entreprise.
\item Je vous remercie de votre attention.
\item Veuillez trouver ci-joint mon CV.
\item J'attends votre réponse avec impatience.
\end{enumerate}
\end{exercice}

\begin{exercice}[Transformation de registre]
Les phrases suivantes sont trop familières pour une lettre formelle. Réécris-les en style épistolaire formel, en chinois avec pinyin.
\begin{enumerate}
\item \zh{我要一个实习。} (Wǒ yào yí ge shíxí.) — « Je veux un stage. »
\item \zh{你给我打电话吧。} (Nǐ gěi wǒ dǎ diànhuà ba.) — « Appelle-moi. »
\item \zh{我很喜欢你们公司。} (Wǒ hěn xǐhuan nǐmen gōngsī.) — « J'aime bien votre boîte. »
\item \zh{快点回信！} (Kuài diǎn huíxìn!) — « Réponds vite ! »
\item \zh{我中文说得不错。} (Wǒ Zhōngwén shuō de búcuò.) — « Je me débrouille en chinois. »
\end{enumerate}
\end{exercice}

\begin{exercice}[Texte à trous : la lettre de demande]
Complète la lettre suivante avec les mots de la liste : \zh{尊敬、申请、机会、感谢、此致、回复}

\begin{textebox}[Lettre à compléter]
\zh{＿＿＿的经理，您好！}\\
\zh{我是雅温得中学的学生。我写这封信是为了＿＿＿贵公司的一个实习＿＿＿。}\\
\zh{我学习中文已经三年了，我非常＿＿＿您给我这个机会。}\\
\zh{期待您的＿＿＿。}\\
\zh{＿＿＿敬礼！}
\end{textebox}
\end{exercice}

\begin{exercice}[Remise en ordre et caractères intrus]
\textbf{A.} Remets les mots dans l'ordre pour former une phrase correcte de lettre formelle :
\begin{enumerate}
\item \zh{是 / 我 / 为了 / 写信 / 申请 / 实习}
\item \zh{您 / 感谢 / 非常 / 的 / 关注}
\end{enumerate}
\textbf{B.} Dans chaque série de 4 caractères, trois partagent le même radical ; trouve l'intrus et justifie :
\begin{enumerate}
\item \zh{语 — 说 — 话 — 手}
\item \zh{信 — 你 — 他 — 心}
\item \zh{想 — 感 — 谢 — 忘}
\end{enumerate}
\end{exercice}

\courssub{Je me prépare au Bac}

\begin{exercice}[Type Bac — Compréhension de texte]
Lis la lettre suivante, puis réponds aux questions.

\begin{textebox}[Une lettre de demande de stage]
\zh{尊敬的经理，您好！}\\
\zh{Zūnjìng de jīnglǐ, nín hǎo!}\\
\zh{我是杜阿拉中学一年级的学生，我叫玛丽。我学习中文已经两年了，成绩很好。我写这封信是为了申请贵公司今年七月的一个实习机会。}\\
\zh{Wǒ shì Dùālā zhōngxué yì niánjí de xuésheng, wǒ jiào Mǎlì. Wǒ xuéxí Zhōngwén yǐjīng liǎng nián le, chéngjì hěn hǎo. Wǒ xiě zhè fēng xìn shì wèile shēnqǐng guì gōngsī jīnnián qīyuè de yí ge shíxí jīhuì.}\\
\zh{我会说中文、法文和英文，也很喜欢电脑。附件里有我的简历。}\\
\zh{Wǒ huì shuō Zhōngwén, Fǎwén hé Yīngwén, yě hěn xǐhuan diànnǎo. Fùjiàn lǐ yǒu wǒ de jiǎnlì.}\\
\zh{感谢您的时间，期待您的回复。}\\
\zh{Gǎnxiè nín de shíjiān, qīdài nín de huífù.}\\
\zh{此致敬礼！}\\
\zh{Cǐ zhì jìnglǐ!}\\
\zh{学生：玛丽　二〇二四年五月十日}\\
\zh{Xuésheng: Mǎlì　èr líng èr sì nián wǔ yuè shí rì}\\
\emph{Traduction : Monsieur le Directeur, bonjour ! Je suis une élève de Première du lycée de Douala, je m'appelle Marie. J'apprends le chinois depuis deux ans et j'ai de bons résultats. Je vous écris pour demander un stage dans votre entreprise au mois de juillet de cette année. Je parle chinois, français et anglais, et j'aime beaucoup l'informatique. Vous trouverez mon CV en pièce jointe. Je vous remercie de votre temps et j'attends votre réponse. Veuillez agréer mes salutations distinguées ! L'élève : Marie, le 10 mai 2024.}
\end{textebox}

\textbf{Questions de compréhension}
\begin{enumerate}
\item QCM : Qui est l'auteur de la lettre ?\\
a. un directeur \quad b. une élève de Première \quad c. un professeur \quad d. un journaliste
\item QCM : Pour quel mois Marie demande-t-elle un stage ?\\
a. mai \quad b. juin \quad c. juillet \quad d. août
\item Question ouverte (en français) : Quelles langues Marie parle-t-elle ?
\item Question ouverte (en français) : Quel document se trouve dans la pièce jointe ?
\item Question ouverte (en chinois) : \zh{玛丽学习中文多长时间了？}
\end{enumerate}
\textbf{Questions de langue}
\begin{enumerate}
\item Releve dans le texte la formule d'appel et la formule finale de politesse.
\item Explique l'emploi du caractère \zh{贵} dans \zh{贵公司}.
\end{enumerate}
\end{exercice}

\begin{exercice}[Type Bac — Thème et version]
\textbf{A. Thème.} Traduis en chinois (caractères et pinyin) :
\begin{enumerate}
\item Je vous écris pour demander un document officiel.
\item Je parle trois langues : le chinois, le français et l'anglais.
\item Je vous remercie de votre réponse.
\end{enumerate}
\textbf{B. Version.} Traduis en français :
\begin{enumerate}
\item \zh{我写这封信是为了申请一个实习机会。} (Wǒ xiě zhè fēng xìn shì wèile shēnqǐng yí ge shíxí jīhuì.)
\item \zh{附件里有我的简历，期待您的回复。} (Fùjiàn lǐ yǒu wǒ de jiǎnlì, qīdài nín de huífù.)
\end{enumerate}
\end{exercice}

\begin{exercice}[Type Bac — Expression écrite : lettre formelle]
\textbf{Sujet.} Tu es élève de Première A au lycée de Yaoundé. Tu écris une lettre au directeur d'une entreprise de télécommunications chinoise installée à Douala pour demander un stage de vacances. Rédige une lettre formelle complète en chinois (100 à 120 caractères) comprenant : l'appel honorifique, la présentation de toi, l'objet de la demande, tes qualités (langues, informatique), la formule finale de politesse, la signature et la date.

\textbf{Barème indicatif.} Format et mise en page de la lettre : 5 points ; vocabulaire honorifique et formules de politesse : 5 points ; syntaxe et connecteurs : 5 points ; exactitude des caractères : 5 points.
\end{exercice}

\newpage

\coursec{Corrigés des exercices}

\begin{corrige}[Exercice 1 — QCM : la lettre formelle]
\textbf{1. b} — \zh{尊敬的经理，您好！} (Zūnjìng de jīnglǐ, nín hǎo!) : c'est l'appel honorifique standard d'une lettre formelle adressée à un supérieur que l'on ne connaît pas. \zh{你好！} est trop neutre, \zh{喂！} sert à répondre au téléphone, et \zh{老朋友} convient à un ami proche.

\textbf{2. c} — \zh{此致敬礼！} (Cǐ zhì jìnglǐ!) : formule finale de politesse épistolaire, équivalent de « Veuillez agréer mes salutations distinguées ». \zh{再见！} signifie « au revoir » (oral), \zh{谢谢你的信。} n'est pas une formule finale, \zh{晚安！} signifie « bonne nuit ».

\textbf{3. b} — \zh{申请} (shēnqǐng) signifie « demander officiellement, faire une demande (de stage, de visa, de poste) ».

\textbf{4. b} — \zh{收} (shōu, « recevoir ») : sur une enveloppe chinoise, on écrit le nom du destinataire suivi de \zh{收} (« à remettre à »). \zh{寄} (jì, « envoyer ») accompagne le nom de l'expéditeur : \zh{寄信人} « expéditeur ».
\end{corrige}

\begin{corrige}[Exercice 2 — Vrai ou faux justifié]
\textbf{1. FAUX.} Dans une lettre formelle chinoise, on vouvoie le destinataire avec \zh{您} (nín), forme polie de \zh{你} (nǐ). Le tutoiement serait une faute de registre grave devant un directeur.

\textbf{2. VRAI.} La date (par exemple \zh{二〇二四年五月十日}) et la signature se placent en bas à droite de la lettre, sous la formule finale \zh{此致敬礼}.

\textbf{3. FAUX.} \zh{此致敬礼} (cǐ zhì jìnglǐ) est la formule de politesse \emph{finale}, placée en toute fin de lettre ; la salutation d'ouverture est \zh{尊敬的……，您好！}.

\textbf{4. VRAI.} Contrairement au français, l'adresse chinoise procède du plus grand au plus petit : pays → ville → quartier → rue → numéro. C'est le même ordre que pour les dates : année → mois → jour.
\end{corrige}

\begin{corrige}[Exercice 3 — Vocabulaire : complète le tableau]
\begin{center}
\begin{tabularx}{\linewidth}{|X|X|X|X|}
\hline
\rowcolor{popblueL} Caractères & Pinyin & Nature & Traduction \\
\hline
\zh{尊敬} & zūnjìng & verbe / adjectif & respecter ; honorable \\
\hline
\zh{实习} & shíxí & nom / verbe & stage ; faire un stage \\
\hline
\zh{机会} & jīhuì & nom & occasion, opportunité \\
\hline
\zh{感谢} & gǎnxiè & verbe & remercier, être reconnaissant \\
\hline
\zh{附件} & fùjiàn & nom & pièce jointe \\
\hline
\zh{回复} & huífù & verbe / nom & répondre ; réponse \\
\hline
\end{tabularx}
\end{center}

\textbf{Points de vigilance.} \zh{机会} est toujours un nom ; on dit \zh{一个机会} avec le classificateur \zh{个}. \zh{感谢} est plus soutenu que \zh{谢谢} : il convient au style épistolaire. Attention au ton de \zh{实习} (shíxí, 2\textsuperscript{e} + 2\textsuperscript{e} ton), à ne pas confondre avec \zh{学习} (xuéxí, « étudier »).
\end{corrige}

\begin{corrige}[Exercice 4 — Associe le pinyin aux caractères]
\begin{enumerate}
\item \zh{信} → \textbf{c. xìn} (« lettre » ; radical de l'homme \zh{亻} + \zh{言} « parole »)
\item \zh{写} → \textbf{a. xiě} (« écrire » ; 3\textsuperscript{e} ton)
\item \zh{寄} → \textbf{b. jì} (« envoyer par la poste » ; radical toit \zh{宀})
\item \zh{收} → \textbf{e. shōu} (« recevoir »)
\item \zh{职} → \textbf{d. zhí} (« poste, fonction », comme dans \zh{职位} « poste »)
\end{enumerate}

\textbf{Point de vigilance.} Ne pas confondre \zh{信} (xìn, 4\textsuperscript{e} ton, « lettre ») et \zh{心} (xīn, 1\textsuperscript{er} ton, « cœur ») : homophones proches, sens très différents.
\end{corrige}

\begin{corrige}[Exercice 5 — Traduction : du français au chinois]
\begin{enumerate}
\item \zh{我给您写信是为了申请贵公司的一个实习。}\\
\emph{Wǒ gěi nín xiěxìn shì wèile shēnqǐng guì gōngsī de yí ge shíxí.}\\
Structure : \zh{给} + personne + \zh{写信} « écrire une lettre à quelqu'un » ; \zh{是为了} « c'est pour » introduit le but ; \zh{贵公司} « votre honorable entreprise ».
\item \zh{感谢您的关注。}\\
\emph{Gǎnxiè nín de guānzhù.}\\
\zh{感谢} + personne + \zh{的} + nom : « remercier quelqu'un de quelque chose ».
\item \zh{附件里有我的简历。}\\
\emph{Fùjiàn lǐ yǒu wǒ de jiǎnlì.}\\
Structure d'existence : lieu (\zh{附件里} « dans la pièce jointe ») + \zh{有} + chose. « Veuillez trouver » se rend naturellement par cette tournure.
\item \zh{期待您的回复。}\\
\emph{Qīdài nín de huífù.}\\
Formule épistolaire consacrée, sans sujet exprimé : « (j')attends votre réponse ». On peut ajouter \zh{我} : \zh{我期待您的回复。}
\end{enumerate}
\end{corrige}

\begin{corrige}[Exercice 6 — Transformation de registre]
\begin{enumerate}
\item \zh{我想申请一个实习机会。}\\
\emph{Wǒ xiǎng shēnqǐng yí ge shíxí jīhuì.}\\
On remplace \zh{要} (« vouloir », trop direct) par \zh{想} ou \zh{希望} et on emploie le verbe officiel \zh{申请}.
\item \zh{请您给我打电话。}\\
\emph{Qǐng nín gěi wǒ dǎ diànhuà.}\\
On ajoute \zh{请} et le vouvoiement \zh{您} ; la particule finale familière \zh{吧} disparaît.
\item \zh{我非常欣赏贵公司。}\\
\emph{Wǒ fēicháng xīnshǎng guì gōngsī.}\\
\zh{很喜欢} → \zh{非常欣赏} (« apprécier beaucoup », registre soutenu) ; \zh{你们公司} → \zh{贵公司} (honorifique).
\item \zh{期待您早日回复。}\\
\emph{Qīdài nín zǎorì huífù.}\\
L'impératif impatient disparaît ; \zh{早日} « au plus tôt » exprime poliment le souhait d'une réponse rapide.
\item \zh{我的中文水平不错，可以进行日常交流。}\\
\emph{Wǒ de Zhōngwén shuǐpíng búcuò, kěyǐ jìnxíng rìcháng jiāoliú.}\\
On objectivise : \zh{中文水平} « niveau de chinois », \zh{可以} « être capable de » ; on évite la première personne trop directe.
\end{enumerate}
\textbf{Point de vigilance.} En style formel : vouvoiement \zh{您}, honorifique \zh{贵}, verbes soutenus (\zh{申请}, \zh{感谢}, \zh{期待}), pas de particules familières (\zh{吧}, \zh{啊}), pas d'impératif sec.
\end{corrige}

\begin{corrige}[Exercice 7 — Texte à trous : la lettre de demande]
\begin{textebox}[Lettre complétée]
\zh{尊敬的经理，您好！}\\
\zh{Zūnjìng de jīnglǐ, nín hǎo!}\\
\zh{我是雅温得中学的学生。我写这封信是为了申请贵公司的一个实习机会。}\\
\zh{Wǒ shì Yǎwēndé zhōngxué de xuésheng. Wǒ xiě zhè fēng xìn shì wèile shēnqǐng guì gōngsī de yí ge shíxí jīhuì.}\\
\zh{我学习中文已经三年了，我非常感谢您给我这个机会。}\\
\zh{Wǒ xuéxí Zhōngwén yǐjīng sān nián le, wǒ fēicháng gǎnxiè nín gěi wǒ zhè ge jīhuì.}\\
\zh{期待您的回复。}\\
\zh{Qīdài nín de huífù.}\\
\zh{此致敬礼！}\\
\zh{Cǐ zhì jìnglǐ!}
\end{textebox}

Ordre des réponses : \zh{尊敬} → \zh{申请} → \zh{机会} → \zh{感谢} → \zh{回复} → \zh{此致}.\\
\textbf{Justification.} \zh{尊敬的} + titre : appel honorifique ; \zh{申请} + \zh{机会} : collocation \zh{申请一个实习机会} « demander une occasion de stage » ; \zh{感谢您给我这个机会} : remercier quelqu'un de quelque chose ; \zh{期待您的回复} : formule d'attente ; \zh{此致敬礼} : salutation finale.
\end{corrige}

\begin{corrige}[Exercice 8 — Remise en ordre et caractères intrus]
\textbf{A. Remise en ordre}
\begin{enumerate}
\item \zh{我写信是为了申请实习。}\\
\emph{Wǒ xiěxìn shì wèile shēnqǐng shíxí.} — « J'écris pour demander un stage. »\\
Ordre : Sujet \zh{我} + verbe \zh{写信} + \zh{是为了} + but \zh{申请实习}.
\item \zh{非常感谢您的关注。}\\
\emph{Fēicháng gǎnxiè nín de guānzhù.} — « Je vous remercie vivement de votre attention. »\\
L'adverbe \zh{非常} précède le verbe \zh{感谢} ; \zh{您的关注} « votre attention ».
\end{enumerate}
\textbf{B. Caractères intrus}
\begin{enumerate}
\item Intrus : \zh{手}. \zh{语}, \zh{说}, \zh{话} portent le radical de la parole \zh{讠} (liés au langage) ; \zh{手} « main » est lui-même un radical différent.
\item Intrus : \zh{心}. \zh{信}, \zh{你}, \zh{他} portent le radical de l'homme \zh{亻} ; \zh{心} « cœur » est un autre radical.
\item Intrus : \zh{谢}. \zh{想}, \zh{感}, \zh{忘} portent le radical du cœur \zh{心} (sentiments, pensée) ; \zh{谢} porte le radical de la parole \zh{讠}.
\end{enumerate}
\textbf{Point de vigilance.} Le radical donne un indice de sens : \zh{讠} → parole, \zh{亻} → personne, \zh{心} → sentiment. C'est un outil précieux pour deviner le sens d'un caractère inconnu au Bac.
\end{corrige}

\begin{corrige}[Exercice 9 — Type Bac : Compréhension de texte]
\textbf{Questions de compréhension}
\begin{enumerate}
\item \textbf{b} — une élève de Seconde (première année de lycée) : \zh{我是杜阿拉中学一年级的学生} (litt. « élève de 1\textsuperscript{re} année » = \zh{高一} / Seconde en système camerounais).
\item \textbf{c} — juillet : \zh{今年七月的一个实习机会}.
\item Marie parle trois langues : le chinois, le français et l'anglais (\zh{我会说中文、法文和英文}).
\item Dans la pièce jointe se trouve son CV (\zh{附件里有我的简历}).
\item \zh{她学习中文已经两年了。}\\
\emph{Tā xuéxí Zhōngwén yǐjīng liǎng nián le.} — « Elle apprend le chinois depuis deux ans. » Structure : durée + \zh{已经……了}.
\end{enumerate}
\textbf{Questions de langue}
\begin{enumerate}
\item Formule d'appel : \zh{尊敬的经理，您好！} ; formule finale : \zh{此致敬礼！}.
\item \zh{贵} (guì, « honorable ») est un préfixe honorifique placé devant un nom appartenant au destinataire : \zh{贵公司} « votre honorable entreprise », \zh{贵校} « votre honorable école ». Il marque le respect et ne s'emploie jamais pour soi-même (pour soi, on dirait \zh{我们公司}).
\end{enumerate}
\textbf{Barème indicatif (20 points).} Compréhension : questions 1-2 : 2 pts chacune ; questions 3-4 : 3 pts chacune ; question 5 : 4 pts (chinois correct exigé). Langue : question 1 : 3 pts ; question 2 : 3 pts.
\end{corrige}

\begin{corrige}[Exercice 10 — Type Bac : Thème et version]
\textbf{A. Thème}
\begin{enumerate}
\item \zh{我给您写信是为了申请一份正式文件。}\\
\emph{Wǒ gěi nín xiěxìn shì wèile shēnqǐng yí fèn zhèngshì wénjiàn.}\\
Classificateur des documents : \zh{份}. « Officiel » se rend ici par \zh{正式} (« officiel, formel »).
\item \zh{我会说三种语言：中文、法文和英文。}\\
\emph{Wǒ huì shuō sān zhǒng yǔyán: Zhōngwén, Fǎwén hé Yīngwén.}\\
Classificateur des langues : \zh{种} ; énumération avec \zh{、} et \zh{和} avant le dernier élément.
\item \zh{感谢您的回复。}\\
\emph{Gǎnxiè nín de huífù.}
\end{enumerate}
\textbf{B. Version}
\begin{enumerate}
\item « Je vous écris cette lettre pour demander une occasion de stage. » (\zh{这封} : classificateur des lettres \zh{封} ; \zh{是为了} : « c'est pour ».)
\item « Vous trouverez mon CV en pièce jointe ; j'attends votre réponse. » (\zh{附件里有} : « il y a dans la pièce jointe » ; \zh{期待} : « attendre avec espoir ».)
\end{enumerate}
\textbf{Barème indicatif (20 points).} Thème : 4 pts par phrase (lexique 2, syntaxe 1, caractères 1) = 12 pts. Version : 4 pts par phrase (fidélité 2, qualité du français 2) = 8 pts.
\textbf{Points de vigilance.} Bien choisir les classificateurs (\zh{封信}, \zh{份文件}, \zh{种语言}, \zh{个机会}) ; ne pas oublier \zh{您} pour le vouvoiement ; en version, ne pas traduire mot à mot \zh{有} par « avoir » mais restituer « il y a / vous trouverez ».
\end{corrige}

\begin{corrige}[Exercice 11 — Type Bac : Expression écrite — lettre formelle]
\textbf{Proposition de rédaction modèle}

\begin{textebox}[Lettre modèle — demande de stage]
\zh{尊敬的经理，您好！}\\
\zh{Zūnjìng de jīnglǐ, nín hǎo!}\\
\zh{我是雅温得中学高二的学生。我学习中文已经三年了，成绩很好。我写这封信是为了申请贵公司今年八月的一个实习机会。}\\
\zh{Wǒ shì Yǎwēndé zhōngxué gāo èr de xuésheng. Wǒ xuéxí Zhōngwén yǐjīng sān nián le, chéngjì hěn hǎo. Wǒ xiě zhè fēng xìn shì wèile shēnqǐng guì gōngsī jīnnián bāyuè de yí ge shíxí jīhuì.}\\
\zh{我会说中文、法文和英文，也很喜欢电脑。附件里有我的简历。}\\
\zh{Wǒ huì shuō Zhōngwén, Fǎwén hé Yīngwén, yě hěn xǐhuan diànnǎo. Fùjiàn lǐ yǒu wǒ de jiǎnlì.}\\
\zh{感谢您的时间，期待您的回复。}\\
\zh{Gǎnxiè nín de shíjiān, qīdài nín de huífù.}\\
\zh{此致敬礼！}\\
\zh{Cǐ zhì jìnglǐ!}\\
\zh{学生：李明　二〇二四年六月五日}\\
\zh{Xuésheng: Lǐ Míng　èr líng èr sì nián liù yuè wǔ rì}
\end{textebox}

\emph{Traduction : Monsieur le Directeur, bonjour ! Je suis un élève de Première du lycée de Yaoundé. J'apprends le chinois depuis trois ans et j'ai de bons résultats. Je vous écris pour demander un stage dans votre entreprise au mois d'août de cette année. Je parle chinois, français et anglais, et j'aime beaucoup l'informatique. Vous trouverez mon CV en pièce jointe. Je vous remercie de votre temps et j'attends votre réponse. Veuillez agréer mes salutations distinguées ! L'élève : Li Ming, le 5 juin 2024.}

\textbf{Barème indicatif (20 points).}
\begin{itemize}
\item Format et mise en page (appel en haut, formule finale, signature et date en bas à droite) : 5 pts.
\item Vocabulaire honorifique et formules (\zh{尊敬的}, \zh{贵公司}, \zh{感谢}, \zh{此致敬礼}) : 5 pts.
\item Syntaxe et connecteurs (\zh{是为了}, \zh{已经……了}, \zh{和}, \zh{也}) : 5 pts.
\item Exactitude des caractères (pas de confusion \zh{信}/\zh{心}, \zh{实习}/\zh{学习}) : 5 pts.
\end{itemize}

\textbf{Points de vigilance.}
\begin{itemize}
\item Respecter les cinq blocs obligatoires : appel → présentation → objet → qualités → formule finale + signature/date.
\item Vouvoiement systématique avec \zh{您} ; jamais \zh{你} dans une lettre officielle.
\item Date dans l'ordre chinois : année → mois → jour, avec \zh{年}, \zh{月}, \zh{日}.
\item Compter ses caractères : entre 100 et 120 ; une lettre trop courte perd des points de contenu.
\item Écrire \zh{此致} seul sur une ligne puis \zh{敬礼！} (ou \zh{此致敬礼！} ensemble) : les deux présentations sont acceptées.
\end{itemize}
\end{corrige}

\newpage

\coursec{Fiche bilan}

\begin{popfigure}
\begin{tikzpicture}[mindmap, grow cyclic, every node/.style={concept, font=\small},
  root concept/.append style={concept color=popblue, font=\bfseries},
  level 1/.append style={sibling angle=60, level distance=3.2cm},
  level 2/.append style={level distance=2.2cm, font=\scriptsize}]
\node[root concept, text=white]{Lettre formelle 正式信函}
  child[concept color=poporange] { node[concept]{Structure 格式}
    child { node[concept]{称呼 Salutation} }
    child { node[concept]{正文 Corps} }
    child { node[concept]{此致敬礼 Formule finale} }
  }
  child[concept color=popgreen] { node[concept]{Lexique 词汇}
    child { node[concept]{申请 postuler} }
    child { node[concept]{实习 stage} }
    child { node[concept]{尊敬 respecté} }
  }
  child[concept color=poppurple] { node[concept]{Connecteurs 连接词}
    child { node[concept]{首先 d'abord} }
    child { node[concept]{因此 donc} }
    child { node[concept]{希望 espérer} }
  }
  child[concept color=poppink] { node[concept]{Caractères 汉字}
    child { node[concept]{讠parole : 请 谢} }
    child { node[concept]{忄cœur : 情 怀} }
  }
  child[concept color=popgold] { node[concept]{Méthode 方法}
    child { node[concept]{Qui ? Quoi ? Pourquoi ?} }
    child { node[concept]{Ton poli 您} }
  }
  child[concept color=popblue!60] { node[concept]{Bac 考试}
    child { node[concept]{150--200 caractères} }
    child { node[concept]{Relire les tons} }
  };
\end{tikzpicture}
\legende{Carte mentale de la leçon : rédiger une lettre formelle en chinois}
\end{popfigure}

\begin{bilan}
\textbf{1. Lexique clé de la correspondance administrative}
\begin{center}
\begin{tabularx}{\linewidth}{|X|X|X|X|}
\hline
\rowcolor{popblueL} Caractères & Pinyin & Nature & Traduction \\
\hline
尊敬 & zūnjìng & verbe / adj. & respecter, respecté \\
\hline
申请 & shēnqǐng & verbe / nom & postuler, demande \\
\hline
实习 & shíxí & verbe / nom & faire un stage, stage \\
\hline
机会 & jīhuì & nom & occasion, opportunité \\
\hline
经验 & jīngyàn & nom & expérience \\
\hline
负责 & fùzé & verbe & être responsable de \\
\hline
附件 & fùjiàn & nom & pièce jointe \\
\hline
此致敬礼 & cǐ zhì jìnglǐ & formule & formule de politesse finale \\
\hline
\end{tabularx}
\end{center}

\textbf{2. Règles à connaître par cœur}
\begin{center}
\begin{tabularx}{\linewidth}{|X|X|X|}
\hline
\rowcolor{popblueL} Règle & Structure & Exemple \\
\hline
Salutation formelle & 尊敬的 + titre/nom + ： & 尊敬的王经理： \\
\hline
Se présenter & 我是 + identité & 我是雅温得中学的学生。 \\
\hline
Exprimer le but & 我想 + verbe / 为了 + but & 我想申请一个实习机会。 \\
\hline
Argumenter & 首先……其次……最后…… & 首先，我学习中文三年了。 \\
\hline
Conclure & 希望 + proposition + 。此致敬礼！ & 希望得到您的回复。 \\
\hline
Signature & nom + date, en bas à droite & 李明 2025年5月10日 \\
\hline
\end{tabularx}
\end{center}

\textbf{3. Méthodes express}
\begin{itemize}
\item \cle{Plan en 4 temps} : salutation $\rightarrow$ présentation $\rightarrow$ demande argumentée $\rightarrow$ formule finale + signature.
\item \cle{Ton soutenu} : utiliser \zh{您} (et jamais \zh{你}), éviter les impératifs directs, préférer \zh{希望} et \zh{请}.
\item \cle{Compréhension de texte} : repérer d'abord l'expéditeur, le destinataire, l'objet, puis les arguments (connecteurs).
\item \cle{Caractères} : mémoriser par clés — \zh{讠} (parole) dans 请、谢、语 ; écrire de haut en bas, de gauche à droite, l'horizontal avant le vertical.
\end{itemize}
\end{bilan}

\begin{aretenir}[Checklist avant le Bac]
\begin{itemize}
\item[$\square$] Je connais par cœur la structure d'une lettre formelle (salutation, corps, formule finale, signature).
\item[$\square$] Je sais écrire et prononcer les 8 mots clés du lexique (tons corrects).
\item[$\square$] J'utilise \zh{您} et un ton poli dans toute la lettre.
\item[$\square$] Je sais ouvrir la lettre avec \zh{尊敬的……：} et la fermer avec \zh{此致敬礼！}.
\item[$\square$] J'annonce mon objectif avec \zh{我想申请……} ou \zh{为了……}.
\item[$\square$] Je relie mes arguments avec \zh{首先}、\zh{其次}、\zh{最后}.
\item[$\square$] Je place la date en chinois : année 年 + mois 月 + jour 日.
\item[$\square$] Je reconnais les clés \zh{讠} et \zh{忄} et je respecte l'ordre des traits.
\item[$\square$] Je vérifie la ponctuation chinoise pleine chasse ：，。！
\item[$\square$] Je relis ma lettre : tons du pinyin, ordre des mots Sujet + Verbe + Complément.
\item[$\square$] Je respecte la longueur demandée (environ 150 à 200 caractères).
\item[$\square$] Je signe en bas à droite avec mon nom et la date.
\end{itemize}
\end{aretenir}

\findecours

\end{document}
